\documentclass[11pt]{article}
\usepackage[T1]{fontenc}
\usepackage{lmodern}
\usepackage{amsmath,amssymb,amsthm,mathtools}
\usepackage[margin=1in]{geometry}
\usepackage{booktabs,microtype,hyperref}
\hypersetup{colorlinks=true,linkcolor=blue,citecolor=blue,urlcolor=blue,
pdftitle={Small-point global maxima and uniform exact certification for the Mordell discriminant problem},
pdfauthor={Fengcan Jiang}}
\setlength{\emergencystretch}{2em}
\newtheorem{theorem}{Theorem}[section]
\newtheorem{proposition}[theorem]{Proposition}
\newtheorem{lemma}[theorem]{Lemma}
\newtheorem{corollary}[theorem]{Corollary}
\theoremstyle{remark}
\newtheorem{remark}[theorem]{Remark}
\newcommand{\C}{\mathbb C}
\newcommand{\R}{\mathbb R}
\newcommand{\Q}{\mathbb Q}
\newcommand{\E}{\mathbb E}
\numberwithin{equation}{section}
\title{Small-point global maxima and uniform exact certification\\for the Mordell discriminant problem}
\author{Fengcan Jiang\thanks{Email:
\href{mailto:2255214397@stu.xjtu.edu.cn}{\texttt{2255214397@stu.xjtu.edu.cn}}.}}
\date{October 5, 2026}
\begin{document}
\maketitle
\begin{abstract}
We determine the six- and seven-point maxima of the squared Vandermonde
product under a prescribed total quadratic moment. Their unique maximizing
shapes are the origin with a regular pentagon or hexagon, with values
$6^{15}/5^{10}$ and $7^{21}/6^{15}$ at moments six and seven.
The proofs invoke the published Faug\`ere--Svartz decomposition and physical
coverage, followed by independent integer and rational root-disk certificates.
The seven-point check covers $174$ algebraic candidates and $1218$ point roots;
the external exclusion of a degree-$356$ branch is stated explicitly.
Small cases supply a cubic selector for all physical four-point stationary
values and explicit exact sum-of-squares certificates.
We prove finite exact certificates after critical-value augmentation, without
assuming isolated stationary shapes, and simplify conjugate pairing by a
bordered Hermitian trace test. A simultaneous pair-replacement argument
strengthens the global separation bound and the compact certification model.
Uniform convergent moment and spectral bounds, deletion monotonicity, radial
cuts and a fixed global disk support the framework for arbitrary point counts.
An explicit rational polar construction matches the first two asymptotic
terms but is strictly suboptimal at linear order. Classical renormalized-energy
theory identifies the global linear coefficient with a microscopic energy
minimum. No finite simple closed expression for every point count is obtained.
\end{abstract}

\section{Scope and normalization}
For $n\ge2$, put
\[
 \Delta_n(z)=\prod_{i<j}|z_i-z_j|^2,\qquad
 M_n=\max_{\sum_i|z_i|^2=n}\Delta_n(z),\qquad
 N=\frac{n(n-1)}2,\qquad P_n=\prod_{k=0}^{n-1}k!.
\]
The maximum exists by compactness; it is positive and has no
collisions, since a unit regular polygon has value $n^n$.
Every maximizer has zero centroid: centering preserves the distances
and decreases the quadratic moment by $n$ times the squared centroid;
rescaling a nonzero centroid configuration to the constraint would
strictly increase its product.

The five-point equality $M_5=5^5$ is proved by Wang \cite{Wang}.
Chalk's center-plus-polygon construction \cite{Chalk} disproves
$M_n=n^n$ for $n\ge6$. The present note provides general constructions
and global bounds. The six-point value is determined below using the
published exhaustive coefficient decomposition and a new independent
rational certificate. The current revision gives additional structural results and an explicit
all-point-count rational construction, while distinguishing attainable
values from global maxima. The seven-point value and finite augmented certificates are established in Section~\ref{sec:examples-bordered-trace}. No finite closed formula for all $M_n$ is claimed.
The asymptotic context is classical Gaussian weighted Fekete theory
and the planar logarithmic gas \cite{AOC,SS}. No claim of priority is
made for these asymptotic terms or elementary bounds.

\section{The exact six-point global maximum}
\label{sec:six-global-exact}

The exhaustive six-vortex coefficient decomposition of
Faug\`ere--Svartz~\cite[Sections~5.3 and~5.7]{FS}, together with
the archived author bases~\cite{FSdata}, supplies eleven triangular
components, covering eighty complex coefficient solutions after the
normalization by the first nonzero elementary symmetric coefficient.
We use the authors' lexicographic bases, rather than a numerical list
of equilibria.  The archived file is reproduced with the certificate.
Its first comment says \texttt{N=7}; its eleven bases involve only
$e_2,\ldots,e_6$ and are the six-point data linked as
\texttt{vortices\_N6.txt} on the authors' index.
The accompanying author solution document, whose six-vortex section
is on pages~7--10, independently identifies these data.
We invoke the published algebraic decomposition as an external
completeness input: every physical stationary six-point shape is
represented by these bases after coefficient normalization.
Only this necessary coverage is used; homogenized coefficient
solutions may also include nonphysical candidates.
The verifier below certifies the subsequent root
coverage and comparisons; it does not recompute that decomposition.
The coefficient convention is
\begin{equation}\label{eq:six-coeff-convention}
 P(t)=t^6+e_2t^4-e_3t^3+e_4t^2-e_5t+e_6.
\end{equation}
The certificate below tests the actual conjugate pairing of every
candidate; an equality between unordered sets of dual roots is
insufficient.

\begin{theorem}\label{thm:six-global}
For $\sum_{i=1}^6|z_i|^2=6$,
\[
 \prod_{i<j}|z_i-z_j|^2\le
 \frac{6^{15}}{5^{10}}=\frac{470184984576}{9765625}.
\]
Equality holds precisely for the origin together with a regular
pentagon of radius $\sqrt{6/5}$, up to rotation and permutation.
\end{theorem}

\begin{proof}
A maximizer is collision-free and centered.  Its coefficients belong
to the exhaustive decomposition just described, after multiplication
of all roots by a common nonzero complex number.  If $w_1,\ldots,w_6$
are the roots in one such coefficient chart, write
$Q=\sum_i|w_i|^2$.  Physical stationarity requires
\[
 \frac{P''(w_i)}{2P'(w_i)}=\kappa\overline{w_i},
 \qquad \kappa>0.
\]
Multiplication by $w_i$ and summation gives $\kappa Q=15$;
therefore every root must satisfy the polynomial residual identity
\begin{equation}\label{eq:six-force-residual}
 QP''(w_i)-30\overline{w_i}P'(w_i)=0.
\end{equation}
The objective of the configuration rescaled to moment six is
\begin{equation}\label{eq:six-chart-value}
 \left(\frac6Q\right)^{15}\prod_{i<j}|w_i-w_j|^2.
\end{equation}

The independent rational-arithmetic verifier in the certificate
bundle exhausts the eighty coefficient solutions.  It excludes
sixty-nine solutions by \eqref{eq:six-force-residual}.  The remaining
eleven consist of the central pentagon and ten candidates with
certified strict upper bounds below $6^{15}/5^{10}$, as follows.
The component numbers are those of the supplied author file;
opposite signs of $e_3$ may describe the same shape after rotation.
\begin{center}
\begin{tabular}{rrrrl}
\toprule
Component&Solutions&Excluded&Not excluded&Upper bound\\
\midrule
1&22&20&2&$3170$\\
2&28&28&0&---\\
3&12&12&0&---\\
4&4&2&2&$42454$\\
5&2&2&0&---\\
6&1&0&1&$18699$\\
7&1&0&1&$18699$\\
8&6&4&2&$22400$ and $133$\\
9&2&1&1&$7225$\\
10&1&0&1&$6^{15}/5^{10}$, exact\\
11&1&0&1&$6^6=46656$, exact\\
\bottomrule
\end{tabular}
\end{center}
Here ``not excluded'' means that the necessary force test did not
exclude the candidate; no sufficiency assertion is needed for the
ten strict upper bounds.
Component ten has $P(t)=t(t^5-1)$, whose roots are the origin and a
unit regular pentagon.  Its moment is five, and its discriminant
product is $5^5$; \eqref{eq:six-chart-value} gives $6^{15}/5^{10}$.
Component eleven has $P(t)=t^6+1$ and is a unit regular hexagon.
All other remaining rows have a strict upper bound smaller than
the central-pentagon value.  This proves the upper bound and the
equality classification.
\end{proof}

\subsection{Rational verification of root coverage and inequalities}
\label{sec:six-disk-certificate}
The frozen centers in \texttt{fs6\_disk\_certificate.json} are
Gaussian rationals.  Parameter disks have radius $10^{-70}$; the
six root disks for each parameter have radius $10^{-12}$.
Floating-point arithmetic is used only to propose these centers.
Verification uses integer and rational arithmetic, and does not
import a numerical package or a computer algebra system.

For a Gaussian rational $c$, let $|c|_1=|\Re c|+|\Im c|$ and
$|c|_\infty=\max(|\Re c|,|\Im c|)$.  Suppose a polynomial
$p_0$ approximates $p$, and the error in its coefficient of $t^j$
is at most $\epsilon_j$.  Put
$a_k=p_0^{(k)}(c)/k!$.  The checker verifies
\begin{equation}\label{eq:six-rouche}
 |a_0|_1+\sum_{k\ge2}|a_k|_1r^k
 +\sum_j\epsilon_j(|c|_1+r)^j
 <r|a_1|_\infty.
\end{equation}
On the circle $|t-c|=r$, this bounds the difference between
$p(t)$ and $a_1(t-c)$ strictly below $|a_1(t-c)|$.
Rouch\'e's theorem proves that each disk contains exactly one
root, counted with multiplicity.  The disks are pairwise disjoint,
and their number equals the degree of the relevant polynomial;
hence no complex roots are omitted.

For a coefficient expressed as $f(b)/A$, where the parameter lies
in $|b-b_0|\le\rho$, its certified error before rational rounding
is bounded by
\[
 \frac1{|A|}\sum_k|f_k|
 \bigl\{(|b_0|_1+\rho)^k-|b_0|_1^k\bigr\}.
\]
This propagates the parameter disk into coefficient errors in
\eqref{eq:six-rouche}, so the six root disks are valid uniformly
for the exact algebraic parameter.

Let $c_1,\ldots,c_6$ be the root-disk centers, and let their common
radius be $r$.  The verifier sets
\[
 Q_0=\sum_i|c_i|^2,\qquad
 E_Q=\sum_i(2|c_i|_1r+r^2).
\]
Then $|Q-Q_0|\le E_Q$.  Taylor bounds for $P'$ and $P''$ and
the coefficient-error bounds give an explicit rational error for
\eqref{eq:six-force-residual}.  If the infinity norm of its center
residual exceeds that error, the candidate is excluded.  The
sixty-nine exclusions all have strictly positive rational margins.
For every remaining candidate outside components ten and eleven,
the checker verifies $Q_0>E_Q$ and the strict rational inequality
\begin{equation}\label{eq:six-disk-upper}
 \left(\frac6{Q_0-E_Q}\right)^{15}
 \prod_{i<j}\left(
 |c_i-c_j|^2+4|c_i-c_j|_1r+4r^2\right)
 <\frac{6^{15}}{5^{10}}.
\end{equation}
The expression on the left bounds \eqref{eq:six-chart-value}.
The integer bounds in the table are the upward-rounded endpoints
of these rational bounds; the hexagon value is computed exactly.
All decisive checks are explicit exceptions, so running Python
with optimization enabled does not disable certification.

\subsection{The exact coefficient equations and the classification input}
\label{subsec:fs6-equations-audit}
For completeness, we specify both the equation convention and the external
classification input. Write
\[
 p(X)=X^6+e_2X^4-e_3X^3+e_4X^2-e_5X+e_6
\]
for the centered polynomial with distinct roots $x_1,\ldots,x_6$, and set
\[
 \delta=\operatorname{Res}(p,p'),\qquad
 g_i=\frac{p''(x_i)}{p'(x_i)},\qquad
 S_k=\sum_i g_i^k,\qquad m_k=\sum_i x_i g_i^k.
\]
The basic partial-fraction identity gives
$g_i=2\sum_{j\ne i}(x_i-x_j)^{-1}$, hence
$S_1=m_0=0$ and $m_1=30$.
Define the dual resultant and its first variation by
\[
 R_0(Y)=\operatorname{Res}_X(p,Yp'-p''),\qquad
 R_t(Y)=\left.\partial_T\operatorname{Res}_X
 (p,(Y+TX)p'-p'')\right|_{T=0}.
\]
Then, directly from the product formula for resultants,
\[
 R_0=\delta\prod_i(Y-g_i),\qquad
 R_t=\delta\sum_i x_i\prod_{j\ne i}(Y-g_j).
\]
Let $R_0/\delta=Y^6+q_2Y^4+q_3Y^3+q_4Y^2+q_5Y+q_6$.
The four polynomial equations used in the certificate are
\begin{align*}
 E_1&=[Y^3]R_t, & E_2&=[Y^2]R_t-12[Y^4]R_0,\\
 E_3&=[Y]R_t-6[Y^3]R_0, & E_4&=[1]R_t-2[Y^2]R_0.
\end{align*}
Synthetic division gives
\begin{align*}
 E_1/\delta&=m_2, & E_2/\delta&=m_3+18q_2,\\
 E_3/\delta&=m_4+q_2m_2+24q_3,
 &E_4/\delta&=m_5+q_2m_3+q_3m_2+28q_4.
\end{align*}
Using $S_2=-2q_2$, $S_3=-3q_3$ and
$S_4=2q_2^2-4q_4$, we obtain the exact equivalence
\[
 E_1=\cdots=E_4=0
 \quad\Longleftrightarrow\quad
 m_2=0,\quad m_3=9S_2,\quad m_4=8S_3,
 \quad m_5=7S_4+S_2^2.
\]
These are the first four nontrivial symmetric vortex equations in the
Faug\`ere--Svartz computation, with the doubled force $g=p''/p'$.
If the $g_i$ are also distinct, define
$h_i=x_i-2\sum_{j\ne i}(g_i-g_j)^{-1}$.
The elementary pairwise-sum identity gives
\[
 2\sum_i g_i^k\sum_{j\ne i}\frac1{g_i-g_j}
 =\sum_{a=0}^{k-1}(S_{k-1-a}S_a-S_{k-1}).
\]
Consequently the four equations, together with $m_0=0,m_1=30$,
are precisely $\sum_i h_i g_i^k=0$ for $0\le k\le5$.
The Vandermonde matrix in the distinct $g_i$ is invertible, so $h_i=0$.
The symmetric equations therefore recover the double-force system on this
collision-free locus. Physical solutions additionally require
\[
 g_i=\frac{30}{Q}\overline{x_i},\qquad Q=\sum_i|x_i|^2,
\]
which is exactly the residual checked by the rational certificate.

The external classification result invoked here is the exhaustive rational
Gr\"obner and triangular computation reported by Faug\`ere and Svartz,
ISSAC 2012, Sections~5.3--5.7, together with their archived author data
\texttt{vortices\_N6.txt}. Section~5.3 splits according to the first nonzero
coefficient, normalizes it to $1$, adds the collision-saturation equation
$h\delta=1$, and decomposes every chart. Section~5.7 reports the exact
$n=6$ rational computation (Table~4). Only its coverage consequence is
needed: every physical stationary six-point shape occurs among the
eleven supplied triangular components. Theorem~5 of that paper supplies
the symmetric equations; it should not be described as a separate numbered
theorem asserting the eleven-component classification.
Complex scaling used for chart normalization changes the force constant
to a positive multiple, so it loses no physical shape. Although the data
file's first comment says $N=7$, the author index names it as the six-point
file, and Section~4 of the author solution appendix explicitly labels the
matching six-point catalogue. This is a harmless header typo.

As an additional independent integrity check, exact rational substitution
verifies $E_1=\cdots=E_4=0$ on the first eight components; direct
power-sum calculations verify the three exceptional components.
The discriminant is relatively prime to the corresponding eliminants.
This check proves containment and consistency of conventions. Exhaustive
coverage uses the published classification computation just cited;
the subsequent root isolation, physical exclusions and objective bounds
are certified independently by the supplied rational checker.

\section{Explicit universal bounds}
\label{general_sec_bounds}
\begin{theorem}\label{general_thm_bounds}
Put $S_m(x)=\sum_{k=0}^m x^k/k!$. For every $\alpha>0$,
\begin{equation}\label{general_eq_bounds}
 L_n:=\frac{n!(n-2)!\,n^N P_n}{(N+n-2)!}
 \le M_n\le U_n(\alpha):=\frac{P_n S_{n-1}(\alpha)^n}{\alpha^N}.
\end{equation}
The upper bound has a unique minimizing parameter $\alpha_n>0$,
characterized by
\begin{equation}\label{general_eq_optimal_alpha}
 \sum_{k=0}^{n-1}(2k-(n-1))\frac{\alpha_n^k}{k!}=0.
\end{equation}
In particular,
\begin{equation}\label{general_eq_exponential_upper}
 M_n\le\widehat U_n:=e^N\left(\frac nN\right)^NP_n.
\end{equation}
\end{theorem}
\begin{proof}
Take independent standard complex Gaussian variables $Z_j$,
each with density $\pi^{-1}e^{-|z|^2}$. Their moments satisfy
$\E(Z^k\overline Z^{\,\ell})=0$ for $k\ne\ell$ and $k!$ for $k=\ell$.
Expansion of the squared Vandermonde determinant gives
$\E\Delta_n(Z)=n!P_n$.

Let $Y_j=Z_j-n^{-1}\sum_iZ_i$ and $R^2=\sum_j|Y_j|^2$.
A unitary change of Gaussian coordinates shows that $R^2$ has the
$\Gamma(n-1,1)$ distribution and is independent of
$W=\sqrt n\,Y/R$, the uniform direction on the zero-centroid sphere.
Translation invariance and homogeneity give
\[
 n!P_n=\E\Delta_n(Y)
 =\frac{\Gamma(N+n-1)}{\Gamma(n-1)n^N}\,\E\Delta_n(W).
\]
Thus $L_n=\E\Delta_n(W)\le M_n$.

For the upper bound, apply row Hadamard to
$A_{jk}=\alpha^{k/2}z_j^k/\sqrt{k!}$, $0\le k<n$.
Writing $t_j=|z_j|^2$, this gives
\[
 \frac{\alpha^N}{P_n}\Delta_n(z)
 \le\prod_{j=1}^nS_{n-1}(\alpha t_j)
 \le S_{n-1}(\alpha)^n.
\]
The second inequality is Jensen, since $\sum_jt_j=n$
and $\log S_m$ is concave on $[0,\infty)$.
For $m=1$ this is immediate. For $m\ge2$, use $S'_m=S_{m-1}$,
$S''_m=S_{m-2}$ and
\[
 S_{m-1}(x)^2-S_m(x)S_{m-2}(x)
 =\frac{x^{m-1}}{(m-1)!}
 \left(1+\sum_{k=1}^{m-1}
         \left(1-\frac km\right)\frac{x^k}{k!}\right)\ge0.
\]
This proves \eqref{general_eq_bounds}.

Set $m=n-1$ and $q(\alpha)=\alpha S'_m(\alpha)/S_m(\alpha)$.
For the distribution on $\{0,\ldots,m\}$ with probabilities
proportional to $\alpha^k/k!$, $q$ is the mean of $k$, and
$dq/d\log\alpha$ is its strictly positive variance.
Thus $q$ increases from $0$ to $m$. The logarithmic derivative of
$U_n(\alpha)$ with respect to $\log\alpha$ is $nq(\alpha)-N$,
so its unique zero is the global minimum, giving
\eqref{general_eq_optimal_alpha}. Finally, $S_m(\alpha)\le e^\alpha$
and $\alpha=N/n$ prove \eqref{general_eq_exponential_upper}.
\end{proof}

The lower bound is an average, not a proposed extremal configuration.
For rational $\alpha$, $U_n(\alpha)$ is rational, giving a fully exact
finite interval. The optimal parameter is algebraic by
\eqref{general_eq_optimal_alpha} and can be enclosed by polynomial
root isolation. For $n=2$, $L_2=U_2(1)=4$.
For example, $\alpha=11/4$ gives the exact bound $M_6<89309$, whereas
the central pentagon gives $M_6\ge48146.9424205824$.
The decimal lower endpoint here is exact; its value is $6^{15}/5^{10}$.

\begin{corollary}\label{general_cor_asymptotic}
With $c=\tfrac12\log2-\tfrac14$,
\begin{equation}\label{general_eq_asymptotic}
 \log M_n=cn^2+\frac n2\log n+O(n).
\end{equation}
\end{corollary}
\begin{proof}
Summing Stirling's formula, or applying Euler--Maclaurin to
$\log P_n=n\log\Gamma(n)-\sum_{k=1}^{n-1}k\log k$, gives
\[
 \log P_n=\frac{n^2}{2}\log n-\frac{3n^2}{4}
            +\frac n2\log(2\pi)+O(\log n).
\]
Substitution in \eqref{general_eq_exponential_upper} yields
\[
 \log\widehat U_n=cn^2+\frac n2\log n
                      +\frac n2\log\pi+O(\log n).
\]
The exact ratio of the upper and lower bounds is
\[
 \frac{\widehat U_n}{L_n}
 =e^NN^{-N}\frac{(N+n-2)!}{n!(n-2)!}.
\]
Stirling gives the logarithm of this ratio as
$(2-\log2)n+O(\log n)$, proving the assertion.
\end{proof}

\section{Global geometry and monotonicity across point counts}
\label{sec:deletion-radius}

The exact six-point value can be propagated to larger point counts
without assuming their maximizing shapes. A separate weighted-polynomial
argument confines every global maximizer to a fixed disk. These are
finite, explicit improvements to the certification framework; neither
asserts a closed formula for every maximum.

\subsection{Subset deletion and a strictly decreasing normalized maximum}
Write $N_j=j(j-1)/2$ and
\[
 d_j=\frac{j-1}{j}M_j^{1/N_j}.
\]
\begin{theorem}\label{thm:subset-monotonic}
For $2\le m<n$,
\begin{equation}\label{eq:subset-upper}
 M_n\le
 \left[\frac{n(m-1)}{m(n-1)}M_m^{1/N_m}\right]^{N_n}.
\end{equation}
Equivalently, $d_n\le d_m$. Moreover,
\[
 d_2=d_3=2,\qquad d_n<d_{n-1}\quad(n\ge4),\qquad
 \lim_{n\to\infty}d_n=\frac2{\sqrt e}.
\]
In particular, the six-point theorem gives the explicit bound
\begin{equation}\label{eq:six-seed-all-n}
 M_n<\left(\frac{n}{n-1}\,5^{1/3}\right)^{N_n}
 \qquad(n>6).
\end{equation}
\end{theorem}
\begin{proof}
Choose a centered maximizer with moment $n$. For each $m$-element
subset $S$, let $\mu_S=m^{-1}\sum_{i\in S}z_i$ and
\[
 Q_S=\sum_{i\in S}|z_i-\mu_S|^2
     =\frac1m\sum_{\{i,j\}\subset S}|z_i-z_j|^2.
\]
All $Q_S$ are positive. Homogeneity and the definition of $M_m$ give
\[
 \Delta_m(z_S)\le M_m(Q_S/m)^{N_m}.
\]
Each pair occurs in exactly $\binom{n-2}{m-2}$ subsets, while
centering gives $\sum_{i<j}|z_i-z_j|^2=n^2$. Therefore
\[
 \prod_{|S|=m}\Delta_m(z_S)
   =\Delta_n(z)^{\binom{n-2}{m-2}},\qquad
 \frac1{\binom nm}\sum_{|S|=m}Q_S
   =\frac{n(m-1)}{n-1}.
\]
Apply the arithmetic--geometric mean inequality to $Q_S/m$ and
use $\binom nm N_m=\binom{n-2}{m-2}N_n$. This proves
\eqref{eq:subset-upper} and monotonicity.

For strictness take $m=n-1$ and suppose equality holds. Equality
in the arithmetic--geometric mean implies all deleted moments are
equal. Since
\[
 Q_{\{1,\ldots,n\}\setminus\{i\}}
  =n-\frac n{n-1}|z_i|^2,
\]
we have $|z_i|^2=1$ for every $i$. Every deleted subset must also
be a global maximizer after centering and scaling. Its centered
coordinates are $v_j=z_j+z_i/(n-1)$, its moment is
$n(n-2)/(n-1)$, and its force constant is
$\kappa=(n-1)^2/(2n)$. Subtract its stationary force equation from
the original one, whose constant is $(n-1)/2$. For every $i\ne j$
this gives
\[
 \frac1{z_j-z_i}
 =\frac{n-1}{2n}(\overline z_j-\overline z_i),\qquad
 |z_j-z_i|^2=\frac{2n}{n-1}.
\]
Four distinct equidistant points cannot lie in the real plane:
after translating one to zero, the Gram matrix of the other three
would have diagonal $a$ and off-diagonal $a/2$ for $a>0$, hence rank
three. This contradicts planarity when $n\ge4$. Thus every such
step is strict. The exact values $M_2=4,M_3=27$ give $d_2=d_3=2$.
Finally Corollary~\ref{general_cor_asymptotic} gives
$\log d_n\to2c=\log2-1/2$. The six-point value gives
$d_6=5^{1/3}$, proving \eqref{eq:six-seed-all-n}.
\end{proof}

There is also a shape-sensitive version of the one-point deletion bound.
For every centered moment-$n$ configuration with distinct points and
$n\ge3$, writing $t_i=|z_i|^2$, one has
\begin{equation}\label{eq:deletion-variance}
\begin{split}
 \log\Delta_n(z)\le{}&\frac n{n-2}\log M_{n-1}
  +N_n\log\frac{n(n-2)}{(n-1)^2}\\
 &-\frac1{4(n-1)}\sum_i(t_i-1)^2.
\end{split}
\end{equation}
Indeed $x_i=Q_{-i}/(n-1)=n/(n-1)-nt_i/(n-1)^2$
lies in $(0,n/(n-1)]$, where $(\log x)''\le-(n-1)^2/n^2$.
Taylor's inequality around the common mean
$n(n-2)/(n-1)^2$, followed by summation and the deletion product
identity, gives precisely the displayed variance penalty. The
resulting recursive bound is useful at modest point counts and can
be intersected with the Hadamard, moment and spectral bounds. Its
six-point seed alone is weaker at the leading asymptotic scale.

\subsection{A fixed disk containing every global maximizer}
\begin{theorem}\label{thm:global-fixed-disk}
Every global maximizer in the original moment-$n$ normalization satisfies
\begin{equation}\label{eq:global-fixed-radius}
                         |z_i|^2\le2\quad(1\le i\le n).
\end{equation}
\end{theorem}
\begin{proof}
Put $\lambda=(n-1)/2$. Homogeneity shows that a moment-$n$
maximizer also maximizes
\[
 \Phi(z)=\Delta_n(z)\exp\left(-\lambda\sum_i|z_i|^2\right)
\]
over the whole space: at a fixed nonzero shape with moment $Q$,
the scale $s$ maximizes $s^{2N_n}e^{-\lambda s^2Q}$ at
$s^2Q=N_n/\lambda=n$.

Fix all coordinates except $z_i$ and put
$p(u)=\prod_{j\ne i}(u-z_j)$, a polynomial of degree $d=n-1$.
The remaining objective is $|p(u)|^2e^{-\lambda|u|^2}$.
For any $R>0$, the maximum-modulus principle, applied to
$p(u)/u^d$ on the exterior disk including infinity, gives
\[
 |p(u)|\le\left(\frac{|u|}R\right)^d
                  \max_{|v|=R}|p(v)|\qquad(|u|\ge R).
\]
Equivalently, apply the usual disk principle to the polynomial
$v^dp(1/v)$. Choose $R^2=d/\lambda=2$. For $r>R$,
\[
 \left(\frac rR\right)^{2d}e^{-\lambda(r^2-R^2)}<1,
\]
because its logarithmic derivative is $2d/r-2\lambda r<0$.
Thus an exterior point has a strictly smaller weighted polynomial
value than some point on $|v|=R$. It cannot be a global maximizer.
This proves \eqref{eq:global-fixed-radius}.
\end{proof}

The fixed-disk argument uses global optimality and must not be applied
indiscriminately to all stationary shapes or local maxima. The radial
stability cuts proved earlier have the broader local scope.

\subsection{A smaller compact trace model}
The fixed disk also strengthens the coefficient model. In force-$1$
coordinates $w_i=\sqrt{(n-1)/2}\,z_i$, it says $|w_i|^2\le n-1$.
Put $a_n=\lceil\sqrt{n-1}\rceil$ and
\[
 \mu'_n=\frac1{n(1+a_n)^{2(n-1)}},\qquad
 C'_n=na_n(1+a_n)^{n-1}.
\]
Every global maximum satisfies the rational closed bounds
\begin{equation}\label{eq:smaller-coefficient-bounds}
 H\succeq\mu'_n I,\qquad
 |c_k|^2\le\binom nk^2(n-1)^k,\qquad
 |r_\ell|^2\le(C'_n)^2.
\end{equation}
To see this, each Lagrange polynomial has coefficient $\ell^1$ norm
at most $(1+a_n)^{n-1}$: every root has modulus at most $a_n$,
and all separating factors have modulus greater than one by
Lemma~\ref{lem:global-separation} after force rescaling. Thus
$\|V^{-1}\|_F^2\le n(1+a_n)^{2(n-1)}$ and
$\|R\|_1\le na_n(1+a_n)^{n-1}$. The elementary symmetric functions
give the middle bound without the integer rounding.

The disk itself can be imposed without choosing roots. Extend the
quotient traces $H_{jk}=\operatorname{Tr}(t^jR^k)$ to indices $n$
by monic reduction, and form the localizing matrix
\[
 J_{jk}=(n-1)H_{jk}-H_{j+1,k+1},\qquad 0\le j,k<n.
\]
On a genuinely conjugate-paired configuration,
\[
 J=V\,\operatorname{diag}(n-1-|w_i|^2)\,V^*.
\]
Since $V$ is invertible, $J\succeq0$ is equivalent to the radius
condition at every root. Replacing the bounds of
Theorem~\ref{thm:compact-trace} by
\eqref{eq:smaller-coefficient-bounds} and adding $J\succeq0$ gives
a smaller compact rational semialgebraic model with exactly the
same maximum. It retains every global maximizer and still excludes
false conjugate pairings. For six points the positivity margin
improves from $1/362797056$ to $1/6291456$, and the auxiliary
coefficient bound decreases from $233280$ to $18432$.

\begin{corollary}\label{cor:fixed-disk-grid}
In the certified grid construction, the coordinate box can be reduced to
\[
 |a|,|b|\le R_q,\qquad
 R_q=\min\{R\in\mathbb Z_{\ge0}:R^2\ge2q^2\}.
\]
The complete grid maximum $B_n^{\rm disk}(q)$ still satisfies
\[
 B_n^{\rm disk}(q)\le M_n
 \le B_n^{\rm disk}(q)e^{E_n(q)}\qquad(q\ge n),
\]
with the same exponent as Theorem~\ref{thm:grid-sharper}. For logarithmic
accuracy $0<\varepsilon\le1$, choosing
$q=\lceil n/\sqrt\varepsilon\rceil$ gives coordinate height
$O(n/\sqrt\varepsilon)$. For $n=6,q=10$, the box radius is $15$.
\end{corollary}
\begin{proof}
Rounding $qz_i$ of a global maximizer lies in this box by
\eqref{eq:global-fixed-radius}. Every subsequent centering, normalization,
distinctness and Taylor estimate is unchanged. Completeness of the
enumeration is still essential; the smaller box alone does not make
high-dimensional exhaustive enumeration inexpensive.
\end{proof}

\section{A deterministic construction for every number of points}
\label{general_sec_lattice}
The uniform disk limit is a classical feature of quadratic weighted
Fekete problems; see \cite{AOC,SS}. The following construction supplies
explicit points and exact finite products. Its optimality statement
concerns the leading $n^2$ scale, not the maximum at each finite $n$.

\begin{theorem}\label{general_thm_lattice}
Order the Gaussian integers $v\in\mathbb Z+i\mathbb Z$ by increasing
$|v|^2$, breaking ties lexicographically by $(\Re v,\Im v)$, and take
the first $n$ distinct points $v_1,\ldots,v_n$. Define
\begin{equation}\label{general_eq_lattice}
 \mu_v=\frac1n\sum_jv_j,\quad
 D_n=n\sum_j|v_j|^2-\left|\sum_jv_j\right|^2,\quad
 w_j=\frac{n(v_j-\mu_v)}{\sqrt{D_n}}.
\end{equation}
Then $D_n$ is a positive integer, $\sum_jw_j=0$,
$\sum_j|w_j|^2=n$, and
\begin{equation}\label{general_eq_lattice_product}
 \Delta_n(w)=\left(\frac{n^2}{D_n}\right)^N
                 \prod_{i<j}|v_i-v_j|^2\in\Q_{>0}.
\end{equation}
Furthermore,
\begin{equation}\label{general_eq_lattice_asymptotic}
 \lim_{n\to\infty}\frac{\log\Delta_n(w)}{n^2}=c,\qquad
 \log M_n-\log\Delta_n(w)=o(n^2).
\end{equation}
\end{theorem}
\begin{proof}
Centering and scaling give the finite identities. Positivity of $D_n$
follows because the points are distinct and $n\ge2$.
Put $R_n=\max_{j\le n}|v_j|$ and $a_n=n/\sqrt{D_n}$.
Covering lattice points by unit squares proves
\[
 \#\{v:|v|\le R\}=\pi R^2+O(R+1),\qquad
 \sum_{|v|\le R}|v|^2=\frac\pi2R^4+O(R^3+1).
\]
Only $O(R_n+1)$ boundary points can be omitted when a radius shell
is cut by the tie rule. The full disk has zero coordinate sum, so
$\sum_{j\le n}v_j=O(R_n^2+1)$ and $\mu_v=O(1)$.
Consequently
\[
 R_n^2\sim\frac n\pi,\quad
 \sum_j|v_j-\mu_v|^2\sim\frac{n^2}{2\pi},\quad
 a_n^2\sim\frac{2\pi}{n},\quad a_nR_n\longrightarrow\sqrt2.
\]
Riemann sums over unit squares show that
$\mu_n=n^{-1}\sum_j\delta_{w_j}$ converges weakly to the uniform
probability measure $\mu$ on $|z|\le\sqrt2$.
Their supports lie in a common bounded disk for large $n$.

The minimum pair distance is at least $a_n$.
Packing lattice squares gives, for $r\ge a_n$, at most
$C(r/a_n)^2\le C'nr^2$ other points within distance $r$ of each point.
Thus the ordered-pair count $H_n(r)$ is zero for $r<a_n$ and
satisfies $H_n(r)\le C'n^2r^2$ for $r\ge a_n$.
For $0<\varepsilon<1$, layer integration gives
\begin{align*}
 \frac1{n^2}\sum_{\substack{i\ne j\\|w_i-w_j|<\varepsilon}}
       -\log|w_i-w_j|
 &=\frac{H_n(\varepsilon)}{n^2}|\log\varepsilon|
       +\frac1{n^2}\int_0^\varepsilon\frac{H_n(r)}r\,dr\\
 &\le C'\varepsilon^2
             \left(|\log\varepsilon|+\frac12\right).
\end{align*}
For fixed $\varepsilon$, the truncated kernel
$\max\{\log|z-u|,\log\varepsilon\}$ is continuous on the common
compact support. Weak convergence gives convergence of its double
integrals, and its diagonal empirical contribution is
$\log\varepsilon/n$.
First let $n\to\infty$ and then $\varepsilon\downarrow0$ to obtain
\[
 \frac{\log\Delta_n(w)}{n^2}
 =\frac1{n^2}\sum_{i\ne j}\log|w_i-w_j|
 \longrightarrow\iint\log|z-u|\,d\mu(z)d\mu(u).
\]
For $R=\sqrt2$, the angular mean identity
\[
 \frac1{2\pi}\int_0^{2\pi}\log|r-se^{i\theta}|\,d\theta
 =\log\max\{r,s\}
\]
gives the disk potential
$\log R-\tfrac12+|z|^2/(2R^2)$ for $|z|\le R$.
Since $\int|z|^2d\mu=R^2/2$, its average is $\log R-1/4=c$.
This proves the limit. The remaining assertion follows from
Corollary~\ref{general_cor_asymptotic} and $\Delta_n(w)\le M_n$.
\end{proof}

\section{Exact multi-ring lower-bound constructions}
Let $\delta\in\{0,1\}$ indicate the presence of the origin.
Ring $a$ consists of $m_a\ge2$ points
$r_a\exp(i(\theta_a+2\pi j/m_a))$, $0\le j<m_a$, with $r_a>0$.
Require
\[
 n=\delta+\sum_a m_a,\qquad \sum_a m_ar_a^2=n.
\]
For two rings define $g_{ab}=\gcd(m_a,m_b)$ and
$L_{ab}=\operatorname{lcm}(m_a,m_b)$.

\begin{proposition}\label{prop:rings}
The exact product of this configuration is
\begin{equation}\label{eq:rings}
\begin{split}
 \Delta_n={}&\prod_a m_a^{m_a}
       r_a^{m_a(m_a-1)+2\delta m_a}\\
 &\quad\cdot\prod_{a<b}
 \left|r_a^{L_{ab}}e^{iL_{ab}\theta_a}
       -r_b^{L_{ab}}e^{iL_{ab}\theta_b}\right|^{2g_{ab}}.
\end{split}
\end{equation}
Consequently this expression is a lower bound for $M_n$.
\end{proposition}
\begin{proof}
Each ring is the root set of
$z^{m_a}-r_a^{m_a}e^{im_a\theta_a}$, whose absolute discriminant
is $m_a^{m_a}r_a^{m_a(m_a-1)}$.
The absolute resultant of $z^m-A$ and $z^k-B$ is
$|A^{k/g}-B^{m/g}|^g$, where $g=\gcd(m,k)$:
the values of $z^k$ on the first root set traverse the roots of
$u^{m/g}-A^{k/g}$, each with multiplicity $g$.
Squaring this resultant gives every cross-ring squared distance.
Each origin-to-ring product contributes $r_a^{2m_a}$.
Multiplying proves the formula, also when cross-ring collisions make
both sides zero.
\end{proof}

With two rings of populations $a,b$ and $\delta$ origin points,
write $g=\gcd(a,b)$, $L=ab/g$, and $x=r^2/s^2$.
The optimal relative phase within this family is
$\theta_2-\theta_1=\pi/L$ modulo $2\pi/L$.
Set
$A=a(a-1)/2+\delta a$ and $B=b(b-1)/2+\delta b$.
Then $N=A+B+ab$, $s^2=n/(ax+b)$, and the family maximum reduces to
\begin{equation}\label{eq:two-rings}
 \max_{x>0} a^ab^bn^N
       \frac{x^A(1+x^{L/2})^{2g}}{(ax+b)^N}.
\end{equation}
The interior critical equation is
\begin{equation}\label{eq:ring-critical}
 Ab-a(B+ab)x+b(A+ab)x^{L/2}-aBx^{L/2+1}=0.
\end{equation}
It follows by logarithmic differentiation and
$gL=ab$. For odd $L$, substituting $x=\rho^2$ makes this a polynomial
with integer exponents, so exact root isolation and value comparison
can certify the optimum within the chosen family.
For three or more rings, the relative phases must be optimized together;
pairwise opposite phases need not be simultaneously possible.
Family optimization alone does not prove unrestricted optimality.

\subsection{A strict ten-point improvement on the central polygon}
The center-plus-$(n-1)$-gon value is
\[
 C_n=(n-1)^{n-1}\left(\frac n{n-1}\right)^N.
\]
For $n=10$, consider instead the roots of
\begin{equation}\label{eq:ten-polynomial}
 (z^2-u)(z^8+v^4),\qquad u=\frac18,\quad v=\frac{39}{32}.
\end{equation}
They are $\pm\sqrt u$ and the vertices of a regular octagon of
radius $\sqrt v$, rotated by $\pi/8$.
They satisfy $2u+8v=10$ and have product
\begin{equation}\label{eq:ten-product}
 D_{10}=4u\,8^8v^{28}(u^4+v^4)^4,\qquad
 \frac{D_{10}}{C_{10}}>\frac{113}{100}>1.
\end{equation}
The last inequality is a rational comparison verified by the
accompanying program. Thus the central polygon cannot be the general
finite-$n$ answer. This two-ring example is a certified lower bound,
and is not asserted to maximize $\Delta_{10}$.

\section{An explicit rational polar construction for every point count}
\label{sec:polar-all-n}

Put $c_*=(\log2)/2-1/4$ and
$\gamma_{\rm pol}=(\log6)/2-1/2$.
The following finite product is an attainable lower bound for every
$n$; it does not assert the unrestricted maximizing geometry.

\begin{theorem}\label{thm:polar-all-n}
For $n\ge2$, let $k$ be the largest nonnegative integer with
$n_0=3k(k+1)\le n$, and put $r=n-n_0$.
For each $1\le j\le k$, take a regular $6j$-gon centered at zero
with squared radius $6j^2/n$ and phase zero.  If $r>0$, add an
outer regular $r$-gon of squared radius $(n+n_0)/n$ and phase zero;
for $r=1$ this means one point on the positive real axis.
These points have quadratic moment $n$ and no collisions.

Write $(m_a,b_a)$ for the population and squared-radius numerator
of each ring, namely $(6j,6j^2)$ and, if present, $(r,n+n_0)$.
For each pair put $g_{ab}=\gcd(m_a,m_b)$ and
$L_{ab}=m_am_b/g_{ab}$.  All $L_{ab}$ are even.  Then
\begin{equation}\label{eq:polar-rational-all-n}
 D_n=\frac1{n^N}
 \prod_a m_a^{m_a}b_a^{m_a(m_a-1)/2}
 \prod_{a<b}(b_b^{L_{ab}/2}-b_a^{L_{ab}/2})^{2g_{ab}}
 \le M_n,\qquad N=\frac{n(n-1)}2.
\end{equation}
In particular $D_n$ is an explicit positive rational number.
Moreover, as $n\to\infty$,
\begin{equation}\label{eq:polar-three-term}
 \log D_n=c_*n^2+\frac n2\log n+\gamma_{\rm pol}n+O(\sqrt n).
\end{equation}
For the subsequence $n=3k(k+1)$ the remainder improves to
$O(\log n)$.
The construction strictly improves the central polygon for every
$n\ge13$:
\[
 D_n>C_n=(n-1)^{n-1}\left(\frac n{n-1}\right)^N.
\]
At $n=18$ it gives the particularly simple bound
\begin{equation}\label{eq:polar-eighteen}
 M_{18}\ge\frac{2^{162}455^{12}}{3^{111}}>970\,C_{18}.
\end{equation}
\end{theorem}

\begin{proof}
The population sum is $n$.  Using
$\sum_{j=1}^k j^3=[k(k+1)/2]^2$, the moment is
\[
 \frac{36}{n}\sum_{j=1}^k j^3+
 \frac{r(n+n_0)}n
 =\frac{n_0^2+r(2n_0+r)}n=n.
\]
The radii increase strictly; the outer radius, if present, is
larger than the last inner radius.  An inner population is even,
so every cross-ring least common multiple is even.  The existing
resultant formula for rings therefore gives
\eqref{eq:polar-rational-all-n}, with integral exponents and an
integer numerator before reduction.  The one-point ring causes
no difficulty.  If desired, centering and renormalizing the
$r=1$ case only increases its product.

We give uniform estimates that prove the asymptotic assertions.
For arbitrary increasing ring radii $u_j=R_j^2$, let
$p_j=m_j/n$ and $s_j=\sum_{i\le j}p_i$.
The present radii satisfy
\[
 u_j=s_{j-1}+s_j,\qquad \sum_jp_ju_j=1.
\]
Factoring the outer radius from every cross-ring resultant gives
the exact identity
\begin{equation}\label{eq:polar-energy-decomposition}
 \log D_n=n^2F+H+T,
\end{equation}
where
\[
 F=\frac12\sum_jp_ju_j\log u_j,
 \qquad H=\sum_jm_j\log(m_j/\sqrt{u_j}),
\]
and
\[
 T=\sum_{a<b}2g_{ab}\log(1-q_{ab}),\qquad
 q_{ab}=(u_a/u_b)^{L_{ab}/2}.
\]
The function $f(s)=s\log(2s)$, with $f(0)=0$, has
$f''(s)=1/s$.  Thus $F$ is its composite midpoint sum on
the partition $s_j$, and
\[
 \int_0^1f(s)\,ds=c_*.
\]

First consider $r=0$, so $n_0=3K$ with $K=k(k+1)$.
Here $p_j=2j/K$, $u_j=2j^2/K$, and
\[
 F_k=\frac12\log(2/K)+\frac4{K^2}\sum_{j=1}^k j^3\log j.
\]
Euler--Maclaurin with the fourth-derivative remainder gives
\begin{equation}\label{eq:polar-em-exact}
 \sum_{j=1}^k j^3\log j
 =\frac{K^2}{4}\log k-\frac{k^4}{16}
  +\frac{k^2}{12}+E_k,
 \quad
 -\frac1{48}-\frac{\log k}{60}
 \le E_k\le-\frac1{48}.
\end{equation}
Indeed, the explicit third-derivative term is
$-(\log k)/120$, while the remainder is bounded in modulus by
$\int_1^k(6/x)\,dx/720=(\log k)/120$.
Writing
\[
 A(t)=\frac12\log(1+t)
 -\frac{t(6+9t+2t^2)}{12(1+t)^2},
\]
one has $A'(t)=-t^3/[6(1+t)^3]$, hence
$-t^4/24\le A(t)\le0$ for $0\le t\le1$.
Substitution into \eqref{eq:polar-em-exact} yields the useful
finite bounds
\begin{equation}\label{eq:polar-midpoint-cost}
 \frac{n_0}{2}-\frac34
 \le n_0^2(c_*-F_k)
 \le\frac{n_0}{2}+\frac34+\frac35\log k.
\end{equation}
Also $H=(n_0/2)\log(6n_0)$ exactly.

The cross-ring correction is bounded independently of $k$.
For inner rings $i<j$, put $d=\gcd(i,j)$ and $h=j-i$.
Then
\[
 q_{ij}=(i/j)^{6ij/d}
 \le e^{-6ih/d}\le e^{-6\max(i,h)},
 \qquad g_{ij}=6d\le6\min(i,h).
\]
The first inequality uses $\log(j/i)\ge(j-i)/j$.
Since $-\log(1-q)\le q/(1-q)$, summing by
$t=\max(i,h)$ and using
$\sum_{\max(i,h)=t}\min(i,h)=t^2$ proves
\begin{equation}\label{eq:polar-cross-constant}
 -B_6\le T\le0,\qquad
 B_6=\frac{12e^{-6}(1+e^{-6})}{(1-e^{-6})^4}<1.
\end{equation}
Together with \eqref{eq:polar-midpoint-cost}, this proves the
three-term expansion with an $O(\log n)$ remainder on the
subsequence.

For general $r>0$, the inner part of the midpoint-error integral
is a scaling of the old partition by $a=n_0/n$.
The affine term introduced by scaling $f(as)$ cancels in midpoint
error.  Consequently its contribution to $n^2(c_*-F)$ is
exactly $n_0^2(c_*-F_k)$.
The last interval has length $p=r/n$ and left endpoint $a$.
Its midpoint error lies between zero and $p^3/(24a)$.
Thus, since $r<6(k+1)$,
\[
 0\le n^2E_{\rm last}\le\frac{r^3}{24n_0}
 <\frac{3(k+1)^2}{k}\le12k.
\]
The angular term now satisfies
\[
 H=\frac n2\log(6n)
 +r\log\frac r{\sqrt{6(n+n_0)}}.
\]
The last term lies between $-6(k+2)/e$ and zero.
For the lower bound use $x\log x\ge-1/e$ and
$\sqrt{6(n+n_0)}\le6(k+2)$.
For the upper bound, the factorization
\[
 r^2-6(n+n_0)=(r-6(k+1))(r+6k)<0
\]
shows $r<\sqrt{6(n+n_0)}$.

For cross terms between inner ring $j$ and the extra ring, put
$U=\sqrt{(n+n_0)/6}$.  Then
$k+2/5\le U<k+1$ and $L\ge6j$.
For $j<U/2$, $q\le2^{-6j}$.
For $j\ge U/2$,
$q\le\exp[-3(k-j+2/5)]$.
With $2g\le12j$, their total logarithmic correction is therefore
at least $-Ak-B$, where
\[
 A=\frac{12e^{-6/5}}
 {(1-e^{-6/5})(1-e^{-3})}<6,
 \qquad
 B=\frac{12/64}{(1-1/64)^3}<1.
\]
Combining all estimates gives, for $n\ge6$,
\begin{equation}\label{eq:polar-uniform-remainder}
 -21k-\log(k+1)-9
 \le\log D_n-c_*n^2-\frac n2\log n-\gamma_{\rm pol}n
 \le3k+4.
\end{equation}
Since $k=O(\sqrt n)$, this proves
\eqref{eq:polar-three-term} for every point count.

Finally $c_*>9/100$, $\gamma_{\rm pol}>7/20$, and
$\log C_n\le n\log n+n/2$.
For $n\ge80$, use $\log n\le n/8$,
$k\le n/12$, and $\log(k+1)\le k$ in
\eqref{eq:polar-uniform-remainder} to obtain
\[
 \log(D_n/C_n)>
 \frac{11n^2}{400}-2n-9>0.
\]
The last polynomial is positive and increasing for $n\ge80$.
The remaining cases $13\le n<80$ are exact rational comparisons
in \texttt{all\_n\_rings\_verify.py}; the verifier also checks
all rational bounds on constants used above.
For $n=18$, the two radii squared are $1/3$ and $4/3$,
with populations six and twelve, giving
\eqref{eq:polar-eighteen} by the resultant formula.
\end{proof}

\subsection{Why the sixfold population slope is selected}
For an integer $c\ge2$, take $n=ck(k+1)/2$, populations $cj$,
and radii $j\sqrt{2/[k(k+1)]}$.
The same calculation yields
\[
 \log D_{c,k}=c_*n^2+\frac n2\log n+
 \left(\frac12\log c-\frac c{12}\right)n+O_c(\log n).
\]
The cross correction is bounded by
$2ce^{-c}(1+e^{-c})/(1-e^{-c})^4$.
The function $(\log c)/2-c/12$ has its unique maximum at $c=6$.
This selects sixfold population growth within this explicit
construction family, without proving that the family itself is
globally optimal.

\subsection{An exact obstruction for two rings with unequal populations}
\begin{proposition}\label{prop:two-ring-population-obstruction}
A physical stationary configuration consisting of exactly two
concentric regular polygons of positive radii, with populations
$a,b\ge2$ and optionally their common center, must have $a=b$.
For $a\ne b$, even the maximum over all radii and relative phases
within this two-ring family is strictly smaller than $M_{a+b+\delta}$.
\end{proposition}

\begin{proof}
Write the polygons as the roots of $z^a-A$ and $z^b-B$,
where $A,B\ne0$, and assume no collisions.
At a point $z$ on the second polygon, multiply the force by $z$:
\[
 z\sum_{w\ne z}\frac1{z-w}
 =\frac{b-1}{2}+\delta+\frac{az^a}{z^a-A}.
\]
Stationarity makes the left side equal to
$\kappa|z|^2$, a constant on this polygon.
The map $t\mapsto at/(t-A)$ is injective away from its pole,
so $z^a$ must be constant on all $b$ points.  Thus $b$ divides $a$.
Interchanging the polygons gives $a$ divides $b$, proving $a=b$.
The family maximum is attained on the compact moment surface;
a positive maximum has neither collapsed rings nor collisions.
If it equalled the unrestricted maximum, its configuration would
be stationary, contradicting the divisibility argument.
\end{proof}

Consequently the eighteen-point construction in
\eqref{eq:polar-eighteen} satisfies the stronger strict inequality
\[
 M_{18}>D_{18}>970C_{18}.
\]
The existing ten-point construction with populations two and eight
is likewise strictly submaximal, even after complete optimization
within that two-ring family.  This obstruction is specific to
two nonzero rings; no analogous assertion is made here for
three or more rings, where additional forces can interact.

\section{A quantitative obstruction to every fixed number of regular rings}
\label{sec:fixed-ring-obstruction}

\begin{proposition}\label{prop:fixed-ring-gap}
Let $R_{n,K}$ be the supremum of the discriminant product over
configurations consisting of at most $K$ concentric regular
polygons and optionally the origin, with quadratic moment $n$.
For every fixed $K$,
\[
 \limsup_{n\to\infty}\frac{\log R_{n,K}}{n^2}
 \le c_*-\frac1{24K^2}<c_*.
\]
Hence no fixed-$K$ regular-ring formula can supply the general
global maximum; its ratio to $M_n$ decays exponentially in $n^2$.
\end{proposition}

\begin{proof}
Order the radii increasingly.  Bounding a cross resultant by
twice its outer-radius power bounds the entire product by
$2^{2\sum_{a<b}g_{ab}}$ times its radial outer-power expression.
For fixed $K$ this factor contributes only $O_K(n)$ to its
logarithm.  The angular factors contribute at most $n\log n$.
For populations $m_j$ and an optional origin indicator $\delta$,
the radial exponent is
\[
 W_j=\frac{m_j}{2}
   \bigl(2\sum_{i<j}m_i+m_j-1+2\delta\bigr),
 \qquad \sum_jW_j=N.
\]
Under $\sum_jm_ju_j=n$, its optimum has
\[
 u_j=\frac{nW_j}{Nm_j}
 =\frac{2\sum_{i<j}m_i+m_j-1+2\delta}{n-1}.
\]
Taking a subsequence of population proportions gives a partition
$0=s_0\le s_1\le\cdots\le s_K=1$, allowing zero intervals.
After division by $n^2$, the radial logarithm tends to
\[
 F=\sum_jp_j f\bigl((s_{j-1}+s_j)/2\bigr),
 \qquad p_j=s_j-s_{j-1},\quad f(s)=s\log(2s).
\]
This remains valid at zero intervals by continuity of
$s\log s$; the finite-$n$ corrections are $o(n^2)$.
Because $f(s)-s^2/2$ is convex on $[0,1]$, Jensen's inequality
on each interval gives
\[
 \int_{s_{j-1}}^{s_j}f(s)\,ds
 -p_jf\bigl((s_{j-1}+s_j)/2\bigr)
 \ge\frac{p_j^3}{24}.
\]
Summing and using $\sum_jp_j^3\ge K^{-2}$ proves
$F\le c_*-1/(24K^2)$.
The already established global asymptotic
$\log M_n/n^2\to c_*$ finishes the comparison.
\end{proof}

For arbitrary global maximizing configurations, the number of
distinct radii about their centroid also tends to infinity,
without an angular regularity assumption.
Indeed, scale a maximizer by
$\sqrt{(n-1)/(2n)}$.  The resulting points maximize the weighted
Vandermonde for the Ginibre potential $Q(z)=|z|^2$.
The classical weighted Fekete equidistribution theorem recalled
in Section~1.2 of Ameur--Ortega-Cerd\`a~\cite{AOC}
(Theorem~1.1 of their arXiv version) implies
convergence back in the original scale to uniform area measure
on the disk of radius $\sqrt2$.
Their squared-radius distributions therefore converge to the
uniform law on $[0,2]$.
Probability measures supported on at most $K$ radii in a fixed
compact interval form a weakly closed set, so bounded $K$ would
contradict this continuous limit.  The explicit gap above is
stronger within the regular-ring family.

\section{Exact six-point subclasses and a majorant obstruction}
Write $C_6=6^{15}/5^{10}$, the value of the origin together with a
regular pentagon of radius $\sqrt{6/5}$. These subclass results
complement Theorem~\ref{thm:six-global}.
The unrestricted six-point value is now certified by the exhaustive
algebraic argument in Section~\ref{sec:six-global-exact}.

\begin{proposition}[Configurations containing the origin]
If $\sum_{i=1}^6|z_i|^2=6$ and some $z_i=0$, then $\Delta_6\le C_6$.
Equality holds precisely for the origin and a regular pentagon of
radius $\sqrt{6/5}$, up to rotation and permutation.
\end{proposition}
\begin{proof}
Assume $z_6=0$. Apply the five-point theorem of Wang \cite{Wang} to
$\sqrt{5/6}(z_1,\ldots,z_5)$, and arithmetic--geometric mean to
the five squared radii. Homogeneity gives
\[
 \Delta_6=\Delta_5(z_1,\ldots,z_5)\prod_{i=1}^5|z_i|^2
 \le5^5(6/5)^{10}(6/5)^5=C_6.
\]
The equality conditions give exactly the asserted configuration.
\end{proof}

\begin{proposition}[Two concentric equilateral triangles]
Among configurations consisting of two equilateral triangles centered
at the origin with $\sum_i|z_i|^2=6$, the largest product is $6^6$,
attained precisely by a unit regular hexagon.
\end{proposition}
\begin{proof}
Let $a,b\ge0$ be the squared radii, so $a+b=2$.
If $ab=0$ there are collisions. Otherwise the ring discriminants are
$3^3a^3$ and $3^3b^3$. Their resultant gives the cross-distance product
$|a^{3/2}-b^{3/2}e^{3i\phi}|^6$, where $\phi$ is the relative rotation.
Thus, with $p=\sqrt{ab}\in(0,1]$,
\[
 \Delta_6\le3^6(ab)^3(a^{3/2}+b^{3/2})^6
 =3^6\{8p^2-6p^4+2p^5\}^3.
\]
The derivative of the expression in braces is
$2p(8-12p^2+5p^3)>0$: the last factor decreases on $[0,1]$
and equals $1$ at $p=1$. Its maximum is $4$, hence
$\Delta_6\le3^6\cdot4^3=6^6$.
Equality requires $a=b=1$ and $3\phi=\pi$ modulo $2\pi$,
giving a regular hexagon. Finally $6^6<C_6$ because
$C_6/6^6=6^9/5^{10}>1$.
\end{proof}

\subsection{Hermite majorants adapted to a candidate}
Suppose a candidate has $r$ distinct positive squared distances
$a_1,\ldots,a_r$. Let $H$ be the Hermite polynomial of degree at
most $2r-1$ satisfying $H(a_j)=\log a_j$ and $H'(a_j)=1/a_j$.
Then
\[
 \log x\le H(x)\quad(x>0),\qquad
 \log\Delta_n\le\sum_{i<j}H(|z_i-z_j|^2).
\]
The interpolation remainder is a nonpositive multiple of
$P(x)=\prod_j(x-a_j)^2$, since
$(\log x)^{(2r)}=-(2r-1)!/x^{2r}<0$.
The sum of $H$ is a finite combination of distance moments;
stationary identities and second-order Gram matrices supply
necessary constraints on those moments.

If an upper endpoint $R>\max_j a_j$ is available, set
$\kappa_R=(H(R)-\log R)/P(R)>0$.
Adding the interpolation node $R$ gives
$\log x\le H(x)-\kappa_RP(x)$ for $0<x\le R$:
the next derivative is positive and the remainder has the sign of
$P(x)(x-R)$. Any verified $0<k\le\kappa_R$ also gives a valid majorant.

For the central pentagon the nodes are
\[
 a_1=6/5,\qquad
 a_2=3(5-\sqrt5)/5,\qquad a_3=3(5+\sqrt5)/5,
\]
each occurring five times.
A regular hexagon has squared-distance multiplicities $6,6,3$ at
$1,3,4$, respectively. The exact interval check
\texttt{verify\_hermite.py} shows that its degree-five majorant satisfies
\[
 \frac2{25}<6H(1)+6H(3)+3H(4)-\log C_6<\frac{83}{1000}.
\]
Even the maximal interpolation correction at $R=15/2$ leaves
this gap greater than $3/100$.
Thus this majorant alone is insufficient. The separate certified
comparison in Section~\ref{sec:six-global-exact} supplies the
unrestricted optimum.

\section{Local certificates and explicit improving perturbations}
\label{sec:local-certificates}
The spectral stability of polygonal vortex equilibria has a substantial
classical literature; see Cabral--Schmidt \cite{CS}.
We give direct constrained-discriminant calculations, with explicit
Taylor bounds and an exact treatment of a degenerate ten-point case.
We make no claim of priority for the polygonal stability thresholds.

\subsection{The complete central-polygon spectrum}
Write $n=m+1$, $m\ge4$, $r^2=n/m$, $\omega=e^{2\pi i/m}$, and
\[
 z^*=(0,r,r\omega,\ldots,r\omega^{m-1}),\qquad \lambda=m/2.
\]
The root polynomial $z(z^m-r^m)$ or a geometric sum verifies
\[
 \sum_{j\ne i}\frac1{z_i^*-z_j^*}=\lambda\overline{z_i^*}.
\]
Define the complex symmetric matrix
\[
 B_{ij}=-\frac1{(z_i^*-z_j^*)^2}\quad(i\ne j),\qquad
 B_{ii}=\sum_{j\ne i}\frac1{(z_i^*-z_j^*)^2}.
\]
Let $\langle x,y\rangle=\sum_i\overline{x_i}y_i$, and put
\[
 E=\{h\in\C^n:\ \sum_i h_i=0,\ \langle z^*,h\rangle=0\}.
\]
This is a local slice for the zero-centroid sphere modulo rotation.
Its normalized chart is
\begin{equation}\label{eq:local-chart}
 z(h)=\frac{\sqrt n\,(z^*+h)}{\sqrt{n+\|h\|^2}}.
\end{equation}
The Taylor expansion of
$F(h)=\log\Delta_n(z(h))-\log C_n$ has quadratic term
\begin{equation}\label{eq:quadratic-slice}
 Q_2(h)=-\Re(h^{\mathsf T}Bh)-\lambda\|h\|^2.
\end{equation}
Indeed, expand each $2\Re\log(1+(h_i-h_j)/(z_i^*-z_j^*))$
and subtract $N\log(1+\|h\|^2/n)$; the linear term vanishes
by stationarity and $h\in E$.

Set
\[
 f_k=\frac1{\sqrt m}(0,1,\omega^k,\ldots,\omega^{(m-1)k}),
 \quad 1\le k<m,\qquad
 u=\frac{(-m,1,\ldots,1)}{\sqrt{m(m+1)}}.
\]
Together with $\mathbf1/\sqrt n$, these form a Hermitian orthonormal
basis. Direct divided-difference expansion and the identity
$\sum_{j=0}^{m-1}\omega^{jk}=0$ for $m\nmid k$ give
\begin{align}
 B\mathbf1&=0,& Bf_1&=\lambda f_{m-1},\notag\\
 Bf_2&=\gamma u,& Bu&=\gamma f_{m-2},\label{eq:central-spectrum-map}\\
 Bf_k&=c_kf_{k-2}\quad(3\le k<m),&
 \gamma&=\frac m{\sqrt{m+1}},\quad
 c_k=\frac{1-(k-2)(m-k)/2}{r^2}.\notag
\end{align}
For example, on ring coordinates and $3\le k<m$, the expansion is
\[
 (B(z^*)^{\circ k})_i
 =\bigl(1-(k-2)(m-k)/2\bigr)(z_i^*)^{k-2};
\]
the center coordinate is zero, since its corresponding power sum
vanishes. For $k=2$, the center coordinate supplies the vector $u$.
The output vectors in \eqref{eq:central-spectrum-map} are distinct
orthonormal basis vectors, so the singular values of $B$ are
\begin{equation}\label{eq:central-singular-values}
 0,\quad\lambda,\quad\gamma,\quad\gamma,
 \quad |c_k|\ (3\le k<m).
\end{equation}
On $E$ the singular value $\lambda$ associated with $f_1$
has been removed. Since
\[
 \max_{3\le k<m}(k-2)(m-k)
 =\left\lfloor\frac{(m-2)^2}{4}\right\rfloor,
\]
the nontrivial singular values are strictly below $\lambda$
for $4\le m\le8$. Thus the central polygon is a strict local
maximum modulo rotations for $5\le n\le9$.
For $m=9$, $c_5=c_6=-\lambda$; the remaining singular values on $E$
are strictly smaller. For $m\ge10$, some $|c_k|>\lambda$.
An explicit improving direction for that case is given below.

\subsection{An explicit six-point neighborhood certificate}
\begin{theorem}\label{thm:six-basin}
Take $n=6$ and $z^*$ as above. For $h\in E$ with
$0<\|h\|\le1/25$, the chart \eqref{eq:local-chart} satisfies
\begin{equation}\label{eq:six-basin}
 \log\Delta_6(z(h))\le\log C_6-\frac14\|h\|^2.
\end{equation}
\end{theorem}
\begin{proof}
For $m=5$, $c_3=c_4=0$, and the norm of $B$ restricted to $E$
is $\gamma=5/\sqrt6$. Hence \eqref{eq:quadratic-slice} gives
$Q_2(h)\le-\delta\|h\|^2$, where
$\delta=5/2-5/\sqrt6>9/20$.
The last strict inequality follows from $25/6<(41/20)^2$.

Write $\rho=\|h\|$ and
$v_{ij}=(h_i-h_j)/(z_i^*-z_j^*)$.
The smallest squared distance at $z^*$ is $r^2=6/5$, so
\[
 \max_{i<j}|v_{ij}|\le\sqrt{5/3}\,\rho<\frac43\rho=:b_\rho,
 \qquad
 \sum_{i<j}|v_{ij}|^2\le\frac6{r^2}\rho^2=5\rho^2.
\]
Here the second inequality uses
$\sum_{i<j}|h_i-h_j|^2=6\|h\|^2$ for $\sum_i h_i=0$.
For $|v|<1$,
\[
 |\log(1+v)-v+v^2/2|
 \le\frac{|v|^3}{3(1-|v|)}.
\]
Also $\log(1+x)\ge x-x^2/2$ for $x\ge0$.
Applying these bounds to the exact chart product gives
\[
 F(h)\le-\delta\rho^2+
 \left(\frac{10b_\rho}{3(1-b_\rho)}+\frac5{24}\rho^2\right)\rho^2.
\]
For $\rho\le1/25$, the coefficient in parentheses is at most
$40/213+1/3000<1/5$. Consequently
$F(h)<-(9/20-1/5)\rho^2=-\rho^2/4$.
This proves the stated non-strict bound.
\end{proof}

This theorem excludes a concrete neighborhood from a search for a
strictly better six-point competitor. It makes no comparison with
configurations outside that neighborhood.

\subsection{Ten points: a strict local maximum that is not global}
\begin{theorem}\label{thm:ten-local-trap}
For $n=10$, the origin plus a regular nine-gon of radius $\sqrt{10/9}$
is a strict local maximum modulo rotations. It is not a global maximum.
\end{theorem}
\begin{proof}
The failure of global optimality is already certified by
\eqref{eq:ten-product}. We prove strict local maximality.
Now $m=9$, $r=\sqrt{10/9}$ and $\omega=e^{2\pi i/9}$.
The only null directions of $Q_2$ on $E$ are
\[
 h(a)_0=0,\qquad
 h(a)_{j+1}=r(a\omega^{5j}+\overline a\,\omega^{6j}),
 \qquad a\in\C.
\]
In fact, on the two relevant Fourier modes, writing their
coefficients as $a,b$ gives $Q_2=-45|a-\overline b|^2$.
Let $W\subset E$ be the real orthogonal complement of this
two-dimensional nullspace. On $W$ write
$Q_2(v)=-(v,Av)_{\R}$, with $A$ positive definite.

For $h=h(1)$ put
$u_{ij}=(h_i-h_j)/(z_i^*-z_j^*)$ and
$a_i=\sum_{j\ne i}u_{ij}^2/(z_i^*-z_j^*)$.
Finite root-of-unity sums give
\begin{align*}
 \sum_{i<j}u_{ij}&=0,&\sum_{i<j}u_{ij}^2&=-90,\\
 \sum_{i<j}u_{ij}^3&=0,&\sum_{i<j}u_{ij}^4&=2970,\\
 a_0&=-9/r,&
 a_{j+1}&=(1-15\omega^{7j}-10\omega^{8j})/r.
\end{align*}
These finite identities are also checked exactly in
$\Q(\omega)$ with $\omega^6+\omega^3+1=0$ by
\texttt{verify\_polygon\_\allowbreak perturbations.py}.
The first four sums can independently be obtained by expanding the
divided differences and using
\[
 \sum_{i\ne j}\omega^{pi+qj}
 =m^2\mathbf1_{m\mid p}\mathbf1_{m\mid q}
       -m\mathbf1_{m\mid p+q}.
\]

The straight null-direction quartic term is
$-2970/2+45\cdot2^2/2=-1395$.
For a transverse correction $th+t^2k$, $k\in W$, the coefficient
of $t^4$ is
\begin{equation}\label{eq:ten-effective-quartic}
 -1395+Q_2(k)+L(k),\qquad
 L(k)=2\Re\sum_i a_i k_i.
\end{equation}
The optimizing correction is
\[
 (k_0)_0=-12r,\qquad
 (k_0)_{j+1}=r(4/3-(17/3)\omega^{2j}).
\]
It lies in $W$. The projected real gradient of $L$ on $E$ is
$2\overline a+18z^*$, and direct substitution in
\eqref{eq:central-spectrum-map} gives
$2Ak_0=2\overline a+18z^*$.
Positive definiteness on $W$ therefore makes $k_0$ the unique maximizer
of $Q_2(k)+L(k)$ there. Writing $k=r(-9c,c+d\omega^{2j})$ reduces the equation for the
correction to $5c+\overline d=1$, $5d+10\overline c=-15$, hence
$c=4/3$, $d=-17/3$. Fourier orthogonality gives
$\|k\|^2=r^2(90c^2+9d^2)$,
$k^TBk=180cd$, and $L(k)=180c-270d$ at this real solution.
The exact constants are
\[
 \|k_0\|^2=\frac{4490}{9},\quad
 \Re(k_0^{\mathsf T}Bk_0)=-1360,\quad
 Q_2(k_0)=-885,\quad L(k_0)=1770.
\]
Thus the largest effective quartic coefficient is
$-1395+885=-510<0$.

For completeness, the negative quartic coefficient controls every
null direction, rather than only $h(1)$.
The cyclic relabeling combined with physical rotation that fixes
$z^*$ acts by $a\mapsto\omega^{-4}a$ on the nullspace.
As this phase has order nine, its invariant homogeneous cubic
polynomial is zero and its invariant homogeneous quartic polynomial
is a multiple of $|a|^4$.
The implicit function theorem on $W$ gives a unique nearby transverse
maximizer $w(a)=O(|a|^2)$ for each small $a$.
Its reduced value is therefore
\[
 F(h(a)+w(a))=-510|a|^4+O(|a|^5).
\]
It is strictly negative for small $a\ne0$.
Continuity preserves strict concavity in the transverse variables,
so any other nearby transverse choice has smaller value.
This proves strict local maximality on the normalized slice.
Nonzero centroids strictly decrease the normalized value relative
to their centered configuration, so the conclusion also holds on
the full constraint, modulo rotations.
\end{proof}

This exact local trap explains why even careful local searches and
second-order tests cannot alone certify the finite global maximum.

\subsection{A certified improvement for every \texorpdfstring{$n\ge11$}{n >= 11}}
\begin{theorem}\label{thm:all-n-improvement}
Let $m=n-1\ge10$, $k=\lfloor(m+2)/2\rfloor$,
$q=\lfloor(m-2)^2/4\rfloor$, and
\[
 b=\begin{cases}1&m\text{ even},\\2&m\text{ odd},\end{cases}
 \qquad A_m=\frac{bm}{2}(q-m-3)>0,\qquad H=m+2.
\]
Choose the explicit rational step
\[
 t_m=\min\left\{\frac1{2H},\frac{A_m}{4NH^3}\right\}.
\]
Set $h_0=0$ and, on the ring,
\[
 h_{j+1}=\begin{cases}
 r\omega^{kj},&m\text{ even},\\
 r(\omega^{kj}+\omega^{(m+2-k)j}),&m\text{ odd}.
 \end{cases}
\]
Then the explicit configuration
\[
 y_j=\frac{z_j^*+t_mh_j}{\sqrt{1+bt_m^2}}
\]
has zero centroid, quadratic moment $n$, distinct points, and
\begin{equation}\label{eq:all-n-strict-gain}
 \log\Delta_n(y)-\log C_n\ge\frac23 A_m t_m^2>0.
\end{equation}
Consequently $M_n>C_n$ for every $n\ge10$, using the separate
ten-point construction when $n=10$.
\end{theorem}
\begin{proof}
Fourier orthogonality gives $\sum h_j=0$,
$\langle z^*,h\rangle=0$, and $\|h\|^2=bn$.
For even $m$, $2k-2=m$, and \eqref{eq:central-spectrum-map} gives
$h^{\mathsf T}Bh=m(1-q/2)$.
For odd $m$, the two modes pair and this value is doubled.
Thus the quadratic coefficient along the normalized direction is
\[
 -\Re(h^{\mathsf T}Bh)-Nb=A_m.
\]
To see its positivity, for even $m=2v\ge10$ one has
$q-m-3=v^2-4v-2>0$; for odd $m=2v+1\ge11$ it is
$(v-4)(v+1)>0$.

Put $u_{ij}=(h_i-h_j)/(z_i^*-z_j^*)$.
Each ring divided difference is a sum of $k$ unit-modulus terms,
or two such sums with total $m+2$ terms. Center-to-ring ratios
are bounded by $b$. Hence $|u_{ij}|\le H$.
The chosen step gives $Ht_m\le1/2$, so every factor
$1+t_mu_{ij}$ is nonzero.
The logarithmic Taylor bound and $\log(1+bt^2)\le bt^2$ give
\[
 \log\Delta_n(y)-\log C_n
 \ge A_mt_m^2-\frac{2NH^3t_m^3}{3(1-Ht_m)}
 \ge A_mt_m^2-\frac43NH^3t_m^3
 \ge\frac23 A_mt_m^2>0.
\]
This also verifies distinctness and proves the assertion.
\end{proof}

The gain bound is deliberately conservative. Its role is to certify
one explicit strictly improving configuration for every point count
in this range, rather than to approximate the unrestricted maximum.

\section{Finite-point numerical candidates}
The optional program \texttt{explore\_candidates.py} uses NumPy,
analytic gradients, BFGS updates and an Armijo line search to minimize
\[
 \frac{n-1}{2}\sum_i|z_i|^2-\log\Delta_n(z).
\]
Its stationary scale is $\sum_i|z_i|^2=n$.
For each $n=6,\ldots,12$, it uses 50 Gaussian random starts and
one center-plus-polygon start, with seed 20261004.
The table shows the largest observed logarithmic product.
These floating-point searches do not certify global optimality.

\begin{center}
\begin{tabular}{rrr}
\toprule
$n$ & Best observed $\log\Delta_n$ & Log gain over $C_n$\\
\midrule
6 & 10.7820129141 & $<10^{-12}$\\
7 & 13.9877210917 & $<10^{-12}$\\
8 & 17.3602500369 & $<10^{-12}$\\
9 & 20.8757216171 & $<10^{-12}$\\
10 & 24.7089672032 & 0.192722803\\
11 & 28.7906758481 & 0.522765029\\
12 & 33.1914966077 & 1.071897730\\
\bottomrule
\end{tabular}
\end{center}

Each best coordinate set is also rounded to Gaussian integers on
a grid of spacing $10^{-6}$, then centered and normalized using
\eqref{general_eq_lattice}, without any lattice-selection assumption.
The exact product formula \eqref{general_eq_lattice_product} applies
to any distinct Gaussian integers. The corresponding rationalized
ten-point candidate strictly exceeds $(121/100)C_{10}$.
Its integer coordinates, normalization denominator and exact rational
product are saved in \texttt{numerical\_candidates.json}.
This last comparison is an exact attainable lower bound; it does
not convert the numerical search into a global-optimality proof.
Gradient residuals and convergence flags are also recorded, including
searches that did not reach the requested gradient tolerance.


\section{Certified approximation of the global maximum}
\label{sec:grid-global}
A grid maximum becomes useful for the unrestricted problem only when
its discretization error is bounded uniformly. At an actual maximum,
stationarity cancels the first-order rounding error.

\begin{lemma}\label{lem:global-separation}
Every maximizer satisfies
\begin{equation}\label{eq:global-separation}
 \sum_{j\ne i}\frac1{z_i-z_j}=\lambda\overline z_i,
 \qquad \lambda=\frac{n-1}{2},\qquad
 |z_i-z_j|^2\ge\frac{2n}{(n-1)^2}.
\end{equation}
\end{lemma}
\begin{proof}
The multiplier follows by multiplying the stationary equations by
$z_i$ and summing, since each pair contributes one.
Homogeneity shows that a constrained global maximizer also maximizes
$\log\Delta_n(z)-\lambda\sum_i|z_i|^2$ without the constraint:
for each nonzero shape, the optimal scale has quadratic moment $n$.
Let $B$ have off-diagonal entries $-(z_i-z_j)^{-2}$ and diagonal
entries $\sum_{j\ne i}(z_i-z_j)^{-2}$.
The quadratic variation is
$-\operatorname{Re}(h^TBh)-\lambda\|h\|^2$.
Varying the phase of $h$ gives $|h^TBh|\le\lambda\|h\|^2$;
Takagi factorization of the complex symmetric matrix implies
$\|B\|\le\lambda$.
Since $B\mathbf1=0$ and $B^T=B$, projecting both coordinate vectors
onto $\mathbf1^\perp$ gives
\[
 |B_{ij}|\le\lambda\left(1-\frac1n\right)
             =\frac{(n-1)^2}{2n}.
\]
This proves the distance bound.
\end{proof}

For a positive integer $q$, put $R=\lceil q\sqrt n\rceil$ and consider
all $n$-element subsets of
\[
 \mathcal G_{n,q}=\{a+ib:\ a,b\in\mathbb Z,\ |a|,|b|\le R\}.
\]
For each such subset $v=(v_1,\ldots,v_n)$, define
\[
 D(v)=n\sum_i|v_i|^2-\left|\sum_iv_i\right|^2,\qquad
 z(v)_i=\frac{n(v_i-\mu_v)}{\sqrt{D(v)}},\qquad
 \mu_v=\frac1n\sum_i v_i.
\]
As in the lattice construction, these points have moment $n$ and their product is
rational. Let
\begin{equation}\label{eq:grid-maximum}
 B_n(q)=\max_{\substack{v\subset\mathcal G_{n,q}\\|v|=n}}
 \left(\frac{n^2}{D(v)}\right)^N
 \prod_{i<j}|v_i-v_j|^2.
\end{equation}
The maximum in this finite set can be found by integer comparisons.

\begin{theorem}\label{thm:grid-global}
For every integer $q\ge2(n-1)$, with
$K_n=n(n-1)(2n-1)/4$, one has
\begin{equation}\label{eq:grid-global}
 B_n(q)\le M_n\le B_n(q)\exp(K_n/q^2).
\end{equation}
If also $q^2>K_n$, the entirely rational enclosure
\begin{equation}\label{eq:grid-rational}
 B_n(q)\le M_n\le\frac{B_n(q)}{1-K_n/q^2}
\end{equation}
holds. In particular, for $0<\varepsilon<1$, choosing
$q\ge\max\{2(n-1),\sqrt{K_n/\varepsilon}\}$ constructs a configuration
with product at least $e^{-\varepsilon}M_n$.
\end{theorem}
\begin{proof}
Fix a centered global maximizer $z^*$. Round the real and imaginary
parts of each $qz_i^*$ to nearest integers, writing $v_i/q=z_i^*+e_i$.
Because $|z_i^*|\le\sqrt n$, this lies in the stated grid.
Let $c=e-n^{-1}(\sum_i e_i)\mathbf1$ and $\rho^2=\|c\|^2$.
Then
\[
 \rho^2\le\frac{n}{2q^2},\qquad
 v_{ij}:=\frac{c_i-c_j}{z_i^*-z_j^*},\qquad
 |v_{ij}|\le b:=\frac{n-1}{\sqrt n\,q}\le\frac12.
\]
Thus the rounded points are distinct. Since $c$ is centered,
$\sum_{i<j}|c_i-c_j|^2=n\rho^2$, so
\[
 \sum_{i<j}|v_{ij}|^2\le\frac{(n-1)^2}{2}\rho^2.
\]
The convergent logarithm series gives
\[
 2\operatorname{Re}\log(1+v)\ge
 2\operatorname{Re}v-\frac{|v|^2}{1-b}\quad (|v|\le b<1).
\]
Put $a=\operatorname{Re}\sum_i\overline z_i^*c_i$.
Stationarity yields $2\operatorname{Re}\sum_{i<j}v_{ij}=(n-1)a$.
The moment before normalization is $n+2a+\rho^2$.
Using $\log(1+x)\le x$ for $x>-1$ in the normalization term, the
linear contributions cancel exactly, and hence
\begin{align*}
 \log\Delta_n(z(v))-\log M_n
 &\ge-\left(\frac{n-1}{2}+
       \frac{(n-1)^2}{2(1-b)}\right)\rho^2\\
 &\ge-\left(\frac{n-1}{2}+(n-1)^2\right)\frac{n}{2q^2}
   =-\frac{K_n}{q^2}.
\end{align*}
The finite grid maximum is at least this rounded value; its own
configuration is feasible, giving both sides of
\eqref{eq:grid-global}. For $0\le x<1$, $e^x\le(1-x)^{-1}$,
which proves \eqref{eq:grid-rational}.
\end{proof}

The guaranteed error concerns a \emph{complete} grid maximum.
Using only a searched subset provides a lower bound and cannot be
substituted for $B_n(q)$ in the upper bound. With the above choice of
$q$, the integer coordinates have height $O(n^2/\sqrt\varepsilon)$,
but the number of subsets is exponential in $n$. High-accuracy
six-point grid enumeration has not been carried out.
The standard-library checker completely enumerates $1176$ subsets
for $n=2,q=2$, giving $B_2(2)=4$, and $1\,873\,200$ subsets for
$n=3,q=4$, giving
\begin{equation}\label{eq:grid-three}
 B_3(4)=\frac{49280427}{1825346}<27.
\end{equation}
An attaining integer triple is $(-7-7i,-7+7i,5)$, with $D=582$.
The corresponding grid enclosure has logarithmic gap at most $15/32$.

% Insert after the original global-grid theorem.  This strengthens its error
% bound without changing B_n(q) or the requirement of COMPLETE enumeration.
\subsection{A sharper grid certificate from the full stability matrix}

The entrywise separation estimate can be supplemented by the entire
second-variation bound.  This controls the quadratic rounding term directly,
and reserves absolute-value estimates for terms of order three and higher.

\begin{theorem}\label{thm:grid-sharper}
For $n\ge2$ and an integer $q\ge n$, put
\[
 b=\frac{n-1}{\sqrt n\,q},\qquad
 E_n(q)=\frac{n(n-1)}{2q^2}
       +\frac{n^{1/4}(n-1)^{5/2}}{3(1-b)q^3},\qquad
 \overline E_n(q)=\frac{n(n-1)}{2q^2}+\frac{35n^3}{72q^3}.
\]
Then
\[
 B_n(q)\le M_n\le B_n(q)e^{E_n(q)},\qquad
 E_n(q)<\overline E_n(q).
\]
If a rational number $\eta$ satisfies $E_n(q)\le\eta<1$, then
\[
 M_n\le\frac{B_n(q)}{1-\eta}.
\]
For $0<\varepsilon\le1$, the choice
\[
 q\ge\frac n{\sqrt\varepsilon}
\]
guarantees $E_n(q)<71\varepsilon/72<\varepsilon$.  Consequently the
integer coordinates needed for relative logarithmic accuracy $\varepsilon$
have height $O(n^{3/2}/\sqrt\varepsilon)$.
\end{theorem}

\begin{proof}
At a centered global maximizer $z^*$, use the matrix $B$ in
Lemma~\ref{lem:global-separation} and put $\lambda=(n-1)/2$.
The relations $\|B\|\le\lambda$ and $B\mathbf1=B^*\mathbf1=0$ imply
\[
 \sum_{j\ne i}\frac1{|z_i^*-z_j^*|^4}
 \le\sum_j|B_{ij}|^2
 \le\lambda^2\left(1-\frac1n\right).
\]
H\"older's inequality therefore gives, with $a_{ij}=|z_i^*-z_j^*|^{-3}$,
\[
 \sum_{j\ne i}a_{ij}
 \le(n-1)^{1/4}
      \left[\lambda^2\left(1-\frac1n\right)\right]^{3/4}
 =\frac{\lambda^{3/2}(n-1)}{n^{3/4}}=:s.
\]
For any $c\in\mathbb C^n$,
\[
 \sum_{i<j}a_{ij}|c_i-c_j|^2
 \le2\sum_i|c_i|^2\sum_{j\ne i}a_{ij}
 \le2s\|c\|^2.
\]

Round $qz^*$ and center the rounding errors exactly as in
Theorem~\ref{thm:grid-global}; write $\rho^2=\|c\|^2\le n/(2q^2)$
and $v_{ij}=(c_i-c_j)/(z_i^*-z_j^*)$.  Coordinate rounding gives
$|c_i-c_j|\le\sqrt2/q$ and $|v_{ij}|\le b<1$.  Hence
\[
 \sum_{i<j}|v_{ij}|^3
 \le\frac{\sqrt2}{q}\sum_{i<j}a_{ij}|c_i-c_j|^2
 \le\frac{(n-1)^{5/2}}{n^{3/4}q}\rho^2.
\]
For $|v|\le b$, the absolutely convergent logarithm series yields
\[
 2\operatorname{Re}\log(1+v)
 \ge2\operatorname{Re}v-\operatorname{Re}(v^2)
      -\frac{2|v|^3}{3(1-b)}.
\]
The sum of the quadratic terms is
\[
 \sum_{i<j}v_{ij}^2=c^{\mathsf T}Bc,
 \qquad|c^{\mathsf T}Bc|\le\lambda\rho^2.
\]
Put $a=\operatorname{Re}\sum_i\overline z_i^*c_i$.  Stationarity
gives $2\operatorname{Re}\sum_{i<j}v_{ij}=(n-1)a$, while the
normalization term satisfies
\[
 -N\log\left(1+\frac{2a+\rho^2}{n}\right)
 \ge-(n-1)a-\lambda\rho^2.
\]
Thus the linear terms cancel, and
\[
 \log\Delta_n(z(v))-\log M_n
 \ge-\left[(n-1)
     +\frac{2(n-1)^{5/2}}{3(1-b)n^{3/4}q}\right]\rho^2
 \ge-E_n(q).
\]
The complete grid maximum is at least this rounded value, proving
the exponential enclosure.  The rational enclosure follows from
$e^x\le(1-x)^{-1}$ for $0\le x<1$.

The assumed lower bound on $q$ gives
\[
 b\le\frac{n-1}{n^{3/2}}\le\frac2{3\sqrt3}<\frac25.
\]
The middle inequality follows by maximizing $(t-1)/t^{3/2}$ on
$t\ge2$, whose derivative is $(3-t)/(2t^{5/2})$.
Also $n^{-1/4}\le2^{-1/4}<7/8$.  Therefore
\[
 \frac{n^{1/4}(n-1)^{5/2}}{3(1-b)}
 <\frac{35n^3}{72},
\]
which proves the stated entirely rational upper bound on $E_n(q)$.
If $q\ge n/\sqrt\varepsilon$, then
\[
 \overline E_n(q)
 \le\frac\varepsilon2+\frac{35}{72}\varepsilon^{3/2}
 \le\frac{71}{72}\varepsilon.
\]
Finally $\lceil q\sqrt n\rceil=O(n^{3/2}/\sqrt\varepsilon)$ for
the displayed choice of $q$.
\end{proof}

The Taylor argument itself only requires $b<1$, so the exponential
enclosure with $E_n(q)$ remains valid whenever $q>(n-1)/\sqrt n$.

For a concrete comparison, when $n=6,q=10$ the previous bound was
$K_6/q^2=33/40=0.825$.  The new exponent satisfies
\[
 E_6(10)=\frac3{20}
 +\frac{25\sqrt5\,6^{1/4}}{3000(1-1/(2\sqrt6))}
 <\frac{7113}{38000}<\frac3{16}.
\]
Indeed $\sqrt5<9/4$, $6^{1/4}<157/100$, and
$1/(2\sqrt6)<5/24$ give the first rational inequality by integer
comparisons.  Consequently a \emph{complete} computation of $B_6(10)$
would certify
\[
 B_6(10)\le M_6\le\frac{16}{13}B_6(10),
\]
whereas the earlier rational enclosure had factor $40/7$.
The new exponent is approximately $0.1866434402$; this decimal is
only a display value.  No complete six-point enumeration is asserted.


\subsection{Radial moment constraints and a finite spectral formulation}

The stability matrix gives necessary conditions in every dimension without
imposing reflection or polygon symmetry.  Set $N=n(n-1)/2$ and use the
force-$1$ normalization
\[
 \sum_{j\ne i}\frac1{w_i-w_j}=\overline w_i,
 \qquad \sum_iw_i=0,\qquad \sum_i|w_i|^2=N.
\]
Write $t_i=|w_i|^2$, $u_k=\sum_i t_i^k$, $s_k=\sum_iw_i^k$,
and $s_0=n$.  The complex symmetric inverse-square Laplacian is
\[
 B_{ij}=-(w_i-w_j)^{-2}\quad(i\ne j),\qquad
 B_{ii}=\sum_{j\ne i}(w_i-w_j)^{-2}.
\]
It satisfies $B\mathbf1=0$ and $Bw=\overline w$.
At a local maximum on the moment sphere its operator norm satisfies
$\|B\|\le1$, by the second-variation argument in
Lemma~\ref{lem:global-separation}.  The argument applies to a local
constrained maximum because, for each nearby shape, the best scale for
$\log\Delta_n(w)-\sum_i|w_i|^2$ has moment $N$.

\begin{proposition}\label{prop:general-radial-stability}
Every local maximum, and hence every global maximum, satisfies
\begin{align}
 3u_2+\frac{|s_2|^2}{n}&\le n(n-1)^2,
 \label{eq:general-fourth-cut}\\
 8u_3-6(2n-3)u_2+(2n-3)^2N+\frac{|s_3|^2}{n}&\le0.
 \label{eq:general-sixth-cut}
\end{align}
In the original normalization $z_i=\sqrt{2/(n-1)}\,w_i$,
\[
 \sum_i|z_i|^4\le\frac{4n}{3},\qquad
 \max_i|z_i|^2\le1+\sqrt{\frac{n-1}{3}}.
\]
\end{proposition}

\begin{proof}
Divided differences, the force equations, and $s_1=0$ give
\[
 B(w_i^2)_i=(2t_i-(n-1))_i,
 \qquad
 B(w_i^3)_i=(3w_it_i-(2n-3)w_i)_i.
\]
Since $B$ annihilates constants and is a contraction,
\begin{align*}
 \sum_i(2t_i-(n-1))^2
 &\le\sum_i\left|w_i^2-\frac{s_2}{n}\right|^2
 =u_2-\frac{|s_2|^2}{n},\\
 \sum_it_i(3t_i-(2n-3))^2
 &\le\sum_i\left|w_i^3-\frac{s_3}{n}\right|^2
 =u_3-\frac{|s_3|^2}{n}.
\end{align*}
The first left-hand side is $4u_2-n(n-1)^2$ and the second is
$9u_3-6(2n-3)u_2+(2n-3)^2N$.  This proves both cuts.
Rescaling gives the stated fourth-moment bound.  If $x=|z_i|^2$,
Cauchy's inequality for the other $n-1$ radii gives
\[
 x^2+\frac{(n-x)^2}{n-1}\le\frac{4n}{3}.
\]
Completing the square yields $(x-1)^2\le(n-1)/3$, proving the
radius bound.
\end{proof}

The same argument supplies higher-order conditions, including the
off-diagonal conditions that scalar radial cuts omit.  Let
$v_k=(w_i^k)_i$ and $\widetilde v_k=v_k-(s_k/n)\mathbf1$.
For $k\ge2$, define the vector $A_k$ by
\begin{equation}\label{eq:power-stability-action}
 (A_k)_i=
 k w_i^{k-1}\overline w_i
 -\sum_{\ell=0}^{k-2}(k-1-\ell)s_\ell w_i^{k-2-\ell}
 +\frac{k(k-1)}2w_i^{k-2},
 \qquad A_1=\overline w.
\end{equation}

\begin{proposition}\label{prop:finite-power-stability}
At every force-$1$ stationary configuration with distinct points,
$Bv_k=A_k$.  For each integer $K\ge1$, form the Hermitian matrix
\[
 \mathcal G^{(K)}_{jk}
 =\langle\widetilde v_j,\widetilde v_k\rangle
  -\langle A_j,A_k\rangle,
 \qquad 1\le j,k\le K,
 \qquad \langle a,b\rangle=\sum_i\overline a_i b_i.
\]
All these matrices are positive semidefinite at a local maximum.
Moreover, $\mathcal G^{(n-1)}\succeq0$ is equivalent to
$\|B\|\le1$.  Thus this one finite matrix expresses the entire
second-variation condition in polynomial moments.
\end{proposition}

\begin{proof}
For $k\ge2$ the elementary identity
\[
 \frac{x^k-y^k}{(x-y)^2}
 =\frac{kx^{k-1}}{x-y}
  -\sum_{\ell=0}^{k-2}(k-1-\ell)x^{k-2-\ell}y^\ell
\]
holds for $x\ne y$.  Sum it over $j\ne i$ with
$x=w_i$, $y=w_j$.  Stationarity gives the first term of
\eqref{eq:power-stability-action}; replacing
$\sum_{j\ne i}w_j^\ell$ by $s_\ell-w_i^\ell$ gives the last term
because $\sum_{\ell=0}^{k-2}(k-1-\ell)=k(k-1)/2$.
The formula for $k=1$ is precisely the force equation.

For any $\alpha\in\mathbb C^K$, the quadratic form of the displayed
matrix equals
\[
 \alpha^*\mathcal G^{(K)}\alpha
 =\left\|\sum_{k=1}^K\alpha_k\widetilde v_k\right\|^2
  -\left\|B\sum_{k=1}^K\alpha_k\widetilde v_k\right\|^2.
\]
This proves necessity.  For the converse, distinctness implies that
$\widetilde v_1,\ldots,\widetilde v_{n-1}$ form a basis of
$\mathbf1^\perp$: a linear dependence would give a polynomial of
degree at most $n-1$ taking the same constant value at $n$ distinct
points, so all its positive-degree coefficients vanish.
Consequently $\mathcal G^{(n-1)}\succeq0$ states that $B$ is a
contraction on $\mathbf1^\perp$.  Since $B\mathbf1=0$, it is then
a contraction on all of $\mathbb C^n$.
\end{proof}

To express the matrix algebraically, use mixed moments
$m_{ab}=\sum_i\overline w_i^{\,a}w_i^b$.
Then
\[
 \langle\widetilde v_j,\widetilde v_k\rangle
 =m_{jk}-\frac{\overline s_js_k}{n},
\]
and \eqref{eq:power-stability-action} expresses every
$\langle A_j,A_k\rangle$ as a polynomial in these moments.
For $\mathcal G^{(n-1)}$ only moments with
$0\le a,b\le n-1$ are needed.
The diagonal cases $k=2,3$ recover
\eqref{eq:general-fourth-cut}--\eqref{eq:general-sixth-cut}.
These constraints can therefore strengthen the coefficient and
moment optimization formulations without introducing inverse
interpoint distances.  They are necessary for a local maximum;
no sufficiency for local or global maximality is claimed.
At a force-$1$ stationary configuration the first row and column
of $\mathcal G^{(n-1)}$ vanish automatically: centering and symmetry
of $B$ give
\[
 \langle w,\widetilde v_k\rangle
 =\overline w^{\mathsf T}v_k
 =(Bw)^{\mathsf T}v_k
 =w^{\mathsf T}Bv_k
 =\langle A_1,A_k\rangle.
\]
For $n\ge3$ it is therefore enough to impose positive
semidefiniteness of the $(n-2)\times(n-2)$ block indexed by
$k=2,\ldots,n-1$.

\begin{corollary}\label{cor:radius-capped-grid}
For an integer $q\ge n$, let $R_n^{\mathrm{cap}}(q)$ be the smallest
integer $R\ge q$ satisfying
\[
 3(R^2-q^2)^2\ge q^4(n-1).
\]
Equivalently,
\[
 R_n^{\mathrm{cap}}(q)
 =\left\lceil q\sqrt{1+\sqrt{(n-1)/3}}\right\rceil.
\]
Define $B_n^{\mathrm{cap}}(q)$ by the same centering and normalization
as $B_n(q)$, but enumerate only the Gaussian integer coordinates
$a_i$ with
$|\operatorname{Re}a_i|,|\operatorname{Im}a_i|
 \le R_n^{\mathrm{cap}}(q)$.
Discard the zero-moment tuples as usual.  Then
\[
 B_n^{\mathrm{cap}}(q)\le M_n
 \le B_n^{\mathrm{cap}}(q)e^{E_n(q)},
\]
with exactly the exponent in Theorem~\ref{thm:grid-sharper}.
For $0<\varepsilon\le1$ and
$q=\lceil n/\sqrt\varepsilon\rceil$, the coordinate height is
$O(n^{5/4}/\sqrt\varepsilon)$ and the logarithmic error is less
than $\varepsilon$.
\end{corollary}

\begin{proof}
The radial bound in Proposition~\ref{prop:general-radial-stability}
applies to a global maximizer $z^*$.  Each real and imaginary
coordinate of $qz_i^*$ therefore lies between
$-q\sqrt{1+\sqrt{(n-1)/3}}$ and its positive counterpart.
Rounding to a nearest integer places every coordinate in the
displayed box.  The same rounded tuple used in the proof of
Theorem~\ref{thm:grid-sharper} is thus present in this smaller
complete enumeration, and every error estimate in that proof is
unchanged.  Conversely, each enumerated tuple of positive centered
moment gives an admissible normalized configuration, so its value
is at most $M_n$.  The integer inequality defining $R$ is equivalent
to the radical bound because $R\ge q$ permits both squarings.
Finally the radial bound is $O(n^{1/4})$, which gives the claimed
height after multiplication by $q$.
\end{proof}

\subsection{Higher radial moments and shell populations at local maxima}

The finite power stability matrix yields additional scalar inequalities
whose mixed terms can be eliminated exactly.  The following results
apply to every local maximum, so they retain information about local
searches as well as the global optimization problem.  Use the force-$1$
normalization and notation of Proposition~\ref{prop:general-radial-stability}.

\begin{proposition}\label{prop:eighth-radial-cut}
At every local maximum,
\begin{equation}\label{eq:eighth-force-cut}
 15u_4-24(n-2)u_3+9(n-2)^2u_2-n|s_2|^2
       +\frac{|s_4|^2}{n}\le0.
\end{equation}
\end{proposition}

\begin{proof}
The power action formula gives
\[
 Bv_4=\bigl(w_i^2(4t_i-3(n-2))-s_2\bigr)_i.
\]
Pairwise summation of the force equations, after multiplication by
$w_i^3$, gives the exact identity
\[
 \sum_i t_iw_i^2
 =\sum_{i<j}(w_i^2+w_iw_j+w_j^2)
 =\left(n-\frac32\right)s_2.
\]
Put $d=3(n-2)$ and $X_i=w_i^2(4t_i-d)$.  This identity implies
$\sum_iX_i=ns_2$.  Therefore
\[
 \|Bv_4\|^2
 =\sum_i|X_i-s_2|^2
 =16u_4-8du_3+d^2u_2-n|s_2|^2.
\]
Since $B\mathbf1=0$ and $\|B\|\le1$, this is at most
$\|v_4-(s_4/n)\mathbf1\|^2=u_4-|s_4|^2/n$.
Rearranging proves the proposition.
\end{proof}

\begin{theorem}\label{thm:combined-higher-radial}
In the original normalization $\sum_i|z_i|^2=n$, write
$T_p=\sum_i|z_i|^{2p}$ and $S_k=\sum_i z_i^k$.
Every local maximum satisfies, for $n\ge2$,
\begin{equation}\label{eq:combined-sixth-radial}
 T_3+\frac{2n-3}{2n(n-1)}|S_2|^2+\frac{|S_3|^2}{8n}
 \le L_{n,3}:=\frac{n(2n-3)(2n-1)}{2(n-1)^2}
 =2n-\frac{n}{2(n-1)^2}<2n.
\end{equation}
For $n\ge3$ it also satisfies
\begin{equation}\label{eq:combined-eighth-radial}
 \begin{split}
 T_4+\frac{4(8n^2-30n+24)}{15n(n-1)^2}|S_2|^2
       +\frac{2(n-2)}{5n(n-1)}|S_3|^2
       +\frac{|S_4|^2}{15n}
 &\le L_{n,4}:=
       \frac{8n(n-2)(2n^2-2n-1)}{5(n-1)^3}\\
 &=\frac{16n}{5}-\frac{24n}{5(n-1)^2}
                    +\frac{8n}{5(n-1)^3}<\frac{16n}{5}.
 \end{split}
\end{equation}
All the additional holomorphic-moment terms have nonnegative
coefficients in the stated ranges.
\end{theorem}

\begin{proof}
Let $c=2n-3$, $d=3(n-2)$ and
$D=6dc-d^2=9(n-2)(3n-4)$.  The existing second- and third-power
cuts, with their left-hand sides denoted by $C_2,C_3$, are
\[
 C_2=3u_2+\frac{|s_2|^2}{n}-n(n-1)^2\le0,
 \qquad
 C_3=8u_3-6cu_2+c^2N+\frac{|s_3|^2}{n}\le0.
\]
The nonnegative combination $C_3+2cC_2$ cancels $u_2$ and gives
\[
 8u_3+\frac{2c}{n}|s_2|^2+\frac{|s_3|^2}{n}
 \le\frac{n(n-1)c(2n-1)}2.
\]
Rescaling $w_i=\sqrt{(n-1)/2}\,z_i$ proves
\eqref{eq:combined-sixth-radial}.

For $n\ge3$, both $d$ and $D$ are positive.  Let $C_4$ be the
left-hand side of \eqref{eq:eighth-force-cut}.  The nonnegative
combination $C_4+dC_3+(D/3)C_2$ cancels $u_3$ and $u_2$, yielding
\[
 15u_4+\frac{8n^2-30n+24}{n}|s_2|^2
       +\frac{3(n-2)}{n}|s_3|^2+\frac{|s_4|^2}{n}
 \le\frac32n(n-1)(n-2)(2n^2-2n-1).
\]
The remaining coefficient $8n^2-30n+24$ is positive at $n=3$
and strictly increases for $n\ge3$.  Rescaling gives
\eqref{eq:combined-eighth-radial}.  Expansion in $n-1$ gives the
displayed simplified bounds and their strict inequalities.
\end{proof}

\begin{corollary}\label{cor:local-radius-and-shells}
Every local maximum satisfies
\[
 \max_i|z_i|\le L_{n,3}^{1/6}\quad(n\ge2),\qquad
 \max_i|z_i|\le L_{n,4}^{1/8}\quad(n\ge3).
\]
The latter bound is $O(n^{1/8})$.  A sharper individual bound is
$|z_i|^2\le\xi_{n,p}$, for $p=3$, or for $p=4$ when $n\ge3$,
where $\xi_{n,p}\in(1,n)$ is the unique solution of
\[
 x^p+\frac{(n-x)^p}{(n-1)^{p-1}}=L_{n,p}.
\]
More generally, if $k$ of the radii satisfy $|z_i|^2\ge a>1$,
then $ka\le n$ and, for $k<n$,
\begin{equation}\label{eq:outer-shell-count}
 k a^p+\frac{(n-ka)^p}{(n-k)^{p-1}}\le L_{n,p}.
\end{equation}
For $0\le a<1$ the same inequality holds if $k<n$ radii satisfy
$|z_i|^2\le a$.  In particular,
for every $a>0$,
$\#\{i:|z_i|^2\ge a\}\le\lfloor L_{n,p}/a^p\rfloor$.
These are necessary shell-population bounds, without any assumption
that the points lie on exact rings.
\end{corollary}

\begin{proof}
Drop the nonnegative holomorphic-moment terms to obtain
$T_p\le L_{n,p}$.  Jensen's inequality on the other $n-1$ radii
gives the individual bound.  The function
$x^p+(n-x)^p/(n-1)^{p-1}$ is strictly increasing for $1<x<n$;
its values at $1$ and $n$ bracket $L_{n,p}$ in the stated ranges.
This proves uniqueness and the sharper bound.

For a group of $k$ points let $S$ be the sum of their squared
radii.  Jensen on this group and its complement gives
\[
 T_p\ge\frac{S^p}{k^{p-1}}
          +\frac{(n-S)^p}{(n-k)^{p-1}}.
\]
This expression increases for $S>k$ and decreases for $S<k$.
For an outer group $S\ge ka>k$; for an inner group
$S\le ka<k$.  Substituting $S=ka$ in either case proves
\eqref{eq:outer-shell-count}.  The final counting bound follows
by retaining only the group's contribution to $T_p$.
\end{proof}

For example, at $n=6$ the new constants are
$L_{6,3}=297/25$ and $L_{6,4}=11328/625$.
Their application to local maxima complements stronger statements
that can use global optimality.  In particular, these local
necessary conditions should not be confused with sufficient
conditions for a global maximum.

\paragraph{Two elementary structural exclusions.}
There is no collinear local maximum when $n\ge3$.
Indeed, at the force-$1$ scale a real-axis configuration has
$s_2=N$ and
$u_2=(n-3/2)N=n(n-1)(2n-3)/4$ by the same pairwise sum as above.
The necessary bound $u_2\le n(n-1)^2/3$ would imply $2n\le5$.
Also, a stationary configuration whose radii are all equal must
be a regular polygon.  Its radius squared is $(n-1)/2$, so its
monic root polynomial satisfies
$tP''(t)-(n-1)P'(t)=0$ at all $n$ distinct roots.  This polynomial
has degree at most $n-2$ and is therefore identically zero.
Comparing coefficients gives $P(t)=t^n+c_0$.

The same scalar power tests also give concise proofs of the familiar
large-polygon instability thresholds.  If
$s_1=\cdots=s_{k-2}=0$, set
$d_k=(k-1)(n-k/2)$.  Then
\[
 Bv_k=(w_i^{k-2}(kt_i-d_k))_i,
\]
so contraction implies
\[
 (k^2-1)u_k-2kd_ku_{k-1}+d_k^2u_{k-2}
       +\frac{|s_k|^2}{n}\le0.
\]
For the regular $n$-gon with $n\ge8$, take $k=4$;
$s_2=s_4=0$ and $t_i=(n-1)/2$.  The left-hand side of
\eqref{eq:eighth-force-cut} becomes
\[
 \frac{3n(n-1)^2(n-3)(n-7)}{16}>0,
\]
excluding a local maximum.  For the origin plus a regular
$(n-1)$-gon with $n\ge11$, take $k=5$;
$s_2=s_3=s_5=0$ and the nonzero squared radii are $n/2$.
The last displayed cut has value
\[
 (n-1)\left(\frac n2\right)^3
       \left\{\left(10-\frac{3n}{2}\right)^2-\frac{n^2}{4}\right\}
 =\frac{n^3(n-1)(n-5)(n-10)}4>0.
\]
These elementary tests recover instability in the stated ranges;
they do not decide the degenerate boundary cases.

\subsection{Six-point radial constraints from the stability operator}

The following necessary conditions apply to the unrestricted global
maximum and do not assume reflection or polygon symmetry.  We use the
force-$1$ normalization
\[
 \sum_{j\ne i}\frac1{w_i-w_j}=\overline w_i,
 \qquad\sum_iw_i=0,\qquad\sum_i|w_i|^2=15.
\]
Write $t_i=|w_i|^2$, $u_k=\sum_i t_i^k$, and
$s_k=\sum_iw_i^k$.

\begin{proposition}\label{prop:six-radial-stability}
Every six-point global maximizer satisfies
\begin{align}
 3u_2+\frac{|s_2|^2}{6}&\le150,\label{eq:six-fourth-cut}\\
 8u_3-54u_2+1215+\frac{|s_3|^2}{6}&\le0.
 \label{eq:six-sixth-cut}
\end{align}
In particular $u_2\le50$.  In the original normalization
$z_i=\sqrt{2/5}\,w_i$, these imply
\[
 \sum_i|z_i|^4\le8,\qquad
 \max_i|z_i|^2\le1+\frac{\sqrt{15}}3.
\]
The same conclusions hold for a local maximum of the discriminant
on the moment sphere.
\end{proposition}

\begin{proof}
Let $B$ be the symmetric inverse-square Laplacian in
Lemma~\ref{lem:global-separation}, now at the force-$1$ scale.
A local constrained maximum is also a local maximum of
$\log\Delta_6(w)-\sum_i|w_i|^2$: for each nearby shape, its best
scale has moment $15$.  Consequently the same second-variation
argument gives $\|B\|\le1$ and $B\mathbf1=0$.
Polynomial divided differences and stationarity give
\[
 B(w_i^2)_i=(2t_i-5)_i,
 \qquad
 B(w_i^3)_i=(3w_it_i-9w_i)_i.
\]
For the first identity, use
$(w_i^2-w_j^2)/(w_i-w_j)^2=2w_i/(w_i-w_j)-1$.
For the second, use
\[
 \frac{w_i^3-w_j^3}{(w_i-w_j)^2}
 =\frac{3w_i^2}{w_i-w_j}-2w_i-w_j
\]
and $\sum_jw_j=0$.
Since $B$ annihilates constants,
\[
 \sum_i(2t_i-5)^2
 \le\sum_i\left|w_i^2-\frac{s_2}{6}\right|^2
 =u_2-\frac{|s_2|^2}{6},
\]
where the left-hand side is $4u_2-150$.
Similarly,
\[
 \sum_it_i(3t_i-9)^2
 \le\sum_i\left|w_i^3-\frac{s_3}{6}\right|^2
 =u_3-\frac{|s_3|^2}{6}.
\]
Expanding proves both inequalities.
Rescaling gives $\sum_i|z_i|^4\le(2/5)^2\cdot50=8$.
For $x=|z_i|^2$, Cauchy's inequality on the other five radii gives
$x^2+(6-x)^2/5\le8$, whose nonnegative solutions satisfy
$x\le1+\sqrt{15}/3$.
\end{proof}

\begin{proposition}\label{prop:six-excluded-subclasses}
No six-point local maximum belongs to any of the following three
subclasses: a collinear configuration, two concentric equilateral
triangles on the same three rays with different radii, or three
centrally symmetric pairs with four points on one line and two on
its perpendicular line.  In the last subclass the two lines pass
through the center of rotation.
\end{proposition}

\begin{proof}
All local maxima satisfy the force equations and the fourth-moment
bound $u_2\le50$.  Each listed subclass violates that bound.

For a collinear configuration, rotate its centered line to the real
axis.  Multiplying the force equation by $w_i^3$ and summing gives
\[
 u_2=\sum_{i<j}\frac{w_i^3-w_j^3}{w_i-w_j}
 =\sum_{i<j}(w_i^2+w_iw_j+w_j^2)
 =\frac92\sum_iw_i^2=\frac{135}{2}>50.
\]

For two aligned triangles, let their radii be $r<s$ and put
$a=r^2$, $b=s^2$.  Their root polynomial is
$(t^3-r^3)(t^3-s^3)$.  At a root on each ring the force equations give
\[
 a=1+\frac{3r^3}{r^3-s^3}<1,
 \qquad
 b=1+\frac{3s^3}{s^3-r^3}>4,
 \qquad a+b=5.
\]
Therefore $u_2=3(a^2+b^2)>3(1^2+4^2)=51>50$.

For the perpendicular-pair configuration rotate so its points are
\[
 \{\pm\sqrt a,\ \pm\sqrt b,\ \pm i\sqrt c\},
 \qquad a,b,c>0,\quad a\ne b.
\]
The moment gives $a+b+c=15/2$.  The force equations are
\[
 a=\frac12+\frac{2a}{a-b}+\frac{2a}{a+c},\qquad
 b=\frac12+\frac{2b}{b-a}+\frac{2b}{b+c},\qquad
 c=\frac12+\frac{2c}{c+a}+\frac{2c}{c+b}.
\]
Put $s=a+b$, $p=ab$, and $D=(a+c)(b+c)=p+cs+c^2>0$.
Adding and subtracting the first two equations yields
\[
 s=3+\frac{2(2p+cs)}D,
 \qquad (s^2-4p)(D-2c)=2sD.
\]
Together with $c=15/2-s$, these give
\[
 p=\frac{c(11s-45)}{2(7-s)},\qquad
 (a-b)^2=\frac{s(15-s)}4,\qquad
 p=\frac{5s(s-3)}{16}.
\]
The exceptional denominator $s=7$ is impossible in the first
summed equation.  Positivity of $p,c$ implies $45/11<s<7$.
Equating the two expressions for $p$ gives
\[
 f(s):=5s^3-138s^2+1125s-2700=0.
\]
On $[50/11,7]$, $f''(s)=30s-276<0$, while
\[
 f(50/11)=\frac{42550}{1331}>0,\qquad f(7)=128>0.
\]
Concavity therefore gives $f>0$ throughout that interval.
Every allowed root consequently has $s<50/11<5$.
Finally,
\[
 u_2-50
 =2(a^2+b^2+c^2)-50
 =\frac{11}{4}\left(s-\frac{50}{11}\right)(s-5)>0.
\]
This excludes the last subclass as well.
\end{proof}

\begin{proposition}\label{prop:six-large-radius-exclusion}
For any centered six-point configuration of moment $6$, if some
$|z_i|^2\ge(9-\sqrt{21})/2$, then
$\Delta_6\le C_6=6^{15}/5^{10}$.
Consequently a hypothetical global maximizer with value strictly above
$C_6$ has
\[
 0<|z_i|^2<\frac{9-\sqrt{21}}2\qquad\text{for every }i.
\]
\end{proposition}

\begin{proof}
Fix $i$ and write $x=|z_i|^2$.  Centering the other five points
gives quadratic moment $S=6-6x/5$.  Wang's five-point theorem gives
\[
 \Delta_5((z_j)_{j\ne i})\le5^5(S/5)^{10}.
\]
Also $\sum_{j\ne i}|z_i-z_j|^2=6(x+1)$, so arithmetic--geometric
mean gives
\[
 \Delta_6\le C_6\{(1-x/5)^2(1+x)\}^5.
\]
The centered moment constraint implies $0\le x\le5$, and
\[
 (1-x/5)^2(1+x)-1=\frac{x}{25}(x^2-9x+15).
\]
This expression is nonpositive when
$(9-\sqrt{21})/2\le x\le5$.
If $x=0$, the same bound gives $\Delta_6\le C_6$ directly.
\end{proof}

These constraints exclude specified families and a radial region.
Their role is to strengthen stationary filtering without prescribing
symmetry; they do not by themselves supply a complete classification.
The exhaustive comparison in Section~\ref{sec:six-global-exact}
establishes the six-point global value. The moment constraints and
finite stability matrix remain available for larger point counts.

\section{An exact polynomial test for stationary configurations}
\label{sec:trace-stationary}
Encoding an unordered configuration by its root polynomial removes
label permutations. A separate positivity condition is needed to
ensure that the force values are paired with the correct conjugate
roots, rather than merely permuted conjugate roots.

\begin{theorem}\label{thm:trace-stationary}
Let $P(t)=t^n+c_2t^{n-2}+\cdots+c_n$ be squarefree, and put
$\mathcal A=\C[t]/(P)$. Let $R$ be the unique polynomial of degree
less than $n$ satisfying
\begin{equation}\label{eq:force-modulo}
 P''-2RP'\equiv0\pmod P.
\end{equation}
For $0\le j,k<n$, define
\begin{equation}\label{eq:trace-matrix}
 H_{jk}=\operatorname{Tr}_{\mathcal A}(t^jR^k).
\end{equation}
The roots $w_i$ of $P$ satisfy
$\sum_{\ell\ne i}(w_i-w_\ell)^{-1}=\overline w_i$ if and only if
\begin{equation}\label{eq:trace-conditions}
 \overline P(R)\equiv0\pmod P,\qquad H=H^*,\qquad H\succ0.
\end{equation}
Here $\overline P$ means coefficientwise conjugation.
In this case $\sum_i|w_i|^2=N$ and the normalized product is
\begin{equation}\label{eq:disc-normalization}
 \Delta_n\!\left(\sqrt{\frac2{n-1}}\,w\right)
   =\left(\frac2{n-1}\right)^N|\operatorname{disc}P|.
\end{equation}
\end{theorem}
\begin{proof}
Squarefreeness makes $P'$ invertible modulo $P$, proving existence
and uniqueness of $R$. At a root,
\[
 R(w_i)=\frac{P''(w_i)}{2P'(w_i)}
       =\sum_{\ell\ne i}\frac1{w_i-w_\ell}.
\]
If this equals $\overline w_i$, both congruence and positivity hold:
the trace is the sum over the roots, so $H=VV^*$ with
$V_{ji}=w_i^j$, an invertible Vandermonde matrix.

Conversely, the congruence says each $R(w_i)$ is a conjugate root.
Positive definiteness makes $H$ nonsingular. Since
$H_{jk}=\sum_i w_i^jR(w_i)^k$, this forces the $R(w_i)$ to be distinct.
Thus $R(w_i)=\overline w_{\sigma(i)}$ for a permutation $\sigma$, and
$H=V\Pi_\sigma V^*$. Consequently $\Pi_\sigma$ is Hermitian positive
definite. A Hermitian permutation is an involution; every
nontrivial transposition has eigenvalue $-1$. It follows that
$\sigma$ is the identity, proving the actual conjugate pairing.
Finally $\sum_i w_iR(w_i)=N$ by pairwise summation.
The discriminant is $\prod_{i<j}(w_i-w_j)^2$, whose absolute value,
not squared absolute value, equals $\Delta_n(w)$.
\end{proof}

This is an exact equivalence, not merely a necessary system.
Let $T$ be the companion matrix representing multiplication by $t$.
The entries $H_{jk}=\operatorname{tr}(T^jR(T)^k)$ are polynomial in
coefficient variables. Taking $R$ as an auxiliary unknown makes both
remainders polynomial. Positive definiteness already enforces
squarefreeness: for distinct roots $\xi_1,\ldots,\xi_r$ with
multiplicities $m_i$, quotient-algebra traces give
$H_{jk}=\sum_{i=1}^r m_i\xi_i^jR(\xi_i)^k$, of rank at most $r$.
Thus $H\succ0$ forces $r=n$.
Hermitian symmetry and positive leading principal minors give real
polynomial conditions, without a separate discriminant denominator.

\begin{corollary}\label{cor:trace-global}
Allow $P$ and the auxiliary polynomial $R$ to vary subject to
\eqref{eq:force-modulo} and \eqref{eq:trace-conditions}. Then
\begin{equation}\label{eq:trace-global}
 M_n=\left(\frac2{n-1}\right)^N\max\det H.
\end{equation}
\end{corollary}
\begin{proof}
Every global maximizer gives a stationary configuration after scaling
by $\sqrt{(n-1)/2}$, and every permitted pair gives a genuine
stationary configuration of moment $N$. For these configurations,
$H=VV^*$, hence $\det H=|\det V|^2=\Delta_n(w)$.
\end{proof}
This is an exact real-polynomial reduction of the global problem;
maximizing its objective has not been carried out for six points.

% A rigorous compact coefficient formulation with a fixed positive margin.
\subsection{A compact coefficient formulation with explicit positivity margin}

The strict trace inequality can be replaced by an explicit closed inequality
without losing any global maximizer.  This gives a compact, polynomial
optimization problem for each $n$, with the same structure for every $n$.
It does not require a classification of the stationary configurations.

\begin{theorem}\label{thm:compact-trace}
Put
\[
 \mu_n=n^{-(2n-1)},\qquad
 b_{n,k}=\binom nk(n-1)^k\quad(2\le k\le n),\qquad
 C_n^{\rm aux}=(n-1)n^n.
\]
Let $\mathcal K_n$ be the set of coefficient pairs
\[
 P(t)=t^n+\sum_{k=2}^n c_kt^{n-k},\qquad
 R(t)=\sum_{\ell=0}^{n-1}r_\ell t^\ell
\]
subject to
\[
 P''-2RP'\equiv0\pmod P,\quad
 \overline P(R)\equiv0\pmod P,\quad
 H=H^*,\quad H-\mu_n I\succeq0,
\]
and the coefficient bounds
\[
 |c_k|^2\le b_{n,k}^2,\qquad
 |r_\ell|^2\le(C_n^{\rm aux})^2.
\]
Here $H_{jk}=\operatorname{Tr}_{\mathbb C[t]/(P)}(t^jR^k)$ as before.
Then $\mathcal K_n$ is a nonempty compact real semialgebraic set,
and
\[
 M_n=\left(\frac2{n-1}\right)^N
               \max_{(P,R)\in\mathcal K_n}\det H.
\]
All defining equations and inequalities become real polynomials with
rational coefficients after separating real and imaginary parts.
\end{theorem}

\begin{proof}
Consider any global maximizer and rescale it to roots $w_i$ with force
constant $1$.  Then $\sum_i|w_i|^2=N$ and
Lemma~\ref{lem:global-separation} gives
\[
 |w_i-w_j|^2\ge\frac n{n-1}>1.
\]
Since $\sum_iw_i=0$, Cauchy's inequality gives
\[
 |w_i|^2\le(n-1)\sum_{j\ne i}|w_j|^2,
 \qquad
 |w_i|^2\le\frac{n-1}{n}\sum_j|w_j|^2
       =\frac{(n-1)^2}{2}.
\]
In particular $|w_i|\le n-1$.
The elementary symmetric functions therefore satisfy
$|c_k|\le\binom nk(n-1)^k$.

Let
\[
 L_i(t)=\frac{\prod_{j\ne i}(t-w_j)}
                   {\prod_{j\ne i}(w_i-w_j)}
\]
be the Lagrange interpolation polynomials.  For a polynomial $f$, denote
the sum of the absolute values of its coefficients by $\|f\|_1$.
The root bounds and separation imply
\[
 \|L_i\|_1
 \le\frac{\prod_{j\ne i}(1+|w_j|)}
             {\prod_{j\ne i}|w_i-w_j|}
 \le n^{n-1}.
\]
As $R(w_i)=\overline w_i$, interpolation gives
\[
 R=\sum_i\overline w_iL_i,\qquad
 \|R\|_1\le(n-1)n^n,
\]
and hence the claimed bounds on all $r_\ell$.

With $V_{ji}=w_i^j$, the columns of $(V^{\mathsf T})^{-1}$ are the
coefficient vectors of the $L_i$.  It follows that
\[
 \|V^{-1}\|^2\le\|V^{-1}\|_F^2
 \le n^{2n-1}.
\]
Since $H=VV^*$, this proves $H\succeq\mu_n I$.
Thus every global maximizer supplies a member of $\mathcal K_n$.

Conversely $H\succeq\mu_n I$ implies $H\succ0$.
Theorem~\ref{thm:trace-stationary} then implies that every member of
$\mathcal K_n$ has distinct genuinely conjugate-paired roots, satisfies
the force equations, and has quadratic moment $N$.
Its normalized determinant is therefore at most $M_n$.
Together with the included global maximizer, this proves equality.

Remainders modulo a monic polynomial and quotient-algebra traces are
polynomials in the coefficient variables.  The Hermitian and positive
semidefinite conditions are closed real polynomial conditions, the latter
being equivalent to nonnegativity of all principal minors of
$H-\mu_n I$.  The displayed coefficient bounds make the set bounded.
Consequently $\mathcal K_n$ is closed and bounded, hence compact.
\end{proof}

For polynomial optimization one may append the redundant rational ball
constraint
\[
 \mathcal R_n-\sum_{k=2}^n|c_k|^2-\sum_{\ell=0}^{n-1}|r_\ell|^2\ge0,
 \qquad\mathcal R_n=\sum_{k=2}^nb_{n,k}^2+n(C_n^{\rm aux})^2.
\]
This supplies an explicit Archimedean quadratic module for a sum-of-squares
hierarchy after encoding the matrix inequality by polynomial principal
minors.  One may use the principal minors of
$(H+H^*)/2-\mu_n I$: these are real polynomials for every real coefficient
vector and equal the desired minors on the equations $H=H^*$.
Let $x$ collect the real and imaginary coefficient variables,
let $e_a(x)=0$ denote the equations, and let $g_b(x)\ge0$ include
the coefficient bounds, the principal minors, and the displayed ball.
For any rational $U>M_n$, Putinar's positivity theorem, in the
Archimedean setting used by Lasserre~\cite{Lasserre}, gives a finite
identity with real coefficients
\[
 U-\left(\frac2{n-1}\right)^N\operatorname{Re}\det H
 =\sigma_0+\sum_b\sigma_b g_b+\sum_a\tau_a e_a,
\]
where the $\sigma_b$ are sums of squares and the $\tau_a$ are arbitrary
polynomials.  Taking the degree truncations of these certificates gives
global upper bounds whose infimum is $M_n$, equivalently the standard
moment and sum-of-squares hierarchy converges to $M_n$.
This theorem establishes existence of real-coefficient certificates;
turning numerical semidefinite output into an exact certificate requires
an additional arithmetic verification.
Finite convergence at the exact optimum requires an additional theorem;
it is not a consequence of compactness alone.

\begin{proposition}[A genuine spurious six-root pairing]
Let $p$ be the unique root in $(11/25,9/20)$ of
$8p^3-60p^2+126p-45=0$, and put
\[
 s=2p^2-10p+\frac{15}{2},\qquad C=\frac{15}{2}-2p,
 \qquad P(t)=(t^4-st^2+p^2)(t^2-C).
\]
The polynomial
\[
 R(t)=\left(p^2-6p+\frac{41}{4}\right)t
   +\left(-\frac{32p^2}{45}+4p-\frac{26}{5}\right)t^3
   +\left(\frac{29}{75}-\frac{4p^2}{225}\right)t^5
\]
satisfies both force and conjugate-root congruences, and its trace
matrix is Hermitian, but is not positive definite.
\end{proposition}
\begin{proof}
The cubic changes sign on the given rational interval and its
 derivative is positive there. Exact bounds give
$17/5<s<347/100$ and $C>33/5$.
If $A<B$ are the roots of $a^2-sa+p^2$, then $0<A<B<C$.
Thus the six roots are $\pm\sqrt A,\pm\sqrt B,\pm\sqrt C$.
Reduction by the cubic verifies
\begin{align*}
 P''-2RP'&\equiv0\pmod P,&P(R)&\equiv0\pmod P,\\
 tR&\equiv p\pmod{t^4-st^2+p^2},&
 tR&\equiv C\pmod{t^2-C}.
\end{align*}
Since $AB=p^2$ and $p>0$, the force map swaps the positive roots
$\sqrt A,\sqrt B$ and their negative counterparts, while fixing
$\pm\sqrt C$. Its permutation is an involution, so $H$ is Hermitian.
Its principal minor with indices $0,2,4$ equals
\[
 -8(s^2-4p^2)(p^2-Cs+C^2)^2
       =-48600+19440p+2592p^2<0.
\]
Consequently $H$ is indefinite, with inertia $(4,2,0)$.
All displayed identities are independently verified in
$\Q[p]/(8p^3-60p^2+126p-45)$ by the ancillary checker.
\end{proof}

\begin{proposition}[Classical collinear benchmark]\label{prop:collinear}
Among all collinear configurations satisfying the original moment
constraint, the exact maximum is
\begin{equation}\label{eq:collinear}
 R_n^{\rm line}=\frac{\prod_{k=1}^n k^k}{(n-1)^N}.
\end{equation}
In particular, $R_6^{\rm line}=(2/729)C_6$, so every collinear
six-point configuration has product strictly below the central pentagon.
\end{proposition}
\begin{proof}
Center and rotate the line to the real axis. In each ordered chamber,
$\log\Delta_n(x)-\sum_i x_i^2$ is strictly concave, tends to $-\infty$
at collisions and infinity, and has a unique maximizer. Its root
polynomial $P_n$ satisfies
\[
 P_n''-2tP_n'+2nP_n=0,\qquad
 P_n'=nP_{n-1},\qquad
 P_n=tP_{n-1}-\frac{n-1}{2}P_{n-2}.
\]
These are the monic physicists' Hermite polynomials, with $P_0=1$,
$P_1=t$. Pairwise summation gives moment $N$ at this scale.
The recurrence for the absolute resultant is
\[
 |\operatorname{Res}(P_n,P_{n-1})|
 =\left(\frac{n-1}{2}\right)^{n-1}
     |\operatorname{Res}(P_{n-1},P_{n-2})|.
\]
Since $P_n'=nP_{n-1}$, it follows that
$|\operatorname{disc}P_n|=2^{-N}\prod_{k=1}^n k^k$.
Rescaling the moment from $N$ to $n$ gives \eqref{eq:collinear}.
\end{proof}

\section{A finite terminating exact-value algorithm}
\label{sec:exact-algorithm}
A second reduction gives low-degree critical equations directly.
For any configuration choose a farthest pair, apply a similarity
sending it to $w_1=-1,w_2=1$, and write
$w_i=x_i+iy_i$ for $3\le i\le n$.
All remaining points satisfy $|w_i-1|,|w_i+1|\le2$, hence
$|x_i|\le1$, $|y_i|\le\sqrt3$. With
\[
 \mu=\frac1n\sum_iw_i,\qquad Q=\sum_i|w_i-\mu|^2\ge2,
\]
the normalized product is the rational function
\begin{equation}\label{eq:farthest-chart}
 F(w)=\frac{n^N\prod_{i<j}|w_i-w_j|^2}{Q^N}.
\end{equation}
Every distinct real configuration in this chart, even outside the
farthest-pair lens, normalizes to a feasible configuration. Therefore
its unrestricted real maximum is exactly $M_n$, in $2n-4$ variables.
No reflection or polygon symmetry is assumed.

For complexification let $Z_i,W_i$ be independent variables,
$Z_1=W_1=-1$, $Z_2=W_2=1$, and put
\[
 \mu_Z=\frac1n\sum_iZ_i,\quad \mu_W=\frac1n\sum_iW_i,\quad
 Q=\sum_iZ_iW_i-n\mu_Z\mu_W,\quad
 V(Z)=\prod_{i<j}(Z_i-Z_j).
\]
Define, for $i\ge3$,
\begin{equation}\label{eq:critical-low-degree}
 G_i=Q\sum_{j\ne i}\prod_{k\ne i,j}(Z_i-Z_k)
       -N(W_i-\mu_W)\prod_{j\ne i}(Z_i-Z_j),
\end{equation}
and define $J_i$ by swapping $Z$ and $W$.
There are $2n-4$ equations, each of total degree at most $n$.
On the set $QV(Z)V(W)\ne0$, they are exactly the critical equations
of $F=n^NV(Z)V(W)/Q^N$.

\begin{theorem}\label{thm:exact-algorithm}
There is a nonzero polynomial $p_n\in\Q[T]$, computable from
\eqref{eq:critical-low-degree} by saturated elimination, such that
\begin{equation}\label{eq:degree-bound}
 p_n(M_n)=0,\qquad \deg p_n\le n^{2n-4}.
\end{equation}
Combining real-root isolation for $p_n$ with the rational enclosures
\eqref{eq:grid-rational} gives a finite terminating algorithm to
identify the exact value $M_n$ for every fixed $n$.
This is an algorithmic characterization; $p_n$ has not been computed
here for $n\ge4$.
\end{theorem}
\begin{proof}
The real maximum in the chart is a critical point, since it is also
a maximum over the entire real chart and has no collisions.
Write $h=QV(Z)V(W)$ and form the graph ideal
\[
 \mathfrak J=\bigl(G_i,J_i,\ TQ^N-n^NV(Z)V(W),\ sh-1\bigr)
 \subset\Q[Z_3,\ldots,Z_n,W_3,\ldots,W_n,T,s].
\]
The equation $sh-1=0$ enforces nonvanishing throughout elimination.
Equivalently, saturate the entire graph ideal by $h$, not just the
critical ideal before adjoining the value equation.
On each irreducible component meeting that nonvanishing set, all
partial derivatives of $F$ vanish. In characteristic zero, a rational
function with zero differential on an irreducible variety is constant.
Thus there are finitely many complex critical values, even if some
critical components have positive dimension. Elimination to $\Q[T]$
gives a nonzero polynomial vanishing on them and on $M_n$.

Repeated affine B\'ezout bounds give at most $n^{2n-4}$ irreducible
components: intersect successively with the $2n-4$ hypersurfaces of
degree at most $n$; the sum of the degrees of retained components
increases by at most a factor $n$ at each step. Saturation removes
components. There are no more distinct critical values than components.
Their set is invariant under the rational coefficient field, so its
squarefree product polynomial is rational and has the stated degree.

Now isolate the distinct real roots of $p_n$ exactly. For increasing
$q$ with $q\ge2(n-1)$ and $q^2>K_n$, completely compute $B_n(q)$.
The rational interval in \eqref{eq:grid-rational} contains $M_n$ and
has length at most $M_n(K_n/q^2)/(1-K_n/q^2)$, which tends to zero.
Eventually it contains exactly one distinct real root of $p_n$.
That root is $M_n$, irrespective of whether other roots arise from
nonphysical complex components. All enumerations and algebraic
operations at each finite stage terminate. Once the value is isolated,
real-algebraic sampling in its maximizing fiber supplies algebraic
coordinates attaining it \cite{Basu}.
\end{proof}

For $n=2$ there are no free variables and $p_2(T)=T-4$.
For $n=3$ write $Z_3=u,W_3=v$; then
\[
 Q=2+\frac23uv,\qquad
 F=\frac{108(1-u^2)(1-v^2)}{Q^3}.
\]
The two critical equations reduce to
\[
 6u-u^2v+3v=0,\qquad6v-v^2u+3u=0.
\]
Multiplying respectively by $v$ and $u$ and subtracting gives
$v^2=u^2$. The case $v=u$ gives $u=0,\pm3$ and value $27/2$;
the case $v=-u$ gives $u=0,\pm i\sqrt3$ and values $27/2,27$.
Consequently
\[
 p_3(T)=(2T-27)(T-27),\qquad M_3=27.
\]
The fully enumerated lower bound \eqref{eq:grid-three} exceeds $27/2$,
so the grid enclosure isolates the larger critical value.
It is attained by normalizing the triangle $(-1,1,i\sqrt3)$.
For six points the chart has eight real variables and the degree
bound is $6^8=1\,679\,616$; this bound is not a computed six-point
polynomial, nor a claim that such elimination is inexpensive.


\paragraph{Precise saturation in the critical-value construction.}
Write $h=QV(Z)V(W)$ and let $I$ be generated by the $G_i,J_i$.
The graph ideal to eliminate is
\[
 \mathfrak J=
 \bigl(I,\ TQ^N-n^NV(Z)V(W),\ sh-1\bigr)
 \subset\mathbb Q[Z_3,\ldots,Z_n,W_3,\ldots,W_n,T,s].
\]
Equivalently, saturate the \emph{whole graph ideal}
$\bigl(I,TQ^N-n^NV(Z)V(W)\bigr)$ by $h$ and then eliminate.
Saturating only $I$ before adjoining the value equation is insufficient:
its closure can contain boundary points with $Q=V(Z)=V(W)=0$, where
the value equation imposes no condition on $T$.  The auxiliary equation
$sh-1=0$ excludes these points throughout elimination.  The supplied
SageMath generator already uses this precise simultaneous construction.

\paragraph{An alternative exact selection of the physical maximum.}
After computing and isolating the real roots of the critical-value
polynomial, complete grid enumeration is optional.  For each isolated
algebraic root $t$, substitute
$Z_i=x_i+iy_i$ and $W_i=x_i-iy_i$, with
$(x_1,y_1)=(-1,0)$ and $(x_2,y_2)=(1,0)$, and test the real algebraic
feasibility of the critical equations, the value equation $F=t$,
and
\[
 \frac4{(n-1)^2}\le(x_i-x_j)^2+(y_i-y_j)^2\le4\qquad(i<j).
\]
This is a compact semialgebraic feasibility problem; its coordinates
satisfy $|x_i|\le1$, $|y_i|\le\sqrt3$, and $Q\ge2$.
Every feasible value is the value of an admissible normalized
configuration and is at most $M_n$.  A global maximizer, placed in
its farthest-pair chart, supplies a feasible value $M_n$: its original
diameter squared is at most $2n$, so Lemma~\ref{lem:global-separation}
gives the displayed lower bound on the chart distances after scaling
the diameter to $2$.
Consequently the largest feasible root is exactly $M_n$.
Real algebraic quantifier elimination or cylindrical algebraic
decomposition terminates for each of the finitely many roots and
can return algebraic coordinates attaining that root.
This alternative avoids replacing a searched grid subset by an
uncertified global enclosure; it does not imply low computational cost.

\section{An integer-coefficient moment framework for every \texorpdfstring{$n$}{n}}
\label{sec:moment-framework}
The previous Gaussian average is the first member of a convergent
family. Higher moments provide a uniform global certificate without
enumerating point configurations. The underlying ingredients are
classical: Gaussian monomial orthogonality, the Bombieri norm, and
the passage from $L^p$ to $L^\infty$. Vandermonde-power expansions
also occur in two-dimensional plasma computations
\cite{TF,Beauzamy}. The following formulation makes their error
and their connection to the exact-value algorithm explicit.

Write $V(z)=\prod_{i<j}(z_i-z_j)$, and for an integer $k\ge1$ expand
\[
 V(z)^k=\sum_{|\alpha|=kN}c_{\alpha}^{(n,k)}z^\alpha,
 \qquad c_{\alpha}^{(n,k)}\in\mathbb Z.
\]
Define the positive integer and the rational number
\begin{equation}\label{eq:coefficient-moment}
 C_{n,k}=\sum_{|\alpha|=kN}
       \bigl(c_{\alpha}^{(n,k)}\bigr)^2\prod_{i=1}^n\alpha_i!,
 \qquad
 A_{n,k}=\frac{n^{kN}(n-2)!\,C_{n,k}}{(kN+n-2)!},
 \qquad A_{n,0}=1.
\end{equation}
Coefficients with value zero contribute nothing. Computing the entire
sum is essential: a partial sum cannot be used for a global upper bound.

\begin{theorem}[Uniform coefficient enclosures]\label{thm:moment-framework}
For every $n\ge2$ and $k\ge1$, put
$D_{n,k}=\binom{kN+n-2}{n-2}$. Then
\begin{equation}\label{eq:moment-enclosure}
 A_{n,k}^{1/k}\le M_n\le
 \left(\frac{n^{kN}C_{n,k}}{(kN)!}\right)^{1/k}
 =(D_{n,k}A_{n,k})^{1/k}.
\end{equation}
The ratio of the two endpoints is exactly $D_{n,k}^{1/k}$;
for fixed $n$ its logarithm is
$O_n(\log(k+1)/k)$ and tends to zero. In particular,
\begin{equation}\label{eq:moment-exact-limit}
 \boxed{\displaystyle
 M_n=n^N\lim_{k\to\infty}
          \left(\frac{C_{n,k}}{(kN)!}\right)^{1/k}.}
\end{equation}
The entirely rational moment-ratio lower bounds
\begin{equation}\label{eq:moment-ratio}
 R_{n,k}:=\frac{A_{n,k+1}}{A_{n,k}}
 =\frac{n^NC_{n,k+1}}{C_{n,k}}
       \frac{(kN+n-2)!}{((k+1)N+n-2)!},\qquad k\ge0,
\end{equation}
are nondecreasing and converge to $M_n$. For $k\ge1$ they satisfy
$R_{n,k}\ge A_{n,k}^{1/k}$.
\end{theorem}

\begin{proof}
For independent standard complex Gaussians $Z_i$, orthogonality of
monomials gives $\mathbb E|V(Z)|^{2k}=C_{n,k}$.
Project $Z$ orthogonally onto
$\mathcal H=\{z:\sum_i z_i=0\}$; translation invariance preserves
$V$. This projection is a standard Gaussian in complex dimension
$d=n-1$. Its squared radius has the $\Gamma(d,1)$ law and is
independent of its uniform spherical direction $U$. Thus
\[
 \mathbb E\Delta_n(\sqrt n\,U)^k
 =\frac{n^{kN}(d-1)!\,C_{n,k}}{(kN+d-1)!}=A_{n,k}.
\]
The first inequality in \eqref{eq:moment-enclosure} is the fact
that an average does not exceed the maximum.

For completeness, any homogeneous holomorphic polynomial
$p(u)=\sum_{|\alpha|=s}b_\alpha u^\alpha$ in orthonormal coordinates
of $\mathcal H$ satisfies the elementary Cauchy--Schwarz inequality
\[
 |p(u)|^2\le
 \left(\sum_{|\alpha|=s}|b_\alpha|^2\frac{\alpha!}{s!}\right)
 \left(\sum_{|\alpha|=s}\frac{s!}{\alpha!}|u^\alpha|^2\right)
 =\|p\|_{B,s}^2\|u\|^{2s}.
\]
Gaussian integration shows that for the restriction $p=V^k$,
$\|p\|_{B,kN}^2=C_{n,k}/(kN)!$. Applying the inequality at
$\|u\|=\sqrt n$ proves the upper bound. The factorial identity
in \eqref{eq:moment-enclosure} follows directly from
\eqref{eq:coefficient-moment}.
Since $D_{n,k}$ is a degree-$(n-2)$ polynomial in $k$ for fixed $n$,
the endpoint ratio tends to one. Squeezing gives
\eqref{eq:moment-exact-limit}.

Let $X=\Delta_n(\sqrt n\,U)$. Then $0\le X\le M_n$ and
$A_{n,k}=\mathbb E X^k>0$. Cauchy--Schwarz gives
$A_{n,k+1}^2\le A_{n,k}A_{n,k+2}$, so $R_{n,k}$ increases;
$\mathbb E X^{k+1}\le M_n\mathbb E X^k$ gives $R_{n,k}\le M_n$.
Since $A_{n,k}=\prod_{j=0}^{k-1}R_{n,j}$ and all ratios are
nondecreasing, $A_{n,k}^{1/k}\le R_{n,k}$. The root bounds already
converge to $M_n$, proving convergence of the ratios.
\end{proof}

At $k=1$, $C_{n,1}=n!P_n$ and $A_{n,1}=L_n$, exactly recovering
the earlier Gaussian lower bound. The $k=1$ coefficient upper bound
can be substantially worse than the earlier Hadamard--Jensen upper
bound. They should be intersected, rather than replacing a better
finite-order bound. Convergence does not mean that a small moment
order is sufficient in six or more points.

% Fragment; classical partition compression with elementary proofs.
\subsection{Partition compression and why a Jack norm does not evaluate the Gaussian norm}
\label{subsec:uniform-partition-compression}
Write $\rho=(n-1,n-2,\ldots,0)$, $N=n(n-1)/2$, and expand
\[
 V_n(z)^k=\sum_{|\alpha|=kN}c_\alpha z^\alpha.
\]
For a decreasing partition $\lambda$ of $kN$ with at most $n$ parts,
let $a_\lambda=c_\lambda$, and let $m_j(\lambda)$ be the multiplicity
of the part $j$, including the zero parts used to reach length $n$.
Then the complete Gaussian coefficient sum can be written
\begin{equation}\label{eq:uniform-partition-norm}
 C_{n,k}=n!\sum_{\substack{\lambda\vdash kN,\ \lambda\preceq k\rho\\
                    \ell(\lambda)\le n}}
 \frac{a_\lambda^2}{\prod_{j\ge0}m_j(\lambda)!}
 \prod_{i=1}^n\lambda_i!.
\end{equation}
When $k$ is odd, only strict partitions occur and every nonzero term
has multiplicity denominator $1$.
\begin{proof}
Permuting variables changes $V_n^k$ by the $k$th power of the
permutation sign. Thus all coefficients in an exponent orbit have equal
squares. The orbit of $\lambda$ has $n!/\prod_jm_j!$ elements,
proving the formula. If $k$ is odd and two exponents coincide, the
transposition exchanging them fixes the monomial and negates its
coefficient, so that coefficient is zero.

For any subset of $r$ variables, the total exponent contributed by
the factors of $V_n^k$ is at most
\[
 k\left[r(n-r)+\binom r2\right]
 =\frac{kr(2n-r-1)}2.
\]
The $r$ largest exponents therefore have sum at most the first $r$
parts of $k\rho$. This proves the support restriction
$\lambda\preceq k\rho$. Cancellation may remove additional
partitions, so dominance is not an assertion that every allowed
coefficient is nonzero.
\end{proof}

This is the classical partition expansion used in T\'ellez--Forrester,
\href{https://arxiv.org/pdf/1204.6003}{\emph{Expanded Vandermonde Powers
and Sum Rules for the Two-Dimensional One-Component Plasma}},
equations~(2.4), (2.9), with the Gaussian radial integral
$G_j=j!$. Their negative-parameter Jack structure supplies coefficient
recurrences and can improve computation. It does not turn the Gaussian
norm into the usual Jack hook-product norm.

Here is a direct check of the mismatch. The standard algebraic Jack
scalar product makes distinct power-sum partitions orthogonal. In
contrast, for $n\ge2$ independent standard complex Gaussian variables,
\[
 \mathbb E\left[\left(\sum_iZ_i^2\right)
 \overline{\left(\sum_iZ_i\right)^2}\right]=2n\ne0.
\]
Thus the Gaussian scalar product used to define $C_{n,k}$ is different.
For this identity, expand the second factor: every term with unequal
holomorphic and antiholomorphic exponents vanishes, and each surviving
$Z_i^2\overline{Z_i}^2$ contributes $2$. Importing a Jack norm formula
without changing the scalar product is therefore invalid.

A second elementary obstruction applies to an ordinary differential
shift argument. Gaussian orthogonality implies
\[
 C_{n,k}=\left.V_n(\partial)^kV_n(z)^k\right|_{z=0},
\]
but one cannot recursively replace
$V_n(\partial)V_n(z)^k$ by a scalar multiple of $V_n(z)^{k-1}$.
For $n=3,k=3$, at $z=(0,0,1)$ one has
\[
 \left.(\partial_1-\partial_2)(\partial_1-\partial_3)
 (\partial_2-\partial_3)V_3(z)^3\right|_{(0,0,1)}=-12,
 \qquad V_3(0,0,1)^2=0.
\]
Only the third derivative of $(z_1-z_2)^3$ can survive there, which
gives the stated value. This does not exclude other exact identities;
it disproves this particular ordinary-differential shift.

One may also refute a specifically normalized Selberg extrapolation.
Suppose a product has the form
\[
 C_{n,k}=a_k^{N}\prod_{j=1}^n\frac{(jk)!}{k!}.
\]
The exact two-point value $C_{2,k}=2^kk!$ forces
$a_k=2^k(k!)^2/(2k)!$. At $n=3,k=2$ the product predicts $1280$,
whereas the exact unitary three-point expansion gives
\[
 C_{3,2}=6!+4(3!)^2+6!=1584.
\]
Consequently this entire one-factor-per-pair product family fails.
Neither this counterexample nor the scalar-product mismatch proves
that no finite closed formula of any kind can exist.

\subsection{From certified approximation to an exact algebraic value}
\begin{corollary}[Exact selection without a geometric grid]
\label{cor:moment-exact-selection}
Let $p_n$ be the nonzero critical-value polynomial constructed
using the saturated graph ideal in
Theorem~\ref{thm:exact-algorithm}. Complete integer-coefficient
moments and rational outward root enclosures give a finite
terminating algorithm to identify the algebraic number $M_n$.
No complete geometric grid enumeration is needed in this version.
\end{corollary}
\begin{proof}
Isolate all distinct real roots of $p_n$. For each increasing $k$,
compute the entire integer $C_{n,k}$ by finite polynomial
multiplication. The two endpoints of
\eqref{eq:moment-enclosure} are positive real $k$th roots of rational
numbers and can be enclosed outward by rational bisection, to
additional absolute error at most $1/k$. Their interval width
tends to zero, and the interval always contains $M_n$.
Eventually it contains exactly one of the finitely many distinct
real roots of $p_n$, which must be $M_n$. One can strengthen the
lower endpoint to the rational $R_{n,k}$ when $C_{n,k+1}$ has also
been computed. Every finite coefficient expansion, bisection,
and root-isolation operation terminates. This argument requires
neither isolated stationary shapes nor a guessed symmetry.
\end{proof}

The exact limit is a uniform formula in finite integer sums and a
limit, rather than a finite closed product for each $n$. The
corollary is a terminating algorithm, not an executed general-$n$
elimination. Its cost still grows rapidly; partition-based
Vandermonde expansions \cite{TF} offer a route to improve that cost.

\subsection{Complete arithmetic examples and the scope of the computation}
The standard-library program \texttt{verify\_moment\_hierarchy.py}
completely expands 25 cases: $n=2$, $k=1,\ldots,4$;
$n=3$, $k=1,\ldots,6$; $n=4$, $k=1,2,3,4,6,8,12,16$;
$n=5$, $k=1,\ldots,4$; and $n=6$, $k=1,2,3$.
All arithmetic in the sums and checks is exact integer or rational
arithmetic. Root endpoints are outward-certified rational numbers.
For example, the six-point third power has $177840$ nonzero
monomials and gives
\[
 6628.71382019\le M_6\le395175.36051878.
\]
This interval is weaker than the already available central-pentagon
lower bound and Hadamard--Jensen upper bound. Its purpose is to
verify the complete coefficient method, not to report a six-point
improvement from these low orders.

For three points, unitary centered coordinates give
$V_3=u^3+v^3$ up to phase, hence independently
\begin{equation}\label{eq:three-coefficient-moment}
 C_{3,k}=\sum_{j=0}^k\binom kj^2(3j)!\,(3k-3j)!.
\end{equation}
The direct three-variable expansion agrees with this independent
formula for every checked order $1\le k\le6$. Formula
\eqref{eq:three-coefficient-moment} permits much higher orders;
the computation includes $k=8,16,32,64,128,256$. Using the rational
moment-ratio lower bound gives, with outward-rounded endpoints,
\[
 \begin{array}{c|cc}
 k&\hbox{lower bound}&\hbox{upper bound}\\\hline
 16&25.41406404&28.22380137\\
 64&26.58521966&27.29554016\\
 256&26.89498543&27.07329696
 \end{array}
\]
These intervals contain the independently known value $M_3=27$.
At order $256$, the coefficient upper bound is within $0.272\%$
of that value. The examples validate the normalization and
convergence without using the known maximum to calculate an upper
bound.

\section{A convergent multiplication-operator spectral hierarchy}
\label{sec:spectral-hierarchy}
The norm of multiplication by a fixed homogeneous polynomial in
Bombieri's norm is a classical finite-dimensional eigenvalue problem;
see Beauzamy \cite{Beauzamy}. We apply this construction to the centered
Vandermonde, prove an explicit error factor, and distinguish convergence
from finite attainment of a sharp bound. No claim of priority is made
for the multiplication-operator framework.

Put $d=n-1$ and $\mathcal H=\{z\in\C^n:\sum_i z_i=0\}$. Choose an
isometry $B:\C^d\longrightarrow\mathcal H$, so $B^*B=I_d$, and set
\[
 f(u)=\prod_{i<j}\big((Bu)_i-(Bu)_j\big)
      =\sum_{|\gamma|=N}a_\gamma u^\gamma.
\]
Centering and homogeneity give
\begin{equation}\label{eq:spectral-normalization}
 S:=\max_{\|u\|=1}|f(u)|^2=\frac{M_n}{n^N}.
\end{equation}
For $s\ge0$, use the normalized monomial column
\[
 h_s(u)=\left(\sqrt{\binom{s}{\beta}}\,u^\beta\right)_{|\beta|=s},
 \qquad D_s=\binom{d+s-1}{s},\qquad
 \binom{s}{\beta}=\frac{s!}{\beta_1!\cdots\beta_d!}.
\]
Thus $\|h_s(u)\|^2=\|u\|^{2s}$. The degree-$s$ Bombieri norm is
\[
 \left\|\sum_{|\beta|=s}b_\beta u^\beta\right\|_{B,s}^2
 =\sum_{|\beta|=s}\frac{|b_\beta|^2}{\binom{s}{\beta}}.
\]
For an integer $r\ge0$, define the $D_{N+r}\times D_r$ matrix
\begin{equation}\label{eq:spectral-matrix}
 (A_r)_{\beta\alpha}
 =\overline{a_{\beta-\alpha}}
   \sqrt{\frac{\binom r\alpha}{\binom{N+r}\beta}},
 \qquad |\beta|=N+r,\quad |\alpha|=r,
\end{equation}
where $a_\gamma=0$ when a component of $\gamma$ is negative. Put
$K_r=A_rA_r^*$ and
$\Lambda_r=\lambda_{\max}(A_r^*A_r)=\lambda_{\max}(K_r)$.
The conjugation in \eqref{eq:spectral-matrix} makes the convention
$p(u)=c^*h_s(u)$ consistent with the usual Hermitian Gram matrix.

\begin{theorem}[Certified spectral bounds]\label{thm:spectral-hierarchy}
For every $r\ge0$,
\begin{equation}\label{eq:spectral-bounds}
 M_n\le n^N\Lambda_r
 \le\frac{D_{N+r}}{D_r}M_n,
 \qquad
 \frac{D_{N+r}}{D_r}
 =\prod_{j=1}^{n-2}\left(1+\frac{N}{r+j}\right).
\end{equation}
Put $T_r=(D_r/D_{N+r})\Lambda_r$. The sequence $\Lambda_r$ is
nonincreasing, $T_r$ is nondecreasing, and both tend to $M_n/n^N$.
Thus $[n^NT_r,n^N\Lambda_r]$ is a nested interval with the explicit
endpoint ratio $D_{N+r}/D_r$.
For $n\ge3$, $r\ge0$, and $\varepsilon>0$, the condition
\[
 r+1\ge\frac{N(n-2)}{\log(1+\varepsilon)}
\]
ensures $M_n\le n^N\Lambda_r\le(1+\varepsilon)M_n$.
Moreover, a proposed $U$ is certified by the exact matrix inequality
$UI_{D_{N+r}}-K_r\succeq0$, and then $M_n\le n^NU$.
Equivalently one may test the smaller matrix
$UI_{D_r}-A_r^*A_r\succeq0$ when $U\ge0$.
\end{theorem}
\begin{proof}
The column indexed by $\alpha$ is the conjugated normalized coefficient
vector of $f(u)\sqrt{\binom r\alpha}u^\alpha$. Consequently
\begin{equation}\label{eq:spectral-gram-identity}
 h_{N+r}(u)^*K_rh_{N+r}(u)
 =|f(u)|^2\sum_{|\alpha|=r}\binom r\alpha|u^\alpha|^2
 =|f(u)|^2\|u\|^{2r}.
\end{equation}
For a unit vector $u$, $h_{N+r}(u)$ is also a unit vector, so
\eqref{eq:spectral-gram-identity} proves $S\le\Lambda_r$.

Let $\sigma$ be probability surface measure on the unit sphere of
$\C^d$. Unitary invariance, or integration of monomials, gives
\begin{equation}\label{eq:bombieri-sphere}
 \int|p(u)|^2\,d\sigma(u)=\frac{\|p\|_{B,s}^2}{D_s}
 \quad\text{for every homogeneous degree-$s$ polynomial }p.
\end{equation}
Since the conjugate matrix $\overline{A_r}$ represents multiplication
by $f$ between the degree-$r$ and degree-$(N+r)$ Bombieri spaces,
\[
 \Lambda_r=\sup_{p\ne0}
   \frac{\|fp\|_{B,N+r}^2}{\|p\|_{B,r}^2}
 =\frac{D_{N+r}}{D_r}\sup_{p\ne0}
   \frac{\int|f|^2|p|^2\,d\sigma}{\int|p|^2\,d\sigma}
 \le\frac{D_{N+r}}{D_r}S.
\]
This proves \eqref{eq:spectral-bounds}; the product formula follows
by cancelling the factorials in $D_{N+r}/D_r$.

For any $U\ge\Lambda_r$, the Hermitian polynomial
\[
 U\|u\|^{2(N+r)}-|f(u)|^2\|u\|^{2r}
 =h_{N+r}(u)^*(UI-K_r)h_{N+r}(u)
\]
is a sum of squared absolute values of homogeneous holomorphic
polynomials. Multiplying this representation by
$\|u\|^2=\sum_j|u_j|^2$ preserves that property. The normalized
monomials form a basis, and the products
$u^\beta\overline{u^\gamma}$ are linearly independent as functions,
so the Hermitian coefficient matrix is unique. At level $r+1$ it is
$UI-K_{r+1}$, and it is positive semidefinite. Taking
$U=\Lambda_r$ gives $\Lambda_{r+1}\le\Lambda_r$.
For completeness, write
\[
 T_r=\sup_{p\ne0,\ \deg p=r}
 \frac{\int|f|^2|p|^2\,d\sigma}{\int|p|^2\,d\sigma}.
\]
For a maximizing $p$, the degree-$(r+1)$ polynomials $u_jp$ satisfy
\[
 \sum_j\int|u_jp|^2\,d\sigma=\int|p|^2\,d\sigma,
 \qquad
 \sum_j\int|f|^2|u_jp|^2\,d\sigma
       =\int|f|^2|p|^2\,d\sigma.
\]
The ratio on the right is a weighted average of the ratios on the
left, so at least one is at least $T_r$. Hence $T_{r+1}\ge T_r$.
The error factor tends to one, proving convergence. Finally,
\[
 \log\frac{D_{N+r}}{D_r}
 \le\sum_{j=1}^{n-2}\frac{N}{r+j}
 \le\frac{N(n-2)}{r+1},
\]
which proves the stated error criterion. The matrix certificate is
the same Gram identity. The two positive-semidefinite tests are
equivalent because the nonzero eigenvalues of $A_rA_r^*$ and
$A_r^*A_r$ coincide.
\end{proof}

At $r=0$ the matrix has a single column and
\begin{equation}\label{eq:spectral-level-zero}
 \Lambda_0=\|f\|_{B,N}^2=\frac{n!P_n}{N!}
             =\frac{C_{n,1}}{N!}.
\end{equation}
Indeed, a standard complex Gaussian $U$ in $\C^d$ satisfies
$\mathbb E|f(U)|^2=N!\|f\|_{B,N}^2$. If $Z$ is a standard Gaussian
in $\C^n$, orthogonal projection onto $\mathcal H$ preserves the
Vandermonde, hence this expectation equals
$\mathbb E|V(Z)|^2=n!P_n$.

\begin{proposition}[The first spectral level]\label{prop:spectral-level-one}
For every $n\ge2$,
\begin{equation}\label{eq:spectral-level-one}
 \Lambda_1=\frac{n+2}{2(N+1)}\Lambda_0.
\end{equation}
\end{proposition}
\begin{proof}
The Gaussian weighted covariance matrix
$\mathbb E(|f(U)|^2UU^*)$ commutes with the standard representation
of $S_n$ on $\mathcal H$. That representation is irreducible over
$\C$ (and one-dimensional when $n=2$), so the matrix is scalar.
Radial integration gives
\[
 \mathbb E(|f(U)|^2\|U\|^2)=(N+d)\mathbb E|f(U)|^2.
\]
The scalar is therefore $(N+d)n!P_n/d$. Gaussian integration and
the degree-$(N+1)$ Bombieri normalization show that $A_1^*A_1$ is
this scalar divided by $(N+1)!$, times the identity. Since
$(N+d)/d=(n+2)/2$, \eqref{eq:spectral-level-one} follows.
\end{proof}

One can choose $B$ with real algebraic entries, for example the
Helmert basis with column $k$ equal to
$(1,\ldots,1,-k,0,\ldots,0)/\sqrt{k(k+1)}$, $1\le k<n$.
Then every matrix entry in \eqref{eq:spectral-matrix} is algebraic.
Exact algebraic linear algebra or certified root isolation can
certify its largest eigenvalue. Numerical eigenvalues alone do not
give rigorous upper bounds. The matrix dimension $D_r$ still grows
rapidly with $n$ and $r$; the convergence theorem is not a claim of
uniformly inexpensive computation. The symmetric group acts on the
degree-$r$ space, and $A_r^*A_r$ commutes with this action, so
representation-theoretic block decomposition is available.

The spectral lower bounds also contain the scalar moment-ratio
lower bounds. If
$A_{n,k}=n^{kN}\int|f|^{2k}\,d\sigma$, using $p=f^k$ in the
variational expression for $T_{kN}$ gives
\[
 n^NT_{kN}\ge\frac{A_{n,k+1}}{A_{n,k}}.
\]
The spectral optimization allows all degree-$kN$ multipliers,
whereas this moment ratio uses one specified multiplier.

\subsection{A closed formula for the second spectral level}
The first two multiplier degrees can be evaluated for every point
count without constructing the full matrix.

\begin{proposition}[Second spectral level]\label{prop:spectral-level-two}
Let $C=n!P_n=C_{n,1}$. On the degree-two centered polynomial space,
the eigenvalues of $A_2^*A_2$ are
\begin{equation}\label{eq:spectral-level-two-eigenvalues}
 \frac{C}{(N+2)!}\left\{
 2(n+1),\quad \frac{n^2+6n+6}{3},\quad
 \frac{(3n+5)(n+2)}6\right\},
\end{equation}
with respective multiplicities $1$, $n-1$, and $n(n-3)/2$ when
$n\ge3$. For $n=3$ the third multiplicity is zero; for $n=2$
only the first eigenvalue occurs. Consequently
\[
 \Lambda_2=
 \begin{cases}
 2,&n=2,\\[2pt]
 11/10,&n=3,\\[2pt]
 \displaystyle\frac{n!P_n}{(N+2)!}
       \frac{(3n+5)(n+2)}6,&n\ge4.
 \end{cases}
\]
\end{proposition}
\begin{proof}
Let $Z$ be a standard complex Gaussian in $\C^n$. All moments in
this proof use the normalized weighted expectation
\[
 \mathcal E[g]=\frac{\mathbb E(|V(Z)|^2g(Z))}{C},
 \qquad \mathcal E[1]=1.
\]
Write $t=\sum_iZ_i$ and $w_i=Z_i-t/n$. Translation invariance of
$V$ implies that $t$ is independent of $w$ under this weighted
measure and has the centered complex Gaussian distribution of
variance $n$. Complex Gaussian integration by parts gives
\begin{equation}\label{eq:weighted-translation-ibp}
 \mathcal E[g\overline t]=
 \mathcal E\left[\sum_i\frac{\partial g}{\partial Z_i}\right]
\end{equation}
for every polynomial $g$; the derivative of the Vandermonde weight
vanishes after summing the translation derivatives.

For pairwise distinct indices as needed, the moments are
\begin{align*}
 m&=\mathcal E[|Z_1|^2]=\frac{n+1}{2},&
 \rho&=\mathcal E[Z_1\overline Z_2]=-\frac12,\\
 a&=\mathcal E[|Z_1|^4]=\frac{(n+1)(n+2)}3,&
 b&=\mathcal E[Z_1^2\overline Z_2^{\,2}]=-\frac{n-2}{3},\\
 c&=\mathcal E[|Z_1Z_2|^2]=\frac{(n+1)(3n+2)}{12},&
 d_0&=\mathcal E[Z_1^2\overline Z_1\overline Z_2]
      =-\frac{n+1}{3},\\
 e&=\mathcal E[Z_1^2\overline Z_2\overline Z_3]=\frac23,&
 g_0&=\mathcal E[Z_1Z_2\overline Z_1\overline Z_3]
       =-\frac{3n+1}{12},\\
 h&=\mathcal E[Z_1Z_2\overline Z_3\overline Z_4]=\frac12.
\end{align*}
Here $e,g_0$ require $n\ge3$, and $h$ requires $n\ge4$.
We give the derivation to fix the labeling and normalization.
Expanding $V$ as an alternating sum over permutations of exponents
$0,\ldots,n-1$ and integrating monomials gives
\[
 m=\frac1n\sum_{k=0}^{n-1}(k+1),\qquad
 a=\frac1n\sum_{k=0}^{n-1}(k+1)(k+2).
\]
For $b$, the only surviving permutation pairs exchange two
exponents differing by two; their signs are opposite. Thus
\[
 b=-\frac1{n(n-1)}\sum_{k=0}^{n-3}(k+1)(k+2).
\]
Applying \eqref{eq:weighted-translation-ibp} first to $Z_1$ and
then to $Z_1^2\overline Z_1$, $Z_1^2\overline Z_2$,
$Z_1Z_2\overline Z_1$, and $Z_1Z_2\overline Z_3$ gives, in order,
\begin{align*}
 m+(n-1)\rho&=1,\\
 a+(n-1)d_0&=2m,\\
 d_0+b+(n-2)e&=2\rho,\\
 d_0+c+(n-2)g_0&=m+\rho,\\
 2g_0+e+(n-3)h&=2\rho.
\end{align*}
Finally, homogeneity gives
$\mathcal E[(\sum_i|Z_i|^2)^2]=(N+n)(N+n+1)$, so
$na+n(n-1)c=(N+n)(N+n+1)$. These identities yield every moment
displayed above, with the last identity used only when $n\ge4$.

Consider the following mutually orthogonal spaces of centered
quadratic polynomials, using the degree-two Bombieri inner product:
\[
 \begin{split}
 \mathcal Q_0&=\operatorname{span}\{q\},\qquad q=\sum_iw_i^2,\\
 \mathcal Q_1&=\operatorname{span}\{w_i^2-w_j^2:i,j\},\\
 \mathcal Q_2&=\left\{p=\sum_{i<j}a_{ij}Z_iZ_j:
                 \sum_{j\ne i}a_{ij}=0\ \text{for every }i\right\}.
 \end{split}
\]
The notation in $\mathcal Q_2$ extends $a_{ij}=a_{ji}$, $a_{ii}=0$.
Its row-sum conditions make $p$ translation invariant.
For $n\ge3$, these spaces have dimensions $1$, $n-1$, and
$n(n-3)/2$, respectively. Indeed, $\mathcal Q_1$ is the standard
representation of $S_n$: the map
$(v_i)\mapsto\sum_iv_iw_i^2$ is injective on $\sum_iv_i=0$,
as its Bombieri norm squared is
$(1-2/n)\sum_i|v_i|^2>0$ there. The unoriented incidence matrix
of the complete graph has row rank $n$ for $n\ge3$, giving the
dimension of $\mathcal Q_2$. Their orthogonality follows by
expanding the quadratic coefficients and using the row sums;
their dimensions total $n(n-1)/2$, the whole centered quadratic
space. When $n=2$ that space consists of $\mathcal Q_0$ alone.

The same spaces are orthogonal for the weighted form
$(p_1,p_2)\mapsto\mathcal E[p_1\overline p_2]$.
The first and second spaces are orthogonal by permutation symmetry.
For $p\in\mathcal Q_2$, the moment table and its row sums give
$\mathcal E[Z_i^2\overline p]=0$ for every $i$. Independence of $t$
and $w$ gives the same result with $w_i^2$ in place of $Z_i^2$.
This proves the remaining orthogonality.

On $\mathcal Q_0$, the identities
\[
 \mathcal E\left[\left|\sum_iZ_i^2\right|^2\right]=2n^2,
 \quad \sum_iZ_i^2=q+t^2/n,
 \quad\mathcal E[|t^2/n|^2]=2
\]
give $\mathcal E[|q|^2]=2(n^2-1)$, whereas
$\|q\|_{B,2}^2=n-1$. Its weighted/Bombieri ratio is $2(n+1)$.
On $\mathcal Q_1$, invariance and irreducibility make that ratio
constant. Using $p=w_1^2-w_2^2$ gives
\[
 \begin{split}
 \mathcal E[|p|^2]
 &=2(a-b)-\frac4n\mathcal E[|w_1-w_2|^2]\\
 &=\frac{2n(n+4)}3-\frac{4(n+2)}n
   =\frac{2(n-2)(n^2+6n+6)}{3n},\\
 \|p\|_{B,2}^2&=\frac{2(n-2)}n.
 \end{split}
\]
The subtraction follows from
$Z_1^2-Z_2^2=p+(2t/n)(w_1-w_2)$ and the independence of $t$;
also $\mathcal E[|w_1-w_2|^2]=2(m-\rho)=n+2$.
This yields the second ratio $(n^2+6n+6)/3$.

For $p\in\mathcal Q_2$, put $D=\sum_{i<j}|a_{ij}|^2$.
The row sums imply that the sum of the ordered coefficient products
for distinct edges sharing a vertex is $-2D$; the sum for disjoint
edges is $D$, since $\sum_{i<j}a_{ij}=0$. The moment table therefore
gives
\[
 \mathcal E[|p|^2]=(c-2g_0+h)D,
 \qquad \|p\|_{B,2}^2=D/2,
\]
so the third ratio is
$2(c-2g_0+h)=(3n+5)(n+2)/6$.
Finally, Gaussian integration converts the output norm by
\[
 \|fp\|_{B,N+2}^2=\frac{C}{(N+2)!}\mathcal E[|p|^2].
\]
The three ratios thus give \eqref{eq:spectral-level-two-eigenvalues}.
The second exceeds the first by $n^2/3$; the third exceeds the
second by $(n-2)(n+1)/6$. Taking the largest present block proves
the stated formula for $\Lambda_2$.
\end{proof}

\subsection{Why this hierarchy cannot prove a sharp finite bound}
The upper bounds converge, but the sharp value itself has an
obstruction to a finite certificate of this particular Hermitian
sum-of-squares form. The following statement holds for every nonzero
homogeneous holomorphic polynomial, not just the Vandermonde.

\begin{proposition}[Finite sharpness obstruction]\label{prop:spectral-no-finite-sharpness}
Let $f$ be nonzero and homogeneous of degree $N\ge1$ on $\C^d$, and
construct $\Lambda_r$ as above. Put
$S=\max_{\|u\|=1}|f(u)|^2$. For a finite $r\ge0$, equality
$\Lambda_r=S$ is possible if and only if
$f(u)=a\ell(u)^N$ for a unit-norm linear functional $\ell$ and
a nonzero scalar $a$. In particular, for the centered Vandermonde
with $n\ge3$,
\begin{equation}\label{eq:spectral-strictness}
 n^N\Lambda_r>M_n\qquad\text{for every finite }r.
\end{equation}
\end{proposition}
\begin{proof}
Suppose $\Lambda_r=S$. The polynomial
$S\|u\|^{2(N+r)}-|f(u)|^2\|u\|^{2r}$ is a Hermitian sum of squares.
Choose a unitary coordinate change sending a maximizing point to
$e_1$. Such a change preserves a sum-of-squares representation.
Write $f(u)=\sum_{|\gamma|=N}a_\gamma u^\gamma$ in the new
coordinates, with $a_0:=a_{(N,0,\ldots,0)}=f(e_1)$; thus
$|a_0|^2=S>0$. In its unique normalized Hermitian coefficient
matrix, the diagonal entry indexed by $u_1^{N+r}$ is zero.
A positive-semidefinite matrix with a zero diagonal entry has the
entire corresponding row and column zero.

For any $\gamma\ne(N,0,\ldots,0)$ with $|\gamma|=N$, let
$\beta=\gamma+(r,0,\ldots,0)$. The entry in that row and column
$\beta$ equals
\[
 -\frac{\overline{a_0}a_\gamma}{\sqrt{\binom{N+r}\beta}}.
\]
Indeed, in \eqref{eq:spectral-matrix} the row indexed by
$(N+r,0,\ldots,0)$ can have only the column
$\alpha=(r,0,\ldots,0)$ nonzero, and the norm-power term is diagonal.
Every $a_\gamma$ other than $a_0$ must consequently vanish. Hence
$f=a_0u_1^N$ in these coordinates, proving necessity for every
$r\ge0$, including $r=0$.

Conversely, after a unitary change let $f=a u_1^N$. Expanding
$|a|^2\big(\|u\|^{2(N+r)}-|u_1|^{2N}\|u\|^{2r}\big)$ in
monomial absolute squares gives nonnegative coefficients: the
coefficients are $|a|^2\binom{N+r}\beta$ minus
$|a|^2\binom r{\beta-(N,0,\ldots,0)}$, when the latter exists.
The first multinomial coefficient is at least the second, by
injecting length-$r$ words with the specified multiplicities into
length-$(N+r)$ words by appending $N$ copies of the first letter.
Thus its Hermitian coefficient matrix is positive semidefinite,
and $\Lambda_r\le |a|^2=S$; evaluation at $e_1$ gives equality.
For $n\ge3$, the linear factors $z_1-z_2$ and $z_1-z_3$ restricted
to $\mathcal H$ are nonproportional, so the centered Vandermonde
cannot be a power of one linear form.
\end{proof}

Thus this spectral hierarchy gives arbitrarily sharp strict upper
bounds. It does not supply a finite sharp Hermitian certificate for
$M_6$ or any other $n\ge3$. This does not exclude other real
sum-of-squares certificates, certificates using additional
relations, or exact-value identification by combining convergent
bounds with a critical-value polynomial.

\subsection{An exact three-point spectral calculation}
For $n=3$, a unitary Fourier basis of the centered plane gives
$f(u,v)=u^3+v^3$ up to a constant phase and a unitary change of $v$.
For example, take $\omega=e^{2\pi i/3}$ and
$(Bu)_j=(\omega^ju+\omega^{-j}v)/\sqrt3$, $j=0,1,2$;
the Vandermonde is a constant of modulus one times $u^3-v^3$.
Changing the phase of $v$ gives the stated plus sign.
At $r=1$, in the degree-four basis
$(u^4,2u^3v,\sqrt6u^2v^2,2uv^3,v^4)$,
\[
 A_1=\begin{pmatrix}
 1&0\\0&1/2\\0&0\\1/2&0\\0&1
 \end{pmatrix},\qquad
 A_1^*A_1=\frac54I_2.
\]
At $r=2$ the three columns have disjoint monomial supports, giving
\[
 A_2^*A_2=\operatorname{diag}\left(\frac{11}{10},
                           \frac45,\frac{11}{10}\right).
\]
Consequently $\Lambda_0=2$, $\Lambda_1=5/4$, and
$\Lambda_2=11/10$, giving the exact certified upper bounds
$54$, $135/4$, and $297/10$ for $M_3$. The unit triangle gives
$M_3\ge27$, and the previously proved exact value is $M_3=27$.
The strict finite gaps illustrate
Proposition~\ref{prop:spectral-no-finite-sharpness}.

A finite sharp certificate becomes possible when squares of mixed
polynomials in $u,v,\overline u,\overline v$ are permitted:
\begin{equation}\label{eq:three-real-sos}
\begin{split}
 (|u|^2+|v|^2)^3-|u^3+v^3|^2
 ={}&|u^2\overline v-\overline u v^2|^2\\
 &+2|u^2\overline v|^2+2|\overline u v^2|^2.
\end{split}
\end{equation}
Expansion proves the identity and hence $S\le1$; equality at
$(u,v)=(1,0)$ gives $M_3=27$. These are real polynomial squares
after splitting each term into its real and imaginary parts.
They are not squares of holomorphic polynomials, so the identity
does not contradict the finite sharpness obstruction.

% Add to the main bibliography:
% \bibitem{Beauzamy}
% B.~Beauzamy, \emph{Extremal products in Bombieri's norm},
% January 1995.
% \url{https://www.scmsa.eu/BB_Publications/BB_Extremal_products_Bombieri_norm.pdf}.

\subsection{Exact six-point spectral certificates at low levels}
For a directly reproducible calculation, retain all six ambient
variables and expand
\[
 V_6(z)=\sum_{\sigma\in S_6}\operatorname{sgn}(\sigma)
             \prod_{i=1}^6z_i^{\sigma(i)-1}
        =\sum_{|\gamma|=15}c_\gamma z^\gamma .
\]
This differs from the centered implementation only in the choice of
ambient space; evaluation on the unit sphere still gives a valid
upper bound for $M_6/6^{15}$. For $|\alpha|=|\eta|=r$, put
\[
 D_{\alpha\alpha}=\alpha!,\qquad
 H_{\alpha\eta}=\sum_{|\delta|=15+r}
            c_{\delta-\alpha}c_{\delta-\eta}\,\delta!,
\]
where coefficients with negative indices vanish and multi-index
factorials denote products of ordinary factorials. The multiplication
Gram matrix is
\[
 \widetilde G_r=\frac{r!}{(15+r)!}D^{-1/2}HD^{-1/2}.
\]
Consequently an integer $U_r$ satisfies $M_6<U_r$ whenever
\[
 U_r(15+r)!\,D-r!\,6^{15}H\succ0.
\]
The ancillary program \texttt{verify\_six\_spectral.py} regenerates
$H$ from the $720$ signed permutation monomials and verifies this
last inequality by fraction-free Bareiss elimination. Every division
is checked to be exact, and every leading principal determinant is
checked to be positive; thus Sylvester's criterion supplies an
integer-arithmetic certificate independent of floating-point
eigenvalue calculations.

\begin{center}
\begin{tabular}{rrrl}
\toprule
$r$ & Matrix size & Certified $U_r$ &
 $6^{15}\lambda_{\max}(\widetilde G_r)$, numerical\\
\midrule
1 & 6   & $2\,236\,740$ & $2\,236\,739.38$\\
2 & 21  & $1\,008\,726$ & $1\,008\,725.60$\\
3 & 56  & $614\,007$    & $614\,006.89$\\
4 & 126 & $405\,332$    & $405\,331.53$\\
5 & 252 & ---           & $290\,511.27$ (estimate)\\
6 & 462 & ---           & $235\,776.01$ (estimate)\\
\bottomrule
\end{tabular}
\end{center}
The default exact verifier runs levels $1$ through $4$ in about
three seconds in the present environment and needs no third-party
Python package or cached matrix. An optional NumPy calculation
reproduces the floating-point estimates; the values displayed for
levels $5$ and $6$ are not used as certified inequalities here.
The JSON record includes SHA-256 hashes of the regenerated integer
Gram matrices and the Bareiss pivots. These hashes identify the
computed data; positivity and exact division, rather than the hashes,
constitute the proof. These low levels demonstrate working
six-point spectral certificates. All their upper bounds exceed
$89\,309$ and the exact value in Theorem~\ref{thm:six-global};
they do not supply the sharp six-point proof.


\section{From the six-point proof to a uniform certification method}
\label{sec:six-to-uniform}

The six-point proof separates three questions: covering every stationary
shape, enforcing actual conjugation, and comparing values.  The first
question uses the published algebraic computation; the last two now
have independent rational certificates.  The same separation gives a
uniform strategy for larger point counts, without prescribing polygon
or reflection symmetry.

\subsection{A coefficient system with only \texorpdfstring{$n-2$}{n-2} nontrivial equations}
Let $P(t)=\prod_{i=1}^n(t-w_i)$ be monic and centered, with distinct
roots, and put
\[
 r_i=\frac{P''(w_i)}{2P'(w_i)},\qquad
 D(Y)=\prod_i(Y-r_i),\qquad
 K(Y)=\sum_i w_i\prod_{j\ne i}(Y-r_j).
\]
Both $D$ and $K$ have coefficients rational in the coefficients of $P$,
with denominators supported on its discriminant.  They can be computed
by resultants and a first derivative, without choosing or ordering roots:
\[
 D(Y)=\frac{\operatorname{Res}_t(P,YP'-P''/2)}
                  {\operatorname{Res}_t(P,P')},\qquad
 K(Y)=\left.\frac{\partial}{\partial T}
 \frac{\operatorname{Res}_t(P,(Y+Tt)P'-P''/2)}
      {\operatorname{Res}_t(P,P')}\right|_{T=0}.
\]
The formula for $K$ follows by differentiating
$\prod_i(Y+Tw_i-r_i)$.

Every physical force-$1$ stationary configuration has $r_i=\overline
w_i$ and consequently satisfies the dual force equations.  If the
$r_i$ are distinct, those dual equations are equivalent to
\begin{equation}\label{eq:uniform-dual-coefficients}
                 D''(Y)=2K(Y).
\end{equation}
Indeed, at $Y=r_i$, $K(r_i)=w_iD'(r_i)$; hence the polynomial
$D''-2K$ vanishes at all $n$ distinct roots if and only if the dual
forces hold.  Its degree is at most $n-2$.
The leading coefficient is automatically zero, because
$\sum_i w_ir_i=n(n-1)/2$.  Thus \eqref{eq:uniform-dual-coefficients}
has only $n-2$ nontrivial scalar coefficient equations.  Clearing
denominators must be followed by discriminant saturation.

There are $n-1$ centered coefficients.  On each chart normalize the
first nonzero elementary coefficient to one, leaving $n-2$ complex
unknowns.  Treat every zero-coefficient chart as well.  The dual
coefficient equations provide a compact permutation quotient of the
necessary stationary equations.  Actual conjugation must still be
checked; the six-point spurious branches show why it cannot be omitted.
The positive trace formulation in Section~\ref{sec:trace-stationary}
provides an exact alternative to rootwise conjugate tests.

\subsection{Finite branch certificates and finite value certificates}
When a saturated chart decomposition is zero-dimensional, isolate
every complex parameter root by disjoint rational disks.  Propagate
coefficient errors and isolate the roots of $P$ uniformly over each
parameter disk, as in Section~\ref{sec:six-disk-certificate}.  For
$Q=\sum_i|w_i|^2$ and $N=n(n-1)/2$, physical stationarity requires
\[
 QP''(w_i)-2N\overline{w_i}P'(w_i)=0,
\]
and the value on the original moment-$n$ sphere is
\[
                 (n/Q)^N|\operatorname{Disc}P|.
\]
A single rigorously nonzero force residual excludes a branch.  For
each remaining branch, an outward product bound proves strict
inferiority when there is a value gap.  Algebraic value comparison or
an exact identity handles branches attaining the same value.
The degree counts and the disk inequalities, rather than numerical
root seeds, certify coverage.

Zero-dimensionality of every shape chart is not assumed for arbitrary
$n$.  If a positive-dimensional component occurs, point disks alone
cannot exhaust it.  The exact critical-value construction of
Theorem~\ref{thm:exact-algorithm} remains available: it eliminates to
a nonzero polynomial in the objective value, since the value is
constant on each relevant irreducible critical component.
This proves finiteness of critical values without requiring a finite
classification of shapes.

Combine that polynomial with either the coefficient-moment enclosure
or the spectral enclosure.  Increase the order until a rational
outward interval contains only one distinct real root of the
critical-value polynomial.  The interval is known to contain $M_n$;
convergence and the positive separation between distinct roots make
this procedure terminate for each fixed $n$.  The output is an exact
algebraic specification of $M_n$, including an isolating interval.
It does not require knowing its extremal geometry in advance.

This gives a unified exact-value framework, together with a practical
branch route demonstrated at $n=6$.  The coefficient equations, trace
positivity, radial stability cuts, moment bounds and spectral bounds
can be combined before expensive elimination.  Their cost still grows
with $n$; the framework is not a finite closed formula for all $M_n$.
The ten-point example and the instability of sufficiently large
central polygons prevent the six-point geometry from serving as a
universal extremal ansatz. The fixed-ring gap in
Section~\ref{sec:fixed-ring-obstruction} proves a stronger limitation:
every fixed number of regular rings loses a positive amount at the
leading $n^2$ scale. The polar construction requires a growing number
of rings and provides explicit all-$n$ lower values. Section~
\ref{sec:uniform-linear-constant} identifies the remaining global
linear coefficient with a universal microscopic energy constant.
Even resolving that constant would leave finite-$n$ boundary corrections
and extremal shapes to determine.

% Fragment for the existing research manuscript; not a standalone document.
\section{The linear asymptotic constant and the crystallization question}
\label{sec:uniform-linear-constant}
The general-$n$ problem is exactly a quadratic weighted Fekete problem,
including its finite-$n$ normalization. This observation permits a sharper
asymptotic statement than an unspecified $O(n)$ remainder. The result
below is a specialization of the classical renormalized-energy theory,
not a new proof of that theory and not a finite closed formula for $M_n$.

Set $N=n(n-1)/2$ and define the unconstrained quadratic energy
\[
 H_n(x)=-\sum_{i\ne j}\log|x_i-x_j|+n\sum_i|x_i|^2,
 \qquad e_n=\min_{x\in\mathbb C^n}H_n(x).
\]
\begin{proposition}[Exact constrained--unconstrained equivalence]
\label{prop:uniform-gaussian-equivalence}
For every $n\ge2$,
\begin{equation}\label{eq:uniform-gaussian-equivalence}
 e_n=-\log M_n+N\left[1-\log\frac{n-1}{2n}\right].
\end{equation}
Every minimizer of $H_n$ is centered and satisfies
$\sum_i|x_i|^2=(n-1)/2$. Multiplication by
$\sqrt{2n/(n-1)}$ maps its shape to a maximizer defining $M_n$, and
this correspondence is reversible.
\end{proposition}
\begin{proof}
Centering preserves pair distances and decreases the quadratic term,
so a minimizer is centered. In the constrained problem, centering and
then rescaling to squared radius $n$ strictly increases the positive
Vandermonde product unless the centroid was already zero. Thus its
maximizers are centered as well. For a centered shape $z$ with
$\sum_i|z_i|^2=n$, scaling it by $t>0$ gives
\[
 H_n(tz)=-\log\Delta_n(z)-N\log t^2+n^2t^2.
\]
The strictly convex function of $t^2$ on the right is minimized at
$t^2=N/n^2=(n-1)/(2n)$, independently of the shape. At this scale,
its last two terms are $N-N\log((n-1)/(2n))$.
Minimizing the remaining shape term proves the identity and the
correspondence. Existence also follows from compactness of the
centered sphere and positivity of its Vandermonde maximum.
\end{proof}

We specify the convention for the universal constant. Let
$\mathcal A_1$ be the class of curl-free planar electric fields satisfying
\[
 \operatorname{div}E=2\pi\left(\sum_{p\in\Lambda}\delta_p-1\right),
\]
with $\Lambda$ discrete and
$\sup_{R\ge1}\#(\Lambda\cap B_R)/R^2<\infty$.
For smooth nonnegative cutoffs $\chi_R$ supported in
$K_R=[-R,R]^2$, define
\[
 W(E)=\limsup_{R\to\infty}\frac1{|K_R|}
 \lim_{\eta\downarrow0}\left[
 \frac12\int_{\mathbb R^2\setminus\bigcup_{p\in\Lambda}B(p,\eta)}
 \chi_R|E|^2+\pi\log\eta\sum_{p\in\Lambda}\chi_R(p)\right],
 \qquad W_*:=\min_{\mathcal A_1}W.
\]
Here $\chi_R=1$ except within a bounded-width boundary layer and has
uniformly bounded gradient. This is the Sandier--Serfaty normalization.
The minimum $W_*$ is finite and attained, but its identification with
the triangular-lattice value is an additional crystallization question.

\begin{theorem}[Existence of the linear asymptotic coefficient]
\label{thm:uniform-linear-coefficient}
Put $c=\tfrac12\log2-\tfrac14$. Then
\begin{equation}\label{eq:uniform-linear-coefficient}
 \boxed{\displaystyle
 \log M_n=cn^2+\frac n2\log n+\gamma n+o(n),\qquad
 \gamma=-\frac{W_*}{\pi}-\frac12\log(2\pi).}
\end{equation}
In particular, the limit defining $\gamma$ exists. Moreover,
\begin{equation}\label{eq:uniform-linear-coefficient-bounds}
 \log\frac{\sqrt3\,\Gamma(1/3)^3}{4\pi^{3/2}}
 \le\gamma\le\frac12\log\pi.
\end{equation}
The lower endpoint is approximately $0.4021788952$ and the upper
endpoint is approximately $0.5723649429$. Equality at the lower endpoint
holds if and only if the triangular lattice of density one attains $W_*$.
\end{theorem}
\begin{proof}
The external input is Sandier--Serfaty's ground-state convergence
theorem, Theorem~2(C) and equations~(1.23), (1.40) in
\href{https://arxiv.org/pdf/1201.3503}{\emph{2D Coulomb Gases and the
Renormalized Energy}}. In the above convention it states, for a regular
confining potential $V$ with equilibrium density $m_0$, that
\[
 \min H_n=n^2I_V(\mu_0)-\frac n2\log n
 +n\left[\frac{W_*}{\pi}-\frac12\int m_0\log m_0\right]+o(n).
\]
For $V(x)=|x|^2$, all regularity and growth assumptions hold. Its
equilibrium measure is the uniform probability measure on the unit disk,
with density $m_0=1/\pi$. For example, its logarithmic potential is
$(1-|x|^2)/2$ inside the disk and $-\log|x|$ outside, so the equilibrium
inequality is verified directly. Its self-energy is $1/4$ and its
quadratic moment is $1/2$, giving $I_V(\mu_0)=3/4$. Also
$\int m_0\log m_0=-\log\pi$. Therefore
\[
 e_n=\frac34n^2-\frac n2\log n
 +n\left(\frac{W_*}{\pi}+\frac12\log\pi\right)+o(n).
\]
On the other hand, elementary expansion gives
\[
 N\left[1-\log\frac{n-1}{2n}\right]
 =\frac{1+\log2}{2}n^2-\frac{\log2}{2}n-\frac14+O(n^{-1}).
\]
Substitution into \eqref{eq:uniform-gaussian-equivalence} proves
\eqref{eq:uniform-linear-coefficient}.

For the upper bound, the elementary Hadamard bound already proved in
the manuscript is
$M_n\le e^N(n/N)^N\prod_{j=0}^{n-1}j!$. Euler--Maclaurin gives
\[
 \log\prod_{j=0}^{n-1}j!
 =\frac{n^2}{2}\log n-\frac{3n^2}{4}
 +\frac n2\log(2\pi)+O(\log n).
\]
It follows that $\log M_n\le cn^2+(n/2)\log n+
(n/2)\log\pi+O(\log n)$; division by $n$ and the limit prove the
upper endpoint.

B\'etermin--Sandier compute, in Section~7.3 of
\href{https://arxiv.org/pdf/1404.4485}{\emph{Renormalized Energy and
Asymptotic Expansion of Optimal Logarithmic Energy on the Sphere}},
the density-one triangular-lattice value
\[
 W_\triangle=\pi\left[\log\pi-\frac12\log3
 -3\log\Gamma(1/3)+\frac32\log2\right].
\]
Since $W_*\le W_\triangle$, substituting this value in the expression
for $\gamma$ proves the lower endpoint. The affine dependence of
$\gamma$ on $W_*$ also proves the stated equality equivalence.
\end{proof}

This result yields, for every $\varepsilon>0$, an asymptotic lower bound
\[
 M_n\ge\exp\left(cn^2+\frac n2\log n+
 (\gamma_\triangle-\varepsilon)n\right)
\]
for all sufficiently large $n$, where
$\gamma_\triangle=\log(\sqrt3\,\Gamma(1/3)^3/(4\pi^{3/2}))$.
No explicit threshold or finite-$n$ equality formula follows from this
statement. Even identifying the linear coefficient with
$\gamma_\triangle$ would leave finite boundary corrections and each
finite-$n$ shape unresolved.

\subsection{The explicit polar value is still strictly suboptimal}
The growing number of rings avoids the fixed-$K$ obstruction, but does
not remove the microscopic energy deficit. This can be proved without
making any floating-point comparison of special functions.
\begin{corollary}\label{cor:polar-linear-deficit}
For the explicit values $D_n$ of Theorem~\ref{thm:polar-all-n},
\[
 \lim_{n\to\infty}\frac1n\log\frac{M_n}{D_n}
   =\gamma-\gamma_{\rm pol}>\frac1{200}.
\]
Consequently $D_n/M_n<e^{-n/200}$ for all sufficiently large $n$.
This construction matches the leading two terms and still cannot
be the requested finite-$n$ maximum formula.
\end{corollary}
\begin{proof}
Subtract \eqref{eq:polar-three-term} from
\eqref{eq:uniform-linear-coefficient}. The triangular trial bound gives
\[
 \gamma-\gamma_{\rm pol}\ge
 \frac12\log\frac{\Gamma(1/3)^6e}{32\pi^3}.
\]
Here is an entirely rational certificate that the right-hand side
exceeds $1/200$. Substitution $t=u^3$ gives
$\Gamma(1/3)=3\int_0^\infty e^{-u^3}\,du$. An odd Taylor truncation
is below the exponential for every nonnegative argument, whence
\[
 \Gamma(1/3)>3\sum_{j=0}^{41}
 \frac{(-1)^j2^{3j+1}}{(3j+1)j!}>\frac{26788}{10000}.
\]
Also $e>\sum_{j=0}^{10}1/j!$ and $\pi<355/113$. The latter can
be verified from Machin's identity and outward alternating-series
bounds on the two arctangents. Integer arithmetic verifies
\[
 \frac{(26788/10000)^6\sum_{j=0}^{10}1/j!}
      {32(355/113)^3}>\frac{100}{99}>e^{1/100},
\]
where the last strict inequality follows by comparing the exponential
series with the geometric series. Taking half the logarithm proves
the stated strict gap. All these rational comparisons are checked
independently by \texttt{verify\_polar\_nonoptimal.py}.
\end{proof}

\subsection{A simultaneous moment limit for the linear coefficient}
Let $C_{n,k}$ be the complete integer-coefficient moment in
\eqref{eq:coefficient-moment}, and put
$U_{n,k}=\bigl(n^{kN}C_{n,k}/(kN)!\bigr)^{1/k}$.
The existing moment enclosure implies the uniform estimate
\begin{equation}\label{eq:uniform-moment-log-gap}
 0\le\log U_{n,k}-\log M_n
 \le\frac{n-2}{k}
 \log\frac{e(kN+n-2)}{n-2}\qquad(n\ge3).
\end{equation}
Indeed, its exact endpoint ratio is
$\binom{kN+n-2}{n-2}^{1/k}$, and
$\binom{s}{r}\le(es/r)^r$.
If $k_n$ is any integer sequence with $k_n/\log n\to\infty$, then
\begin{equation}\label{eq:uniform-moment-linear-limit}
 \boxed{\displaystyle
 \gamma=\lim_{n\to\infty}\left[
 \frac1{k_n n}\log\frac{n^{k_nN}C_{n,k_n}}{(k_nN)!}
 -cn-\frac12\log n\right].}
\end{equation}
To prove this, note that $N/(n-2)\le n$ for $n\ge3$, so
\eqref{eq:uniform-moment-log-gap} divided by $n$ is at most
$k_n^{-1}\log(e k_n(n+1))$, which tends to zero under the stated
condition. Equation~\eqref{eq:uniform-linear-coefficient} finishes the
argument. For example, $k_n=\lceil(\log n)^2\rceil$ is sufficient.
This supplies a rigorous connection between the coefficient hierarchy
and the microscopic energy constant. It remains a limit of finite sums,
with rapidly growing coefficient calculations, rather than a finite
closed product.

\section{Reproducibility and interpretation}
The standard-library program \texttt{verify\_constructions.py} computes
the normalized integer-lattice products exactly, evaluates the
rational bounds, and checks the ten-point comparison.
It uses explicit exceptions, which remain active under Python
\texttt{-O}. Run it with
\begin{verbatim}
python verify_constructions.py --max-n 30
python verify_hermite.py
python verify_saved_candidates.py
python verify_polygon_perturbations.py
python verify_global_grid.py
python verify_stationary_polynomial.py
python verify_critical_value.py
python -O verify_moment_hierarchy.py
python -O verify_framework.py
python -S -O certify_fs6_disks.py --verify
python -S -O verify_fs6_components_exact.py
python -S -O verify_general_radial_caps.py
python -S -O verify_all_n_deletion.py
python -S -O all_n_structure_radial_verify.py
python -S -O all_n_uniform_checks.py
python -S -O all_n_rings_verify.py
python -S -O verify_polar_nonoptimal.py
\end{verbatim}
The optional numerical search is reproduced separately with
\texttt{python explore\_candidates.py}; it requires NumPy.
The JSON file separates exact fractions from floating-point display
values. Finite arithmetic checks support the formulas; the Gaussian,
interpolation and limiting arguments are proved analytically in this note.
The new perturbation checker uses exact cyclotomic field arithmetic
for the ten-point constants and independent root-of-unity power sums,
and verifies the rational six-point neighborhood constants.
The global-grid checker completely enumerates the two stated grids.
The stationary-polynomial checker verifies force remainders and trace
matrices, including an exact spurious six-root pairing. The critical-value
checker verifies the complete three-point elimination identities and
the Hermite resultant recurrence. An optional SageMath eliminant
generator is included as \texttt{critical\_value\_eliminant.sage.py};
it has not been run in this environment and is not a six-point
completeness certificate.
The new coefficient checker completely expands the 25 stated powers,
verifies the Gaussian normalization and an independent three-point
formula, and records rational outward enclosures. The framework checker
verifies the exact three-point spectral matrices, the mixed-variable
sharp identity, the first spectral-level formula, and the rational
six-point improvement in the grid constant. The six-point root-disk verifier supplies a genuine maximum
certificate when combined with the explicitly cited exhaustive
Faug\`ere--Svartz decomposition. Its verification uses only standard
library integer and rational arithmetic, including with Python
\texttt{-S -O}. The separate exact component checker verifies the
coefficient equations and discriminant nonvanishing; it does not
recompute the external decomposition. The general radial checker
checks the polynomial divided-difference formulas and integer grid
caps. The general limiting and finite-value arguments are proved
analytically in this note.

The October 5 extension adds exact subset-deletion identities,
stronger local radial moments, complete partition-orbit checks and the
rational polar construction. The polar checker covers finite comparisons
through 120 points; the uniform remainder and larger-$n$ comparison
are proved analytically. The fixed-disk proof is self-contained. The
global linear asymptotic uses the cited external renormalized-energy
theorem, not the finite arithmetic checks.

% Fragment: independently checked small-n examples and a bordered trace test.
% This file does not modify the main manuscript.
\section{Exact examples and uniform finite certificates}
\label{sec:examples-bordered-trace}

The following examples separate a short global inequality from an exact
stationary classification. They do not assert that the regular polygon
remains optimal for arbitrary $n$.

\subsection{Two, three, and four points}
For two points,
\[
 |z_1-z_2|^2\le2(|z_1|^2+|z_2|^2)=4,
\]
with equality precisely for $z_2=-z_1$ and $|z_1|=1$.
Thus $M_2=4$. For three points, center first and use orthonormal
Fourier coordinates $u,v$ on the centered complex plane. The Vandermonde is
$u^3+v^3$ up to a constant of modulus one, and
\[
 (|u|^2+|v|^2)^3-|u^3+v^3|^2
 =|u^2\overline v-\overline u v^2|^2
   +2|u^2\overline v|^2+2|\overline u v^2|^2.
\]
Consequently $M_3=27$, with equality for a regular triangle.

Here is an independent short proof of the four-point case. Its
Hadamard change of variables is also used in Appendix C of
J.~Wang, \emph{The Five-Point Case of Mordell's Discriminant
Inequality}, arXiv:2610.00358. The four-point value is classical;
no novelty claim is made for it.

\begin{lemma}\label{lem:l1-three-variational}
For $A,B,C\in\C$ with $|A|+|B|+|C|=1$,
\[
 |(A-B)(B-C)(C-A)|^2\le\frac1{16}.
\]
Equality holds precisely when, after permutation,
$(A,B,C)=(0,\eta/2,-\eta/2)$ with $|\eta|=1$.
\end{lemma}
\begin{proof}
The compact constraint set has a positive maximum. A collision
cannot attain that maximum. If $A=0$, writing $x=|B|$, $y=|C|$
gives
\[
 |(A-B)(B-C)(C-A)|^2
 =x^2y^2|B-C|^2\le x^2y^2(x+y)^2\le\frac1{16}.
\]
The two equalities require $x=y=1/2$ and $C=-B$.

At an interior maximum all three moduli are positive. Differentiating
the logarithm of the product under the modulus constraint gives
\[
 \sum_{j\ne i}\frac1{A_i-A_j}
   =\kappa\frac{\overline A_i}{|A_i|},\qquad \kappa=3.
\]
The value of $\kappa$ follows by multiplying by $A_i$ and summing.
Summing the force equations instead shows that the three unit
directions add to zero. They therefore have mutual angles $120^\circ$.
After rotation and relabeling, write
\[
 A=x,\quad B=y\omega,\quad C=z\omega^2,
 \qquad x,y,z>0,\quad x+y+z=1,
 \quad \omega=e^{2\pi i/3}.
\]
The force at $A$ is real. Its imaginary part vanishes exactly when
\[
 \frac{y}{x^2+xy+y^2}-\frac{z}{x^2+xz+z^2}=0,
 \quad\text{or}\quad (y-z)(x^2-yz)=0.
\]
The cyclic equations hold as well. If two of $x,y,z$ agree and the
third differs, the equation at either of the equal radii forces the
third to agree also. If all three differ, then $x^2=yz$ and
$y^2=zx$, which again force $x=y$. Thus $x=y=z=1/3$.
The interior stationary product is $1/27<1/16$, so the global
maximum and its equality cases are the boundary ones already found.
\end{proof}

\begin{proposition}\label{prop:four-variational-exact}
$M_4=256$, with equality precisely for a regular square on the
unit circle, up to rotation and permutation.
\end{proposition}
\begin{proof}
Center a maximizing configuration. Centering lowers its quadratic
moment, and a positive rescaling would increase the product unless
the original center were zero. Every centered quadruple has the
unique representation
\[
 (z_1,z_2,z_3,z_4)
 =(a+b+c,a-b-c,-a+b-c,-a-b+c),
\]
where $\sum_i|z_i|^2=4(|a|^2+|b|^2+|c|^2)$. Moreover,
\[
 V_4=64(a^2-b^2)(a^2-c^2)(b^2-c^2).
\]
Apply Lemma~\ref{lem:l1-three-variational} to
$A=a^2,B=b^2,C=c^2$. Their moduli sum to one, giving
$\Delta_4\le4096/16=256$. Equality means that one of $a,b,c$ is
zero and the other two have squared values opposite each other,
of modulus $1/2$. The resulting four points form a centered
unit square. Such a square attains $256$.
\end{proof}

\subsection{All physical four-point stationary values in a cubic selector}
\label{sec:four-critical-values}
The same Hadamard variables give a complete list of real physical
stationary values, not merely an upper bound. A centered
collision-free stationary quadruple has
$|a|^2+|b|^2+|c|^2=1$.
If $abc\ne0$, the map $(a,b,c)\mapsto(a^2,b^2,c^2)$ is locally
invertible over $\C$. The preceding interior force calculation
therefore applies to every stationary point, regardless of stability,
and gives $|A|=|B|=|C|=1/3$ with directions separated by
$120^\circ$. The discriminant is $4096/27$.

If, say, $c=0$, then $a,b\ne0$: otherwise some point roots collide.
The first variation in $c$ is automatically zero, since both the
Vandermonde expression and the constraint are even in $c$.
Put $x=|a|^2$, $y=|b|^2$, $x+y=1$, and let $\theta$ be the
argument difference of $A=a^2$ and $B=b^2$. The product divided
by $4096$ is
\[
                 x^2y^2(x^2+y^2-2xy\cos\theta).
\]
Its phase derivative forces $\sin\theta=0$.
For opposite directions it is $x^2y^2$, whose stationary point
is $x=y=1/2$, giving the square value $256$.
For the same direction it is $p^2(1-4p)$, where $p=xy$.
The collision value $x=y$ is excluded; hence differentiation
forces $p=1/6$, giving the collinear value $1024/27$.
Each configuration satisfies the force equations: the square and
central triangle by geometric sums, and the collinear case by the
Hermite polynomial identity. Thus the complete physical
stationary-value set is
\[
                      \left\{256,\frac{4096}{27},\frac{1024}{27}\right\}.
\]
In particular, a sufficient exact value-selector polynomial is
\begin{equation}\label{eq:four-physical-selector}
          p_4^{\rm phys}(T)=(T-256)(27T-4096)(27T-1024).
\end{equation}
This degree-three real selector is not asserted to equal the
eliminant of the full complexified stationary system. It is
obtained by classifying the actual conjugate-paired configurations.

\subsection{A fully executed five-point classification}
Theorem 1.1 of Wang's arXiv:2610.00358 proves $M_5=3125$,
with equality precisely for a regular unit pentagon. Its
accompanying \texttt{n5\_complete\_verification.py} was run in full
on 5 October 2026 with Python assertions enabled. The run rebuilt
the resultant equations and saturated Gr\"obner bases, applied
FGLM and Sturm counting, and verified the exact discriminant bounds;
it exited successfully. The complete output is saved separately.
This is an executed exact classification, rather than a numerical
search for the pentagon.

There are five physical stationary shapes modulo rotations and
permutations. Their objective values are
\[
 3125,\qquad \frac{9765625}{4096},\qquad
 \frac{84375}{1024},\qquad D(H_-),\qquad D(H_+),
\]
where $H_-,H_+$ are the two real roots of
\[
 10H^4-65H^3+140H^2-97H+6=0
\]
and
\[
 D(H)=\frac{1084589409375}{385875968}
 -\frac{1391385515625}{192937984}H
 +\frac{2547927984375}{385875968}H^2
 -\frac{321987234375}{192937984}H^3.
\]
An independent standard-library checker uses rational Sturm and
interval arithmetic to verify
$2347<D(H_-)<2348$ and $801<D(H_+)<802$.
It also verifies both classical rational values are below $3125$.
Completeness here comes from the fully rebuilt classification, not
from this smaller interval checker alone.


\subsection{An elementary sharp bound when one point is at the origin}
The following result removes all origin-containing branches in the seven-point
comparison by a single inequality.
\begin{theorem}[The origin-constrained maximum for $2\le n\le7$]
Let $2\le n\le7$, and let $z_1,\ldots,z_n\in\mathbb C$ satisfy
$z_1=0$ and $\sum_i|z_i|^2=n$. Then
\[
 \prod_{i<j}|z_i-z_j|^2\le
 (n-1)^{n-1}\left(\frac n{n-1}\right)^{n(n-1)/2}.
\]
Equality holds precisely when the other $n-1$ points form a regular polygon
on the circle of radius $\sqrt{n/(n-1)}$, up to rotation and permutation.
For $n=2$, a regular one-gon means a single point on that circle.
\end{theorem}
\begin{proof}
Put $m=n-1$, $u_i=\sqrt{m/n}\,z_{i+1}$ and $t_i=|u_i|^2$.
Thus $\sum_i t_i=m$. If any $t_i=0$, or if two points coincide, the claimed
upper bound is immediate. Otherwise, the Vandermonde determinant and
Hadamard's inequality give
\[
 \prod_i|u_i|^2\prod_{i<j}|u_i-u_j|^2
 =\left|\det(u_i^j)_{1\le i,j\le m}\right|^2
 \le\prod_{i=1}^m S_m(t_i),\qquad
 S_m(t)=\sum_{j=1}^m t^j.
\]
For $1\le m\le6$, $\log S_m$ is strictly concave on $(0,\infty)$.
Indeed, write $R_m=S_m'^2-S_mS_m''$, so
$(\log S_m)''=-R_m/S_m^2$. Direct exact expansion shows that
$R_m(1+x)$ has nonnegative coefficients and positive constant term.
On $[0,1]$, all Bernstein coefficients of $R_m$ are positive when
$m\le5$. For $m=6$, this holds on each of $[0,1/2]$ and $[1/2,1]$.
These finite integer and rational coefficient checks are reproduced by
\texttt{origin\_polygon\_verify.py}, using only the Python standard library.
For example,
\[
 R_6(t)=1+2t+2t^2-5t^4-14t^5+8t^7+11t^8+10t^9+6t^{10}.
\]
Jensen's inequality now gives $\prod_iS_m(t_i)\le S_m(1)^m=m^m$.
Restoring scale multiplies the squared distance product by
$(n/m)^{n(n-1)/2}$, proving the bound.
Strict concavity forces $t_i=1$ at equality. For $m\ge2$, equality in
Hadamard additionally requires distinct rows to be orthogonal. With
$q=u_i\overline{u_j}$, this means
$\sum_{\ell=1}^m q^\ell=0$. Since $|q|=1$ and $q\ne1$, it follows that
$q^m=1$. Thus the $m$ distinct points are precisely a rotated set of
$m$th roots of unity. The case $m=1$ is immediate.
\end{proof}
In particular, the origin-constrained seven-point maximum is exactly
$7^{21}/6^{15}$. This theorem alone does not establish the unrestricted
seven-point maximum: one must still compare every stationary configuration
having no point at the origin. Translation of an arbitrary point to the
origin changes the moment constraint and cannot be used to bypass that step.

The short argument has a precise limit. For $S_7$,
$(\log S_7)''(1)=0$ and $(\log S_7)'''(1)=-4$, so the second derivative is
positive just to the left of $1$. For $S_8$, it is already $3/4$ at $1$.
This proves failure of the log-concavity mechanism beyond $m=6$; it makes
no assertion about whether the center-plus-polygon configuration is
optimal for eight points.

\subsection{The exact unrestricted seven-point maximum}
\label{sec:seven-global-exact}

The six-point rational-disk method extends to the next point count.
The external coverage input is the published physical classification
of Faug\`ere--Svartz~\cite[Sections 5.3--5.7 and Section 7,
Figure 4]{FS}. It is important to distinguish this input from an
independent recomputation of the entire complexified decomposition.
Figure 4 displays an $e_2=1$ component of degree $526$ and labels
its degree-$356$ branch physically empty. The archived author
file~\cite{FSdataSeven} omits that branch and supplies seven
triangular charts with parameter degrees
\[
                         1,1,2,6,18,82,64.
\]
The degree-$18$ branch is also labeled physically empty in the
paper; it is nevertheless included in our independent force check.
We use the following precise external statement: every physical
collision-free stationary seven-point shape occurs in one of the
supplied charts, up to complex scaling and permutation. Our
verifier does not reconstruct the missing degree-$356$ branch or
independently prove its physical emptiness. The author numerical
solution catalogue is not used as a substitute for this coverage
statement.

\begin{theorem}\label{thm:seven-global}
For $\sum_{i=1}^7|z_i|^2=7$,
\[
 \Delta_7(z)\le \frac{7^{21}}{6^{15}}
              =\frac{558545864083284007}{470184984576}.
\]
Equality holds precisely at the origin together with a regular
hexagon of radius $\sqrt{7/6}$, up to rotation and permutation.
The proof invokes the published physical coverage just specified.
\end{theorem}
\begin{proof}
A maximizer is centered and collision-free. In a supplied coefficient
chart its root polynomial has the convention
\[
 P(t)=t^7+e_2t^5-e_3t^4+e_4t^3-e_5t^2+e_6t-e_7.
\]
If $w_i$ are its roots and $Q=\sum|w_i|^2$, physical stationarity
requires, at every root,
\begin{equation}\label{eq:seven-force-residual}
                  QP''(w_i)-42\overline w_iP'(w_i)=0.
\end{equation}
Indeed the chart force constant is positive and its product with
$Q$ is $N_7=21$. The original normalized objective is
\begin{equation}\label{eq:seven-chart-value}
             (7/Q)^{21}\prod_{i<j}|w_i-w_j|^2.
\end{equation}

The exact standard-library checker \texttt{certify\_fs7\_disks.py}
covers every root of each supplied parameter polynomial with
disjoint rational disks and then every root of its seven-point
polynomial. All inequalities are strict integer or rational
comparisons, including when Python assertions are disabled.
There are $174$ parameter disks and $1218$ point-root disks.
The results are:
\begin{center}
\begin{tabular}{rrrrl}
\toprule
Chart&Parameters&Force excluded&Not excluded&Certified objective\\
\midrule
1&1&0&1&$823543=7^7$\\
2&1&0&1&$7^{21}/6^{15}$, exact\\
3&2&1&1&$<40253$\\
4&6&4&2&$<159666$\\
5&18&18&0&---\\
6&82&78&4&$<268962$\\
7&64&54&10&$<822629$\\
\bottomrule
\end{tabular}
\end{center}
Thus $155$ candidates are excluded by
\eqref{eq:seven-force-residual}, $18$ have strict objective bounds
below the target, and one is evaluated exactly. No sufficiency
claim is needed for the candidates bounded strictly below the target.

Chart two is $P=t(t^6+1)$: it has the origin and a rotated unit
hexagon. Its moment is six and its discriminant product is $6^6$,
so \eqref{eq:seven-chart-value} gives $7^{21}/6^{15}$.
The root polynomial or a geometric sum verifies its physical force
equations directly. Every other covered physical shape has a
strictly smaller value. External coverage now proves the upper
bound and the stated equality classification.
\end{proof}

For reproducibility the archived data have SHA-256
\begin{center}\footnotesize\ttfamily
d8cd0bc601071a7b4c0f7e9124c613eca9e3902cc58b29d8871686d541e73530.
\end{center}
The certificate contains $300$-digit rational parameter centers,
radius $10^{-250}$, and $45$-digit rational point centers, radius
$10^{-10}$. Their generation uses numerical arithmetic only to
propose the centers. Verification reparses the exact author data,
checks its hash, and uses the integer-arithmetic parameter evaluator
to avoid expensive repeated rational gcd computations.
For a parameter polynomial $f$, a center $b$ and radius $r$, it
checks
\[
 |f(b)|_{\rm up}+\frac{r^2}{2}
            \sum_{k\ge2}|a_k|k(k-1)B^{k-2}
       <|f'(b)|_{\rm low}r,
 \qquad B\ge |b|+r.
\]
The parameter-root count follows from Rouch\'e's theorem and the
degree of $f$. Derivative majorants propagate each parameter disk
through the rational polynomial coordinate formulae. The
seven-point polynomial checks and the force and objective estimates
are the same as the six-point rational-disk argument, with the
constants $42$, $7$ and $21$ in place of $30$, $6$ and $15$.

\begin{corollary}\label{cor:seven-seed-all-n}
The seven-point result strengthens the propagated bound to
\[
 M_n<\left(\frac n{n-1}\,6^{2/7}\right)^{n(n-1)/2}
                    \qquad(n>7).
\]
\end{corollary}
\begin{proof}
The normalized value in Theorem~\ref{thm:subset-monotonic} is
$d_7=(6/7)M_7^{1/21}=6^{2/7}$. Strict deletion monotonicity gives
$d_n<d_7$ for $n>7$.
\end{proof}
The agreement of $M_6$ and $M_7$ with the central-polygon values
does not justify extrapolation to all larger $n$.
The certified ten-point improvement and the fixed-ring obstruction
remain explicit limits of that template.

\subsection{Conjugate pairing from one extra trace row}
The following formulation is a supplement to the existing exact
trace test. It replaces the composition congruence
$\overline P(R)\equiv0$ by Hermitian symmetry of one bordered
trace matrix. It does not provide a formula for its largest determinant.

\begin{theorem}[Bordered trace test]\label{thm:bordered-trace}
Let $P(t)$ be monic of degree $n$, let $\lambda>0$ be real,
and let $R(t)$ have degree less than $n$. In
$\mathcal A=\C[t]/(P)$ define
\[
 K_{jk}=\operatorname{Tr}_{\mathcal A}(t^jR^k),
 \qquad 0\le j,k\le n,
 \qquad H=(K_{jk})_{0\le j,k<n}.
\]
Then the following conditions are equivalent to $P$ having distinct
roots $w_i$ satisfying
\[
 R(w_i)=\overline w_i,
 \qquad \sum_{\ell\ne i}\frac1{w_i-w_\ell}
   =\lambda\overline w_i:
\]
\[
 P''-2\lambda RP'\equiv0\pmod P,
 \qquad K=K^*,\qquad H\succ0.
\]
When they hold, $K\succeq0$ has rank $n$,
$\sum_iw_i=0$, $\sum_i|w_i|^2=N/\lambda$, and
\[
 \det H=|\operatorname{disc}P|,\qquad
 \Delta_n\!\left(\sqrt{\frac{n\lambda}{N}}\,w\right)
 =\left(\frac{n\lambda}{N}\right)^N\det H.
\]
All matrix entries and remainder equations are polynomial in
the coefficients of $P,R$ and their conjugates; no discriminant
denominator is needed when $R$ is retained as an auxiliary polynomial.
\end{theorem}
\begin{proof}
For distinct roots with the stated conjugate pairing,
$K=\widetilde V\widetilde V^*$, where
$\widetilde V_{ji}=w_i^j$ for $0\le j\le n$. Its leading block
is positive definite, and the force congruence follows by evaluation.

Conversely, for distinct values $\xi_i$ occurring as roots of $P$
with multiplicities $m_i$, quotient-algebra traces are
\[
 K_{jk}=\sum_i m_i\xi_i^jR(\xi_i)^k.
\]
Thus $\operatorname{rank}H$ is at most the number of distinct roots.
Positive definiteness forces $n$ distinct roots $w_i$. The factorization
$H=V W^{\mathsf T}$, with $V_{ji}=w_i^j$ and
$W_{ki}=R(w_i)^k$, then also forces the $R(w_i)$ to be distinct.

Write $P(t)=t^n+\sum_{a=0}^{n-1}c_at^a$. The last row of $K$
satisfies $K_{n,k}=-\sum_a c_aK_{a,k}$. Hermitian symmetry yields
\[
 K_{j,n}=-\sum_a\overline c_aK_{j,a},\qquad0\le j<n.
\]
Equivalently,
\[
 \sum_i w_i^j\overline P(R(w_i))=0,\qquad0\le j<n.
\]
Invertibility of $V$ proves $\overline P(R(w_i))=0$ for every $i$.
Since the force values are distinct, they are a permutation of the
conjugate roots. Consequently $H=V\Pi V^*$ for a permutation matrix
$\Pi$. Hermitian positive definiteness passes to $\Pi$. A Hermitian
permutation is an involution, and each nontrivial transposition has
eigenvalue $-1$. Therefore $\Pi=I$, proving the actual conjugate pairing.
The force congruence gives the claimed force equations. Summing them
gives centering; multiplying by $w_i$ and summing gives
$\lambda\sum_i|w_i|^2=N$. The Gram determinant and scaling identities
follow immediately.
\end{proof}

The extra row is essential. For $n=2$, $\lambda=1$,
$P=t^2-4$, $R=t/8$, the force congruence holds and
$H=\operatorname{diag}(2,1)\succ0$, but the roots are not
physical stationary roots at that scale. The border detects this:
$K_{20}=8$ whereas $K_{02}=1/8$.

\subsection{An explicit exact-output specification for every point count}
The existing critical-value elimination and moment selection already
give termination without assuming that stationary components are
zero-dimensional. To make the output unambiguous, after saturation
and elimination take the squarefree primitive integer generator
$p_n(T)$, factor it over $\Q$, and isolate all its real roots.
Successive rational moment enclosures identify exactly one such root.
Output the primitive irreducible factor $m_n\in\mathbb Z[T]$ containing
that root, with positive leading coefficient, together with a rational
open interval containing exactly one real root of $m_n$.
This pair is an exact representation of $M_n$, even if $M_n$ is not
expressible by radicals. Algebraic sampling of the corresponding
real maximizing fiber can also output attaining algebraic coordinates.

For the executed or independently proved examples above, the
minimal polynomials are $T-4$, $T-27$, $T-256$, and $T-3125$.
The previously certified six-point value has primitive minimal
polynomial $9765625T-470184984576$. This specification adds a
reproducible output format to the existing algorithm; it is not a
new general finite product or a computed elimination for every $n$.

\subsection{A stronger separation bound from simultaneous pair replacement}
\label{sec:pair-separation}

The compact disk bound can be complemented by a strictly positive
separation constant independent of the point count at force-one scale.
This estimate uses global optimality and is not an assertion about all
stationary solutions.

\begin{lemma}\label{lem:weighted-zero-separation}
Let $p$ be a nonzero polynomial and $\alpha>0$. Suppose $a$ maximizes
$|p(z)|^2e^{-\alpha|z|^2}$ over $\C$. For every zero $a+\delta$ of $p$,
\[
                         |\delta|^2>\frac{2}{e\alpha}.
\]
\end{lemma}
\begin{proof}
The maximum is positive, so $p(a)\ne0$. The entire function
\[
 g(u)=\frac{p(a+u)}{p(a)}e^{-\alpha\overline a u}
\]
satisfies $g(0)=1$, $g'(0)=0$ and
$|g(u)|\le e^{\alpha|u|^2/2}$. The derivative vanishes by the
stationary equation $p'(a)/p(a)=\alpha\overline a$.
Choose $R^2=2/\alpha$ and $b=e^{-1}$.
On $|u|\le R$, $f(u)=bg(u)$ has modulus at most one, $f(0)=b$,
and $f'(0)=0$. The disk map
\[
 h(u)=\frac{f(u)-b}{1-bf(u)}
\]
has a double zero at zero. Applying Schwarz's lemma twice gives
$|h(u)|\le |u/R|^2$.
If $|\delta|\ge R$, the claimed strict inequality is immediate.
Otherwise $f(\delta)=0$ implies
$b=|h(\delta)|\le |\delta/R|^2$.
Equality would force $h(u)=\xi(u/R)^2$ for $|\xi|=1$.
Then $f=(b+h)/(1+bh)$ on the disk, and hence by analytic continuation
where defined. The latter rational function has nonremovable poles
at the two zeros of $1+b\xi(u/R)^2$, since $0<b<1$.
This contradicts the fact that $f$ is entire. Thus the inequality
is strict.
\end{proof}

\begin{theorem}\label{thm:pair-separation}
Every global maximizer with quadratic moment $n$ satisfies
\begin{equation}\label{eq:pair-separation-new}
             |z_i-z_j|^2>\frac{8}{e(n-1)}\quad(i\ne j).
\end{equation}
For $n\ge4$ this improves the separation bound
$2n/(n-1)^2$ in Lemma~\ref{lem:global-separation}.
\end{theorem}
\begin{proof}
As in Theorem~\ref{thm:global-fixed-disk}, the configuration globally
maximizes $\Delta_n\exp(-\lambda\sum|z_i|^2)$ on the full space,
where $\lambda=(n-1)/2$.
Fix a pair, its midpoint $m$, and all other points. Write the pair
as $m+w,m-w$, and put
\[
 p(w)=w\prod_{k\ne i,j}\bigl(w^2-(m-z_k)^2\bigr).
\]
Up to a positive factor independent of $w$, the weighted objective
is $|p(w)|^2e^{-(n-1)|w|^2}$.
The original $w_*=(z_i-z_j)/2$ is a global maximizer of this
one-variable problem. Since $p(0)=0$, the preceding lemma gives
$|w_*|^2>2/(e(n-1))$. Multiplication by four proves
\eqref{eq:pair-separation-new}.
Finally $e<3$ and $4(n-1)/n\ge3$ for $n\ge4$, so the new
bound is strictly stronger there.
\end{proof}

\subsection{Rational improvements to the compact physical trace model}
Put $w_i=\sqrt{(n-1)/2}\,z_i$ and
\[
\begin{gathered}
 a_n=\lceil\sqrt{n-1}\rceil,\qquad
 \eta_n=\max\left\{\frac n{n-1},\frac{16}{11}\right\},\\
 \widehat\mu_n=\frac{\eta_n^{n-1}}{n(1+a_n)^{2n-2}},
 \qquad
 \widehat B_n=\frac{n^2a_n^2(1+a_n)^{2n-2}}{\eta_n^{n-1}}.
\end{gathered}
\]
The elementary rational estimate
\[
 e<\sum_{j=0}^5\frac1{j!}+\frac{1/6!}{1-1/7}<\frac{11}{4}
\]
shows $4/e>16/11$. Combining the new pair bound with
Lemma~\ref{lem:global-separation} gives
$|w_i-w_j|^2\ge\eta_n$ at every global maximum.
Thus the Lagrange polynomials have coefficient $\ell^1$ norm at most
$(1+a_n)^{n-1}\eta_n^{-(n-1)/2}$.
The same inverse-Vandermonde argument as before gives the stronger
rational inequalities
\begin{equation}\label{eq:new-pair-compact-bounds}
 H\succeq\widehat\mu_n I,\qquad |r_\ell|^2\le\widehat B_n,
 \qquad |c_k|^2\le\binom nk^2(n-1)^k.
\end{equation}
Replacing the corresponding inequalities in the compact trace model
preserves its exact maximum: every global maximizer remains feasible,
and every feasible conjugate-paired stationary configuration is still
a configuration in the original problem.
For $n=6$ the numbers simplify to
\[
                  \widehat\mu_6=\frac1{966306},\qquad
                  \widehat B_6=52180524.
\]
The previous disk-only margin was $1/6291456$ and its squared
auxiliary bound was $18432^2=339738624$.

A further $n$-by-$n$ polynomial localizing inequality can be added
without enumerating pairs. In the quotient algebra put
\[
 D(t)=P'(t)\,\overline{P'}(R(t)),\qquad
 K_{jk}=\operatorname{Tr}\bigl(t^jR(t)^k
                  [D(t)-\eta_n^{n-1}]\bigr),\quad 0\le j,k<n.
\]
Here the bar conjugates coefficients. On physical roots,
\[
 K=V\operatorname{diag}\left(|P'(w_i)|^2-\eta_n^{n-1}\right)V^*.
\]
Since $|P'(w_i)|^2=\prod_{j\ne i}|w_i-w_j|^2$, one has
$K\succeq0$ for every global maximizer. Adding this closed rational
polynomial constraint again preserves the maximum of the compact
model. It controls products of separations, rather than imposing
each individual separation equivalently.
The standard-library checker \texttt{verify\_examples\_separation.py}
checks the rational constants and matrix-factor identities independently.

\subsection{Exact finite certificates after critical-value augmentation}
\label{sec:finite-value-sos}

The critical-value method can be made into a finite exact
sum-of-squares certificate, even when the stationary configurations
themselves are not isolated. The additional equation is essential:
this statement does not assert finite convergence of the unaugmented
moment hierarchy.

\begin{theorem}\label{thm:finite-value-sos}
Let $K=\{x:e_a(x)=0,\ g_b(x)\ge0\}$ be nonempty and compact,
with an Archimedean quadratic module modulo the equality ideal.
Let $F$ be a polynomial taking only finitely many values on $K$.
Suppose a square-free polynomial
\[
                  p(t)=\prod_{j=1}^s(t-\lambda_j),
                 \qquad \lambda_j\in\R\text{ distinct},
\]
satisfies $p(F)=0$ on $K$. It may have real roots that are not
attained on $K$. Put $M=\max_K F$. After adding the redundant
equation $p(F)=0$, there is a finite polynomial identity
\begin{equation}\label{eq:finite-value-exact-sos}
 M-F=\sigma_0+\sum_b\sigma_b g_b+\sum_a\tau_a e_a
                              +\tau_p p(F),
\end{equation}
where all $\sigma_b$ are sums of squares.
If all input coefficients and $M,\lambda_j$ are real algebraic,
the certificate can be chosen with real algebraic coefficients.
\end{theorem}
\begin{proof}
Let $E_j(t)$ be the Lagrange polynomial with
$E_j(\lambda_k)=\delta_{jk}$. Modulo $p$,
\[
 E_j^2=E_j,\qquad (t-\lambda_j)E_j=0,\qquad
              M-t=\sum_j(M-\lambda_j)E_j^2.
\]
When $\lambda_j\le M$, the corresponding term is a square with
a nonnegative constant multiplier. When $\lambda_j>M$, the set
$K\cap\{F=\lambda_j\}$ is empty. Its quadratic module remains
Archimedean. The empty-set form of Putinar's theorem gives
\[
 -1=q_j+\sum_a u_{ja}e_a+v_j(F-\lambda_j),
\]
where $q_j$ belongs to the quadratic module generated by the $g_b$.
Multiplying by $E_j(F)^2$ preserves membership in that module, and
the last term is a multiple of $p(F)$ because
$(t-\lambda_j)E_j(t)^2\equiv0\pmod p$.
Consequently $-E_j(F)^2$ has such a certificate modulo the augmented
ideal. Multiplication by $\lambda_j-M>0$ handles each formerly
negative term. Adding the finitely many identities proves
\eqref{eq:finite-value-exact-sos}.

For the coefficient assertion, take the finite degrees and numbers
of squares in one real certificate. The coefficient identity is
then a finite polynomial system over the real algebraic numbers,
with the square factors as variables. The real algebraic numbers
form a real closed field; transfer for real closed fields supplies
a solution in that field~\cite{Basu}. This proves existence of algebraic
coefficients, without claiming that a numerical SDP output is
automatically an exact certificate.
\end{proof}

For the compact physical trace model, $F$ is the normalized
determinant and every feasible configuration is stationary.
Section~\ref{sec:exact-algorithm} establishes finiteness of its
critical values without assuming zero-dimensional shape charts.
Isolate the real roots of the rational critical-value polynomial
and form their square-free product over the real algebraic numbers.
The theorem therefore supplies a finite exact certificate for
each $n$ after critical-value augmentation. In principle one can
increase the certificate degree and use exact real algebraic
feasibility until it succeeds. No uniform small degree is obtained,
and construction of the critical polynomial may still be costly.
The result is a uniform terminating certification mechanism, not
a short expression for the numerical value of every $M_n$.

\subsection{A fully explicit three-point demonstration}
For $n=3$, every collision-free centered stationary configuration
has value $27$ or $27/2$. Here is a short derivation. Write
$P(t)=t^3+at+b$. The force constant is one, so at each root
\[
 3z_i=\overline z_i(3z_i^2+a),\qquad
             3(1-|z_i|^2)z_i^2=a|z_i|^2.
\]
If $a=0$, none of the roots is zero, and all have modulus one;
centering forces an equilateral triangle. If $a\ne0$, any nonzero
root satisfies $z_i^2/a\in\R$. After a rotation the roots lie on
the two coordinate axes. A centered triple cannot occupy both
axes with nonzero points, so it is collinear. For real roots the
stationary polynomial identity is
\[
                 P''-2tP'+6P=0,
\]
which gives $P=t^3-3t/2$ and value $27/2$.
The equilateral value is $27$. Thus on the stationary set,
\[
                p_3(F)=(F-27)(2F-27)=0.
\]
Set $E=2(27-F)/27$. Direct rational polynomial expansion gives
\begin{equation}\label{eq:three-point-value-sos}
 27-F=(3E)^2+\left(\frac32 E\right)^2
          +\left(\frac32 E\right)^2-\frac1{27}p_3(F).
\end{equation}
After imposing the valid critical-value equation this is an exact
upper-bound certificate, attained only at the equilateral value.
It exhibits the univariate certificate construction explicitly;
the elementary direct triangle proof remains shorter if one needs
only $M_3$. The verifier \texttt{verify\_examples\_value\_sos.py}
uses rational arithmetic for the identity.

\subsection{An explicit six-square certificate for the four-point value}
Section~\ref{sec:four-critical-values} supplies the complete physical
value set $256,4096/27,1024/27$. Put $p_4=p_4^{\rm phys}$ and
let $E_\triangle,E_{\rm line}$ be the Lagrange selectors for the
last two values:
\[
 E_\triangle(T)=-\frac{27(T-256)(27T-1024)}{2816\cdot3072},
 \qquad
 E_{\rm line}(T)=\frac{27(T-256)(27T-4096)}{5888\cdot3072}.
\]
The positive gaps have rational square decompositions
\[
 \frac{2816}{27}=\left(\frac{80}{9}\right)^2
                      +2\left(\frac{32}{9}\right)^2,
 \qquad
 \frac{5888}{27}=\left(\frac{128}{9}\right)^2
                   +\left(\frac{32}{9}\right)^2
                   +\left(\frac{16}{9}\right)^2.
\]
Define the sum of six squares
\[
\begin{split}
 S(T)={}&\left(\frac{80}{9}E_\triangle\right)^2
       +2\left(\frac{32}{9}E_\triangle\right)^2
       +\left(\frac{128}{9}E_{\rm line}\right)^2\\
       &+\left(\frac{32}{9}E_{\rm line}\right)^2
       +\left(\frac{16}{9}E_{\rm line}\right)^2.
\end{split}
\]
An exact polynomial identity, verified by rational expansion, is
\begin{equation}\label{eq:four-value-six-squares}
 256-T=S(T)+
       \frac{11008-51T}{33957085184}\,p_4(T).
\end{equation}
On every physical stationary configuration the last term vanishes.
Thus this finite certificate proves the global upper bound once
stationary-value coverage is established, and the unit square
provides the attaining lower bound. It is a complete executed
example of Theorem~\ref{thm:finite-value-sos}, with no numerical SDP
or assumption of isolated complexified stationary shapes.

\section*{Acknowledgement of assistance}
AI assistance through OpenAI's Codex was used for mathematical
exploration, verification support and manuscript preparation.
The author is responsible for checking the final claims and sources.

\begin{thebibliography}{99}
\small
\setlength{\itemsep}{2pt}
\bibitem{Chalk}
J.~H.~H. Chalk, \emph{A note on a discriminant inequality},
Canadian Mathematical Bulletin \textbf{3} (1960), no.~1, 7--9.
\href{https://doi.org/10.4153/CMB-1960-002-x}{doi:10.4153/CMB-1960-002-x}.
\bibitem{Wang}
J.~Wang, \emph{The Five-Point Case of Mordell's Discriminant Inequality},
arXiv:2610.00358v1 (2026). \url{https://arxiv.org/abs/2610.00358}.
\bibitem{AOC}
Y.~Ameur and J.~Ortega-Cerd\`a,
\emph{Beurling--Landau densities of weighted Fekete sets and correlation
kernel estimates}, Journal of Functional Analysis \textbf{263}
(2012), no.~7, 1825--1861. \url{https://arxiv.org/abs/1110.0284}.
\bibitem{SS}
E.~Sandier and S.~Serfaty, \emph{2D Coulomb Gases and the Renormalized
Energy}, arXiv:1201.3503v2 (2013).
\url{https://arxiv.org/abs/1201.3503}.
\bibitem{Basu}
S.~Basu, \emph{Algorithms in Real Algebraic Geometry: A Survey},
arXiv:1409.1534 (2014). \url{https://arxiv.org/abs/1409.1534}.
\bibitem{CS}
H.~E. Cabral and D.~S. Schmidt,
\emph{Stability of relative equilibria in the problem of $N+1$ vortices},
SIAM Journal on Mathematical Analysis \textbf{31} (1999), no.~2, 231--250.
\href{https://doi.org/10.1137/S0036141098302124}{doi:10.1137/S0036141098302124}.
\bibitem{FS}
J.-C. Faug\`ere and J.~Svartz,
\emph{Solving polynomial systems globally invariant under an action of
the symmetric group and application to the equilibria of $N$ vortices
in the plane}, Proceedings of ISSAC 2012, 170--178.
\href{https://doi.org/10.1145/2442829.2442856}{doi:10.1145/2442829.2442856}.
\bibitem{Beauzamy}
B.~Beauzamy, \emph{Extremal products in Bombieri's norm},
January 1995.
\url{https://www.scmsa.eu/BB_Publications/BB_Extremal_products_Bombieri_norm.pdf}.
\bibitem{TF}
G.~T\'ellez and P.~J. Forrester,
\emph{Expanded Vandermonde powers and sum rules for the
two-dimensional one-component plasma}, Journal of Statistical Physics
\textbf{148} (2012), 824--855.
\url{https://arxiv.org/abs/1204.6003}.
\bibitem{Lasserre}
J.~B. Lasserre, \emph{Global optimization with polynomials and the
problem of moments}, SIAM Journal on Optimization \textbf{11}
(2001), no.~3, 796--817.
\href{https://doi.org/10.1137/S1052623400366802}{doi:10.1137/S1052623400366802}.
\bibitem{FSdata}
J.-C. Faug\`ere and J.~Svartz, \emph{Six-vortex lexicographic bases
and solution appendix}, supplementary author data for \cite{FS}.
\href{https://web.archive.org/web/20240423101624id_/https://www-polsys.lip6.fr/~jcf/vortices/vortices_N6.txt}{Archived six-point bases};
\href{https://web.archive.org/web/20220305130108id_/https://www-polsys.lip6.fr/~jcf/vortices/SOLUTIONS\%20N7.pdf}{archived author appendix}, Section~4.
The unchanged bases and rational certificates are included in the source bundle.
\bibitem{FSdataSeven}
J.-C. Faug\`ere and J.~Svartz, \emph{Seven-vortex lexicographic bases},
supplementary author data for \cite{FS}.
\href{https://web.archive.org/web/20240423101625id_/https://www-polsys.lip6.fr/~jcf/vortices/vortices_N7.txt}{Archived seven-point bases}.
The file omits the degree-$356$ branch declared physically empty in
\cite[Section 7, Figure 4]{FS}; this external coverage input is retained.
\end{thebibliography}
\end{document}
